% =====================================================================
%  PLANTILLA DE PLAN DE PRUEBAS DE SOFTWARE
%  Basada en ISO/IEC/IEEE 29119 (partes 2, 3 y 4)
%
%  Universidad Veracruzana · Facultad de Estadística e Informática
%  Licenciatura en Ingeniería de Software · EE. Pruebas de Software
%
%  Compilación:   latexmk -pdf main.tex      (o simplemente: make)
%  Limpieza:      make clean
%
%  Este archivo NO contiene contenido del plan: sólo configuración,
%  metadatos y el orden de inclusión de los módulos.
%  El contenido está en sections/ , appendices/ y diagrams/.
% =====================================================================
\documentclass[11pt,a4paper,oneside]{article}

% =====================================================================
%  config/preamble.tex
%  ---------------------------------------------------------------
%  Paquetes, geometría, tipografía, encabezados y estilo de tablas.
%  Todo lo aquí usado forma parte de una instalación estándar de
%  TeX Live / MacTeX. No se requieren paquetes exóticos ni servicios web.
% =====================================================================

% ------------------------- Codificación e idioma ---------------------
\usepackage[T1]{fontenc}
\usepackage[utf8]{inputenc}
\usepackage{lmodern}
\usepackage[spanish,mexico,es-nolists,es-noindentfirst]{babel}
% Nota: la opción «mexico» usa punto decimal.
%
% Rótulos: la plantilla Word original habla de «cuadros y figuras», y así
% se citan en el texto. Las versiones recientes de babel-spanish rotulan
% las tablas como «Tabla», por lo que se fija explícitamente «Cuadro»
% para que el rótulo automático y las menciones del texto coincidan.
% Si prefiere «Tabla», cambie las dos líneas siguientes y, con ellas,
% las menciones «Cuadro~\ref{...}» del texto.
\addto\captionsspanish{%
  \renewcommand{\tablename}{Cuadro}%
  \renewcommand{\listtablename}{Índice de cuadros}%
  \renewcommand{\listfigurename}{Índice de figuras}%
  \renewcommand{\contentsname}{Tabla de contenido}%
}

% ---- Símbolos Unicode de uso frecuente ------------------------------
% Permiten que el equipo escriba directamente ×, ÷, ≥, ≤, ≠, → … en el
% texto y en las tablas sin recurrir a modo matemático.
\usepackage{amsmath}
\DeclareUnicodeCharacter{00D7}{\ensuremath{\times}}
\DeclareUnicodeCharacter{00F7}{\ensuremath{\div}}
\DeclareUnicodeCharacter{2264}{\ensuremath{\leq}}
\DeclareUnicodeCharacter{2265}{\ensuremath{\geq}}
\DeclareUnicodeCharacter{2260}{\ensuremath{\neq}}
\DeclareUnicodeCharacter{2192}{\ensuremath{\rightarrow}}
\DeclareUnicodeCharacter{2190}{\ensuremath{\leftarrow}}
\DeclareUnicodeCharacter{2248}{\ensuremath{\approx}}
\DeclareUnicodeCharacter{2212}{\ensuremath{-}}
\DeclareUnicodeCharacter{2022}{\textbullet}

% --------------------------- Tipografía ------------------------------
\usepackage{helvet}
\renewcommand{\familydefault}{\sfdefault}   % aproxima la Arial del Word
\usepackage{setspace}
\onehalfspacing
\usepackage{microtype}

% ---------------------------- Geometría ------------------------------
\usepackage{geometry}
\geometry{a4paper,
  left=2.5cm, right=2.5cm, top=2.4cm, bottom=2.2cm,
  headheight=46pt, headsep=12pt, footskip=34pt}

% ------------------------------ Color --------------------------------
\usepackage{xcolor}
\definecolor{uvazul}{HTML}{1F3864}   % azul institucional para títulos
\definecolor{uvgris}{HTML}{D9D9D9}   % encabezados de tabla
\definecolor{uvgrisc}{HTML}{F2F2F2}  % relleno suave
\definecolor{instrverde}{HTML}{2E6B34}
\definecolor{instrfondo}{HTML}{EDF6EE}
\definecolor{polinaranja}{HTML}{9C5A00}
\definecolor{polfondo}{HTML}{FDF3E3}
\definecolor{ejemazul}{HTML}{1F4E79}
\definecolor{ejemfondo}{HTML}{EAF1F8}
\definecolor{normgris}{HTML}{4A4A4A}
\definecolor{normfondo}{HTML}{F4F4F4}
\definecolor{obsmorado}{HTML}{6B3FA0}  % observación de la revisión (variante anotada)
\definecolor{obsfondo}{HTML}{F3EEFA}
\definecolor{campoamar}{HTML}{FFF2CC}
\definecolor{campotexto}{HTML}{9A3412}  % marcador de campo por rellenar

% ------------------------- Tablas y figuras --------------------------
\usepackage{array}
\usepackage{tabularx}
\usepackage{longtable}
\usepackage{xltabular}   % longtable + columnas X: fichas que pueden partirse entre páginas
\usepackage{multirow}
\usepackage{colortbl}
\usepackage{booktabs}
\usepackage{caption}
\usepackage{graphicx}
\usepackage{float}

\captionsetup{font=small,labelfont=bf,justification=centering,skip=6pt}
\captionsetup[table]{position=above}

% Tipos de columna reutilizables ---------------------------------------
% El \hspace{0pt} inicial permite que TeX divida con guion la PRIMERA
% palabra de la celda (sin él, TeX nunca parte la primera palabra de un
% párrafo, y las columnas estrechas se desbordan).
\newcolumntype{Y}{>{\RaggedRight\arraybackslash\hspace{0pt}}X}                       % flexible
\newcolumntype{C}[1]{>{\Centering\arraybackslash\hspace{0pt}}p{#1}}                  % centrada fija
\newcolumntype{L}[1]{>{\RaggedRight\arraybackslash\hspace{0pt}}p{#1}}                % izquierda fija
\newcolumntype{E}[1]{>{\columncolor{uvgrisc}\bfseries\RaggedRight\arraybackslash\hspace{0pt}}p{#1}} % etiqueta

\renewcommand{\arraystretch}{1.25}
\setlength{\tabcolsep}{5pt}

% ------------------------- Índice (TOC) ------------------------------
\usepackage{tocloft}
\setlength{\cftbeforesecskip}{2pt}
% La numeración de los anexos («ANEXO A») es mucho más ancha que la de
% una sección numérica, por lo que \iniciaranexos ensancha la caja del
% número a partir de ese punto del índice (véase config/commands.tex).

% --------------------------- Listas ----------------------------------
\usepackage{ragged2e}   % \RaggedRight conserva la partición de palabras
\usepackage{enumitem}
\setlist{nosep,topsep=3pt,itemsep=2pt,parsep=0pt,leftmargin=1.4em}

% ------------------------ Encabezado y pie ---------------------------
\usepackage{fancyhdr}
\usepackage{lastpage}

% Los preliminares se numeran en romanos y el cuerpo reinicia en arábigos,
% de modo que «de N» sólo tiene sentido dentro del cuerpo: \pageref{LastPage}
% devuelve el último número arábigo, no el total físico de páginas.
% main.tex activa este interruptor al llamar a \pagenumbering{arabic}.
\newif\ifcontarpaginas

\newcommand{\bandaEncabezado}{%
  \setlength{\tabcolsep}{4pt}%
  \renewcommand{\arraystretch}{1.05}%
  \footnotesize
  \begin{tabularx}{\textwidth}{|>{\raggedright\arraybackslash}X|C{3.2cm}|>{\raggedleft\arraybackslash}X|}
    \hline
    \uvPrograma & \textbf{\uvActividad} & \uvEE \\ \hline
    \multicolumn{3}{|c|}{\uvBandaHeader} \\ \hline
  \end{tabularx}%
}

\newcommand{\bandaPie}{%
  \setlength{\tabcolsep}{4pt}%
  \renewcommand{\arraystretch}{1.05}%
  \footnotesize
  \begin{tabularx}{\textwidth}{|>{\raggedright\arraybackslash}X|>{\raggedleft\arraybackslash}p{5.4cm}|}
    \hline
    \textbf{\uvUniversidad}\ \textperiodcentered\ \uvFacultad
      & Página \thepage\ifcontarpaginas\ de \pageref{LastPage}\fi\quad \docVersionPlantilla \\ \hline
  \end{tabularx}%
}

\fancypagestyle{uvestilo}{%
  \fancyhf{}%
  \renewcommand{\headrulewidth}{0pt}%
  \renewcommand{\footrulewidth}{0pt}%
  \fancyhead[C]{\bandaEncabezado}%
  \fancyfoot[C]{\bandaPie}%
}
\pagestyle{uvestilo}
\fancypagestyle{plain}{\pagestyle{uvestilo}}

% ------------------------ Títulos de sección -------------------------
\usepackage{titlesec}
\titleformat{\section}[block]
  {\normalfont\large\bfseries\color{uvazul}}{\thesection.}{0.6em}{}
\titleformat{\subsection}[block]
  {\normalfont\normalsize\bfseries\color{uvazul}}{\thesubsection}{0.6em}{}
\titleformat{\subsubsection}[block]
  {\normalfont\normalsize\bfseries}{\thesubsubsection}{0.6em}{}
\titlespacing*{\section}{0pt}{14pt}{6pt}
\titlespacing*{\subsection}{0pt}{12pt}{4pt}
\titlespacing*{\subsubsection}{0pt}{10pt}{3pt}

% ---------------------- Cajas de la plantilla ------------------------
\usepackage{etoolbox}
\usepackage[most]{tcolorbox}
\usepackage{comment}

% ----------------------------- Enlaces -------------------------------
\usepackage[hidelinks,bookmarksnumbered,bookmarksopen]{hyperref}
\usepackage{bookmark}

% ------------------------------ TikZ ---------------------------------
\usepackage{tikz}
\usetikzlibrary{positioning,arrows.meta,shapes.geometric,fit,backgrounds,calc,decorations.pathreplacing}

% ------------------- Ajustes finales de composición ------------------
\setlength{\parskip}{6pt}
\setlength{\parindent}{0pt}
\setlength{\emergencystretch}{3em}   % reduce overfull hboxes en texto justificado
\widowpenalty=10000
\clubpenalty=10000
    % paquetes, geometría, encabezados, estilos
% =====================================================================
%  config/metadata.tex
%  ---------------------------------------------------------------
%  ÚNICO archivo que el equipo necesita editar para personalizar la
%  identificación del documento. No contiene contenido del plan.
% =====================================================================

% ---------------------------------------------------------------------
% 1. Identificación institucional (encabezado y pie de página)
% ---------------------------------------------------------------------
\newcommand{\uvUniversidad}{UNIVERSIDAD VERACRUZANA}
\newcommand{\uvFacultad}{Facultad de Estadística e Informática}
\newcommand{\uvPrograma}{Licenciatura en Ingeniería de Software}
\newcommand{\uvEE}{EE. Pruebas de Software}
\newcommand{\uvActividad}{PROYECTO FINAL}
\newcommand{\uvBandaHeader}{Proyecto de Pruebas}

% ---------------------------------------------------------------------
% 2. Identificación del documento
% ---------------------------------------------------------------------
\newcommand{\docTitulo}{Proyecto de pruebas de software}
\newcommand{\docProyecto}{\campo{Nombre del proyecto}}
\newcommand{\docFecha}{\campo{dd/mm/aaaa}}
\newcommand{\docEquipo}{\campo{Nombre del equipo}}
\newcommand{\docIntegrantes}{\campo{Matrícula y nombre de cada integrante}}
\newcommand{\docVersionPlantilla}{Ver. 01/10/25}   % versión de la PLANTILLA
\newcommand{\docVersionDoc}{\campo{0.1}}           % versión del DOCUMENTO del equipo

% ---------------------------------------------------------------------
% 3. Autoridad de aprobación
%    Quién autoriza desviaciones respecto de la estrategia acordada.
%    En el curso es el profesor; en un proyecto profesional cámbielo.
% ---------------------------------------------------------------------
\newcommand{\autoridadAprobacion}{\campo{Profesor de la Experiencia Educativa}}
\newcommand{\autoridadAceptacion}{\campo{Rol con autoridad para aceptar el riesgo residual}}

% ---------------------------------------------------------------------
% 4. Interruptores de la plantilla
%    Cambie a \...false para producir la versión final de entrega.
% ---------------------------------------------------------------------
\newif\ifmostrarinstrucciones
\mostrarinstruccionestrue      % \mostrarinstruccionesfalse  -> oculta las guías
\newif\ifmostrarejemplos
\mostrarejemplostrue           % \mostrarejemplosfalse       -> oculta los ejemplos

% «make final» compila una vez más definiendo \VERSIONFINAL, lo que apaga
% ambos interruptores sin tocar este archivo y deja el PDF de entrega en
% plan-pruebas-final.pdf. Editar las dos líneas de arriba produce el mismo
% resultado de forma permanente.
\ifdefined\VERSIONFINAL
  \mostrarinstruccionesfalse
  \mostrarejemplosfalse
\fi

% Las cajas de "Política académica" NO se ocultan con los interruptores
% anteriores: son contenido vigente del plan mientras el curso las exija.
% Si reutiliza la plantilla fuera del curso, elimine o sustituya esas cajas.
    % <-- EDITE AQUÍ los datos del equipo/proyecto
% =====================================================================
%  config/commands.tex
%  ---------------------------------------------------------------
%  Sistema de marcado de la plantilla. Define CINCO tipos de bloque
%  visualmente distinguibles, de modo que quien lea el documento
%  identifique de inmediato la naturaleza de cada texto:
%
%    instruccion        -> qué debe escribir el equipo (guía de la plantilla)
%    ejemplo            -> ilustración; NO es contenido a conservar
%    politicaacademica  -> regla local del curso, NO requisito de la norma
%    notanormativa      -> referencia o aclaración conceptual con fuente
%    observacion        -> origen de una corrección respecto de la plantilla
%                          anterior; SÓLO aparece en la variante anotada
%
%  Requiere que config/metadata.tex se cargue ANTES que este archivo,
%  porque aquí se consultan los interruptores \ifmostrarinstrucciones
%  y \ifmostrarejemplos.
% =====================================================================

% ---------------------------------------------------------------------
% Estilo base común a todas las cajas
% ---------------------------------------------------------------------
\tcbset{
  cajaplantilla/.style={
    enhanced, breakable,
    boxrule=0.6pt, arc=2pt,
    left=6pt, right=6pt, top=5pt, bottom=5pt,
    fonttitle=\bfseries\footnotesize,
    fontupper=\small,
    attach boxed title to top left={xshift=6pt, yshift=-2pt},
    boxed title style={boxrule=0pt, arc=1pt, left=4pt, right=4pt, top=1pt, bottom=1pt},
    toptitle=2pt,
    before skip=8pt, after skip=8pt
  }
}

% ---------------------------------------------------------------------
% 1. INSTRUCCIÓN — guía de la plantilla (se oculta en la versión final)
% ---------------------------------------------------------------------
\ifmostrarinstrucciones
  \newtcolorbox{instruccion}[1][Qué debe documentarse en este apartado]{%
    cajaplantilla,
    colback=instrfondo, colframe=instrverde,
    coltitle=white, colbacktitle=instrverde,
    title={#1}}
\else
  \excludecomment{instruccion}
\fi

% ---------------------------------------------------------------------
% 2. EJEMPLO — ilustrativo, se elimina en la versión final
% ---------------------------------------------------------------------
\ifmostrarejemplos
  \newtcolorbox{ejemplo}[1][Ejemplo ilustrativo]{%
    cajaplantilla,
    colback=ejemfondo, colframe=ejemazul,
    coltitle=white, colbacktitle=ejemazul,
    title={#1}}
\else
  \excludecomment{ejemplo}
\fi

% ---------------------------------------------------------------------
% 3. POLÍTICA ACADÉMICA — regla local del curso. SIEMPRE visible.
%    Se conserva porque el curso la exige, pero queda marcada para que
%    nunca se confunda con un requisito de ISO/IEC/IEEE 29119.
% ---------------------------------------------------------------------
\newtcolorbox{politicaacademica}[1][Política académica de la Experiencia Educativa]{%
  cajaplantilla,
  colback=polfondo, colframe=polinaranja,
  coltitle=white, colbacktitle=polinaranja,
  title={#1}}

% ---------------------------------------------------------------------
% 4. NOTA NORMATIVA / CONCEPTUAL — con fuente identificable
% ---------------------------------------------------------------------
\newtcolorbox{notanormativa}[1][Referencia normativa]{%
  cajaplantilla,
  colback=normfondo, colframe=normgris,
  coltitle=white, colbacktitle=normgris,
  title={#1}}

% ---------------------------------------------------------------------
% 5. OBSERVACIÓN DE LA REVISIÓN — trazabilidad del cambio.
%    Sólo se compone en la variante anotada (make annotated). No forma
%    parte del plan: documenta por qué la plantilla es como es.
%
%    \begin{observacion}[Título corto]
%      \obsproblema{Qué estaba mal en la plantilla anterior.}
%      \obscambio{Qué se hizo en su lugar.}
%      \obsjustificacion{Por qué ese cambio resuelve el problema.}
%      \obsfuente{OBS §4.4}
%    \end{observacion}
%
%    Las claves de \obsfuente siguen el criterio de docs/MATRIZ_CAMBIOS.md:
%      OBS §n  -> Observaciones_plantilla_pruebas_ISO_29119.docx, apartado n
%      PDF §n  -> Plan_de_trabajo_29119_AngelAguilar.pdf, sección n
%      DIAG    -> análisis de los diagramas de referencia
%      PROPIA  -> criterio de esta reingeniería
% ---------------------------------------------------------------------
\newcommand{\obsitem}[2]{%
  \par\addvspace{2pt}%
  \noindent\textbf{\textcolor{obsmorado}{#1}}\enspace #2\par}
\newcommand{\obsproblema}[1]{\obsitem{Problema identificado.}{#1}}
\newcommand{\obscambio}[1]{\obsitem{Cambio aplicado.}{#1}}
\newcommand{\obsjustificacion}[1]{\obsitem{Justificación.}{#1}}
\newcommand{\obsfuente}[1]{\obsitem{Fuente.}{#1\ \textperiodcentered\ detalle completo en \texttt{docs/MATRIZ\_CAMBIOS.md}}}

\ifmostrarobservaciones
  \newtcolorbox{observacion}[1][]{%
    cajaplantilla,
    colback=obsfondo, colframe=obsmorado,
    coltitle=white, colbacktitle=obsmorado,
    title={Observación de la revisión\ifstrempty{#1}{}{: #1}}}
\else
  \excludecomment{observacion}
\fi

% ---------------------------------------------------------------------
% Menciones a las figuras
%   Los diagramas son material de apoyo y no se componen en la variante
%   final. \sidiag envuelve las frases (o los incisos) que remiten a una
%   figura, de modo que el texto siga siendo correcto cuando no está.
%   Redacte siempre la frase de manera que se sostenga sin el inciso.
% ---------------------------------------------------------------------
\newcommand{\sidiag}[1]{\ifmostrardiagramas #1\fi}

% ---------------------------------------------------------------------
% Marcador de campo por rellenar
% ---------------------------------------------------------------------
% Se compone como texto normal (no como caja) para que pueda partirse
% entre líneas dentro de las celdas estrechas de las tablas. El color y
% los corchetes bastan para localizarlo; al sustituirlo por contenido
% real, el marcador desaparece por completo.
\newcommand{\campo}[1]{{\color{campotexto}\itshape[\hspace{0pt}#1\hspace{0pt}]}}

% Marcador breve dentro de tablas (sin caja, para no romper filas)
\newcommand{\pordef}[1][\ ]{\textit{\textcolor{gray}{#1}}}

% ---------------------------------------------------------------------
% Divisor de PARTE (PARTE I / PARTE II)
% ---------------------------------------------------------------------
\newcommand{\parteplan}[2]{%
  \clearpage
  \phantomsection
  \addcontentsline{toc}{part}{#1. #2}%
  \vspace*{5\baselineskip}
  \begin{center}
    {\Large\bfseries\color{uvazul} #1}\\[10pt]
    {\LARGE\bfseries\color{uvazul} #2}
  \end{center}
  \vspace{2\baselineskip}
  \setcounter{section}{0}%
  \clearpage
}

% ---------------------------------------------------------------------
% Inicio de los anexos: la numeración pasa de 1,2,3 a ANEXO A, B, C…
% ---------------------------------------------------------------------
\newcommand{\iniciaranexos}{%
  \clearpage
  \setcounter{section}{0}%
  \renewcommand{\thesection}{ANEXO \Alph{section}}%
  \renewcommand{\thesubsection}{\Alph{section}.\arabic{subsection}}%
  % ensancha la caja del número en el índice sólo a partir de los anexos
  \addtocontents{toc}{\protect\cftsetindents{section}{1.5em}{6.4em}}%
  \phantomsection
  \addcontentsline{toc}{part}{ANEXOS}%
}

% ---------------------------------------------------------------------
% Nota al pie de una tabla (fuente / aclaración)
% ---------------------------------------------------------------------
\newcommand{\notatabla}[1]{\par\vspace{3pt}{\footnotesize\itshape #1\par}}

% ---------------------------------------------------------------------
% Encabezado de tabla sombreado
% ---------------------------------------------------------------------
\newcommand{\filaenc}{\rowcolor{uvgris}}

% ---------------------------------------------------------------------
% Abreviatura de la norma, para citarla siempre igual
% ---------------------------------------------------------------------
\newcommand{\norma}{ISO/IEC/IEEE 29119}
\newcommand{\normap}[1]{ISO/IEC/IEEE 29119-#1}
    % cajas de la plantilla y comandos propios

\hypersetup{
  pdftitle={Plan de pruebas de software},
  pdfsubject={Plantilla basada en ISO/IEC/IEEE 29119},
  pdfkeywords={pruebas de software, ISO 29119, plan de prueba, plantilla}
}

\begin{document}

% ---------------------------------------------------------------------
% Preliminares
% ---------------------------------------------------------------------
\pagenumbering{roman}

% =====================================================================
%  Portada y consideraciones generales de uso de la plantilla
% =====================================================================

\thispagestyle{uvestilo}

\begin{instruccion}[Consideraciones generales para el uso de esta plantilla]
\begin{itemize}
  \item Cuide la redacción y la ortografía; respete tipografía, interlineado y justificación.
  \item No pierda el encabezado ni el pie de página al agregar contenido.
  \item Actualice los cuadros de \emph{Historial} y \emph{Revisiones} en cada entrega.
  \item El índice, los índices de cuadros y figuras y las referencias cruzadas se generan
        solos: basta con recompilar el documento dos veces.
  \item \textbf{Todos} los cuadros y figuras deben estar referidos desde el texto
        (\texttt{\textbackslash ref\{...\}}), indicando qué información presentan y una
        breve explicación.
  \item Para producir la versión final de entrega basta con ejecutar \texttt{make final}:
        las guías verdes y los ejemplos azules desaparecen y el resto del documento no
        cambia. Para hacerlo permanente, ponga a \texttt{false} los dos interruptores
        \texttt{mostrarinstrucciones} y \texttt{mostrarejemplos} de
        \texttt{config/metadata.tex}.
  \item Sustituya cada marcador \campo{así} por el contenido que corresponda.
\end{itemize}
\end{instruccion}

\begin{notanormativa}[Cómo leer los bloques de color de esta plantilla]
\begin{itemize}
  \item \textcolor{instrverde}{\textbf{Verde}} — guía de la plantilla: explica qué debe
        escribirse. Se elimina en la versión final.
  \item \textcolor{ejemazul}{\textbf{Azul}} — ejemplo ilustrativo. No es contenido a conservar.
  \item \textcolor{polinaranja}{\textbf{Naranja}} — política académica del curso.
        Es una regla \emph{local}, no un requisito de \norma.
  \item \textcolor{normgris}{\textbf{Gris}} — referencia normativa o aclaración conceptual,
        con su fuente identificada.
\end{itemize}
Esta separación responde a una observación explícita de la revisión: el lector debe poder
distinguir de inmediato un requisito del curso de una recomendación de la norma.
\end{notanormativa}

\clearpage

% --------------------------- PORTADA ---------------------------------
\begin{center}
  \vspace*{1.5\baselineskip}

  {\large\bfseries \uvUniversidad}\\[2pt]
  {\normalsize \uvFacultad}\\[2pt]
  {\normalsize \uvPrograma}\\[2pt]
  {\normalsize \uvEE}

  \vspace{3\baselineskip}

  {\Huge\bfseries\color{uvazul} \docTitulo}

  \vspace{1.2\baselineskip}
  {\LARGE \docProyecto}

  \vspace{3\baselineskip}

  \begin{tabularx}{0.82\textwidth}{|E{3.6cm}|Y|}
    \hline
    Fecha            & \docFecha \\ \hline
    Equipo           & \docEquipo \\ \hline
    Integrantes      & \docIntegrantes \\ \hline
    Versión del documento & \docVersionDoc \\ \hline
  \end{tabularx}

  \vspace{2.5\baselineskip}

  \begin{minipage}{0.86\textwidth}\small\centering
    El presente documento se basa en la estructura de documentación de la prueba indicada en
    \normap{3}, \emph{Software and systems engineering — Software testing — Part~3: Test
    documentation}, y en el modelo de procesos de \normap{2}.
  \end{minipage}

  \vfill
\end{center}

\clearpage

% --------------------- INFORMACIÓN DEL PROYECTO ----------------------
\section*{Información del proyecto}
\phantomsection
\addcontentsline{toc}{part}{Información del proyecto}

\begin{instruccion}
Identifique el proyecto que será sometido a prueba y su situación actual. El
\emph{estatus} sirve para justificar decisiones posteriores: no se planifica igual la prueba
de un sistema terminado que la de uno en desarrollo activo.
\end{instruccion}

\begin{table}[H]
\centering
\caption{Identificación del proyecto sometido a prueba.}
\label{tab:info-proyecto}
\begin{tabularx}{\textwidth}{|E{5cm}|Y|}
  \hline
  Equipo desarrollador            & \campo{Quién construyó el software} \\ \hline
  Proyecto                        & \campo{Nombre y versión del sistema} \\ \hline
  Estatus actual del proyecto     & \campo{En desarrollo / entregado / en mantenimiento} \\ \hline
  Experiencia Educativa           & \uvEE \\ \hline
  Equipo de prueba                & \docEquipo \\ \hline
  Periodo del esfuerzo de prueba  & \campo{dd/mm/aaaa – dd/mm/aaaa} \\ \hline
\end{tabularx}
\end{table}

El Cuadro~\ref{tab:info-proyecto} identifica el sistema que será probado, quién lo desarrolló
y en qué estado se encuentra al iniciar el esfuerzo de prueba; estos datos delimitan lo que el
resto del plan puede asumir como disponible.

% =====================================================================
%  Control de versiones del documento, revisiones y aprobaciones
% =====================================================================
\clearpage
\section*{Control de versiones}
\phantomsection
\addcontentsline{toc}{part}{Control de versiones}

\begin{instruccion}
Registre aquí la evolución de \emph{este documento}, distinguiendo tres cosas: el
\textbf{historial} de cambios (quién modificó qué), la \textbf{revisión} (alguien leyó el
documento y emitió observaciones) y la \textbf{aprobación formal} (alguien con autoridad acepta
el documento como válido). Cada una tiene su propio cuadro.
\end{instruccion}

\begin{notanormativa}[Alcance de este apartado]
Este control de versiones cubre únicamente el \emph{documento}. El control de versiones de los
demás activos de prueba —software bajo prueba, casos, procedimientos, guiones, datos y
configuración del entorno— se define en el apartado~\ref{sec:configuracion} y se registra en
cada ejecución (\ref{anx:ejecuciones}).
\end{notanormativa}

\begin{observacion}[Revisión y aprobación no son lo mismo]
\obsproblema{La plantilla anterior resolvía las revisiones con un único cuadro cuya columna
«Estatus» mezclaba el resultado de una revisión con la aprobación formal del documento, y
limitaba todo el control de versiones al propio archivo Word.}
\obscambio{El apartado se divide en historial, revisiones y aprobación formal; el control de
versiones del resto de los activos pasa al apartado~\ref{sec:configuracion}.}
\obsjustificacion{Una revisión puede cerrarse con observaciones sin que el documento quede
aprobado, y un resultado de prueba no es reproducible si sólo se versiona el documento que lo
describe.}
\obsfuente{OBS §2 y §3}
\end{observacion}

% ------------------------------ Historial ----------------------------
\subsection*{Historial de cambios}

\begin{table}[H]
\centering
\caption{Historial de cambios del documento.}
\label{tab:historial}
\begin{tabularx}{\textwidth}{|C{2.5cm}|C{1.8cm}|L{3.4cm}|Y|}
  \hline
  \filaenc \textbf{Fecha} & \textbf{Versión} & \textbf{Autor} & \textbf{Descripción del cambio} \\ \hline
  \campo{dd/mm/aaaa} & \campo{0.1} & \campo{Nombre} & \campo{Qué se agregó, modificó o eliminó} \\ \hline
  & & & \\ \hline
  & & & \\ \hline
  & & & \\ \hline
\end{tabularx}
\end{table}

El Cuadro~\ref{tab:historial} deja constancia de cada versión emitida del plan; debe
actualizarse \emph{antes} de cada entrega, no después.

% ------------------------------ Revisiones ---------------------------
\subsection*{Revisiones}

\begin{table}[H]
\centering
\caption{Revisiones realizadas al documento.}
\label{tab:revisiones}
\begin{tabularx}{\textwidth}{|C{2.5cm}|C{1.5cm}|L{2.9cm}|Y|C{2.4cm}|}
  \hline
  \filaenc \textbf{Fecha} & \textbf{Versión} & \textbf{Revisor} & \textbf{Observaciones} & \textbf{Resultado} \\ \hline
  \campo{dd/mm/aaaa} & \campo{0.1} & \campo{Nombre y rol} & \campo{Hallazgos de la revisión} & \campo{Ver nota} \\ \hline
  & & & & \\ \hline
  & & & & \\ \hline
\end{tabularx}
\notatabla{Resultado: \textit{Sin observaciones} · \textit{Con observaciones menores} · \textit{Con observaciones mayores}. La revisión no equivale a la aprobación.}
\end{table}

El Cuadro~\ref{tab:revisiones} registra quién leyó el documento y qué observó. Una revisión
puede cerrarse sin que el documento quede aprobado.

% ----------------------------- Aprobaciones --------------------------
\subsection*{Aprobación formal}

\begin{table}[H]
\centering
\caption{Aprobación formal del documento.}
\label{tab:aprobaciones}
\begin{tabularx}{\textwidth}{|C{2.5cm}|C{1.5cm}|L{3.4cm}|Y|C{2.2cm}|}
  \hline
  \filaenc \textbf{Fecha} & \textbf{Versión} & \textbf{Autoridad que aprueba} & \textbf{Ámbito de la aprobación} & \textbf{Estatus} \\ \hline
  \campo{dd/mm/aaaa} & \campo{1.0} & \autoridadAprobacion & \campo{Plan completo / apartado} & \campo{Ver nota} \\ \hline
  & & & & \\ \hline
\end{tabularx}
\notatabla{Estatus: \textit{Aprobado} · \textit{Aprobado con correcciones} · \textit{No aprobado}. Corresponde a la tarea TP8 «Lograr consenso sobre el plan» de \normap{2}.}
\end{table}

El Cuadro~\ref{tab:aprobaciones} identifica a quién autorizó el plan y sobre qué versión lo
hizo; sin esa firma, el plan permanece en borrador y no debería dirigir la ejecución.

\begin{politicaacademica}
Durante el curso, la autoridad de aprobación del plan y de cualquier desviación respecto de la
estrategia acordada es \autoridadAprobacion.\ifversionfinal\else\ Al reutilizar la plantilla
fuera del curso, sustituya ese valor en \texttt{config/metadata.tex} por el rol que corresponda
en la organización (líder de pruebas, responsable de calidad, patrocinador del proyecto).\fi
\end{politicaacademica}


\clearpage
\phantomsection
\addcontentsline{toc}{part}{Tabla de contenido}
\tableofcontents

\clearpage
\phantomsection
\addcontentsline{toc}{part}{Índice de cuadros y figuras}
\listoftables
\vspace{1\baselineskip}
\listoffigures

% ---------------------------------------------------------------------
% PARTE I. PLAN DE PRUEBA
% ---------------------------------------------------------------------
\clearpage
\pagenumbering{arabic}

\parteplan{PARTE I}{PLAN DE PRUEBA}

% =====================================================================
%  1. Introducción
% =====================================================================
\section{Introducción}

\subsection{Presentación del documento}

\begin{instruccion}
Explique el objetivo del plan de prueba, su estructura y, de forma general, el contenido de
cada apartado. Indique además a quién va dirigido el documento y en qué momento del proyecto
se emite. Conviene señalar explícitamente qué partes de \norma\ se han tomado como referencia
y cuáles no, para no dar a entender una conformidad que no se ha demostrado.
\end{instruccion}

\campo{Redacte aquí la presentación del documento}

\begin{notanormativa}[Relación entre las partes de la norma y este documento]
\normap{2} describe los \emph{procesos} de prueba; \normap{3}, la \emph{documentación} que
esos procesos producen; \normap{4}, las \emph{técnicas} de diseño de pruebas. Por lo tanto,
el contenido de este plan \textbf{no sustituye} a los registros que se generan al diseñar,
ejecutar, controlar y cerrar las pruebas: los complementa.\sidiag{ La
Figura~\ref{fig:estructura} muestra esa relación y señala qué apartados de este documento se
apoyan en cada parte.}
\end{notanormativa}

\ifmostrardiagramas
% =====================================================================
%  diagrams/estructura-29119.tex
%  Figura: estructura general de ISO/IEC/IEEE 29119 y su relación
%          con las partes de este documento.
%
%  CORRECCIONES respecto del diagrama original («normal general.png»):
%   1. El original colgaba «Procesos organizacionales · política y
%      estrategia» de la Parte 1. Los procesos —incluido el
%      organizacional— pertenecen a la Parte 2; la Parte 1 aporta
%      conceptos generales y vocabulario.
%   2. Se elimina toda sugerencia de que una parte sustituye o contiene
%      a otra: son partes complementarias de una misma norma.
%   3. Se añade la correspondencia explícita con las secciones y anexos
%      de este documento, que el original no representaba.
%
%  Nota de mantenimiento: los nodos se colocan con «positioning»
%  (below=… of …) en lugar de coordenadas absolutas, de modo que al
%  cambiar el texto de un cuadro los demás se desplacen y nunca se
%  superpongan.
% =====================================================================
\begin{figure}[htbp]
\centering
\begin{tikzpicture}[
  font=\scriptsize,
  node distance=4mm,
  raiz/.style   ={draw=uvazul, very thick, fill=uvazul!12, rounded corners=3pt,
                  text width=5.0cm, align=center, inner sep=4pt},
  parte/.style  ={draw=uvazul, thick, fill=uvazul!7, rounded corners=2pt,
                  text width=2.35cm, align=center, inner sep=4pt},
  sub/.style    ={draw=gray!70, fill=gray!8, rounded corners=2pt,
                  text width=2.35cm, align=center, inner sep=4pt},
  doc/.style    ={draw=instrverde, thick, fill=instrverde!8, rounded corners=2pt,
                  text width=3.7cm, align=center, inner sep=4pt},
  fl/.style     ={-{Latex[length=4pt]}, draw=uvazul!80},
  mapa/.style   ={-{Latex[length=4pt]}, draw=instrverde!85, dashed}
]

% ------------------------------- raíz --------------------------------
\node[raiz] (n) {\textbf{ISO/IEC/IEEE 29119}\\ Prueba de software (norma multiparte)};

% ------------------------------ partes -------------------------------
\node[parte, below=10mm of n] (p3) {\textbf{Parte 3}\\ Documentación de prueba};
\node[parte, left=6mm of p3]  (p2) {\textbf{Parte 2}\\ Procesos de prueba};
\node[parte, left=6mm of p2]  (p1) {\textbf{Parte 1}\\ Conceptos generales y vocabulario};
\node[parte, right=6mm of p3] (p4) {\textbf{Parte 4}\\ Técnicas de prueba};
\node[parte, right=6mm of p4] (p5) {\textbf{Parte 5}\\ Pruebas por palabras clave};

\foreach \d in {p1,p2,p3,p4,p5} { \draw[fl] (n) -- (\d); }

% --------------------------- contenido -------------------------------
% Parte 1
\node[sub, below=of p1] (v1) {Vocabulario y conceptos que usan las demás partes};
\draw[fl] (p1) -- (v1);

% Parte 2: las tres capas de procesos
\node[sub, below=of p2] (c1) {Capa organizacional};
\node[sub, below=of c1] (c2) {Gestión de la prueba};
\node[sub, below=of c2] (c3) {Procesos dinámicos};
\draw[fl] (p2) -- (c1); \draw[fl] (c1) -- (c2); \draw[fl] (c2) -- (c3);

% Parte 3: tipos de documentación
\node[sub, below=of p3] (d1) {Documentos de gestión};
\node[sub, below=of d1] (d2) {Especificaciones de nivel};
\node[sub, below=of d2] (d3) {Registros de ejecución e incidentes};
\draw[fl] (p3) -- (d1); \draw[fl] (d1) -- (d2); \draw[fl] (d2) -- (d3);

% Parte 4: familias de técnicas
\node[sub, below=of p4] (t1) {Basadas en especificación};
\node[sub, below=of t1] (t2) {Basadas en estructura};
\node[sub, below=of t2] (t3) {Basadas en experiencia};
\draw[fl] (p4) -- (t1); \draw[fl] (t1) -- (t2); \draw[fl] (t2) -- (t3);

% Parte 5
\node[sub, below=of p5] (k1) {Enfoque opcional; se usa sólo si el proyecto lo adopta};
\draw[fl] (p5) -- (k1);

% ------------------- correspondencia con este documento --------------
\node[fit=(v1)(c3)(d3)(t3)(k1), inner sep=0pt, draw=none] (todos) {};

\node[align=center, font=\scriptsize\bfseries, text=instrverde,
      below=7mm of todos.south] (rot) {ESTE DOCUMENTO};

\node[doc, below=3mm of rot] (e2) {\textbf{Parte II}\\ Reporte de finalización};
\node[doc, left=5mm of e2]   (e1) {\textbf{Parte I}\\ Plan de prueba (secciones 1--9)};
\node[doc, right=5mm of e2]  (e3) {\textbf{Anexos A--G}\\ casos, procedimientos, ejecuciones, incidentes, informes, entorno, trazabilidad};

\draw[mapa] (c2.west) to[out=190,in=100] (e1.north);
\draw[mapa] (d1.east) to[out=0,in=90]    (e2.north east);
\draw[mapa] (d3.east) to[out=350,in=110] (e3.north west);
\draw[mapa] (t3.south) to[out=280,in=80] (e3.north);

% ------------------------------ leyenda ------------------------------
\node[align=left, font=\scriptsize, text width=15.2cm, below=5mm of e2] (leg)
  {\textbf{Cómo leer la figura.} Las flechas azules indican \emph{qué contiene} cada parte de la
   norma; las verdes discontinuas, \emph{dónde} se refleja ese contenido en este documento.
   Ninguna parte sustituye a otra: el plan no reemplaza a los registros que producen los procesos
   dinámicos, ni las técnicas de la Parte~4 son obligatorias en su totalidad.};

\end{tikzpicture}
\caption[Estructura de ISO/IEC/IEEE 29119 y su relación con este documento]%
{Estructura general de \norma\ y correspondencia de sus partes con las secciones y anexos de
este documento.}
\label{fig:estructura}
\end{figure}

\begin{observacion}[Diagrama de estructura de la norma]
\obsproblema{El diagrama de referencia colgaba «Procesos organizacionales · política y
estrategia» de la Parte~1, y sugería una jerarquía en la que unas partes sustituían a otras.}
\obscambio{Los procesos —incluido el organizacional— cuelgan de la Parte~2; la Parte~1 aporta
conceptos y vocabulario. Se añade la correspondencia con las secciones y anexos de este
documento y una leyenda que aclara que ninguna parte reemplaza a otra.}
\obsfuente{OBS §1 · DIAG}
\end{observacion}


La Figura~\ref{fig:estructura} presenta la relación conceptual entre las partes de \norma\ y
las secciones y anexos de esta plantilla. Debe leerse como un mapa de correspondencias, no
como una jerarquía de sustitución: ninguna parte de la norma reemplaza a otra, y disponer de
este plan no acredita por sí solo la aplicación de los procesos que la norma describe.

% =====================================================================
%  diagrams/procesos-prueba.tex
%  Figura: flujo del proceso de prueba en tres capas, con sus bucles
%          de retroalimentación.
%
%  CORRECCIONES respecto del diagrama original («Procesos de prueba.png»):
%   1. El original rotulaba «verificar correcciones» como prueba de
%      regresión. Aquí se representan por separado la prueba de
%      CONFIRMACIÓN (el defecto quedó corregido) y la de REGRESIÓN (el
%      cambio no dañó lo que ya funcionaba).
%   2. El original se leía como una secuencia lineal. Aquí el monitoreo
%      y control aparece como proceso CONCURRENTE y se dibujan los tres
%      bucles de retroalimentación: control, replanificación y mejora.
%   3. El original situaba el entorno y los datos dentro de la
%      preparación de pruebas; en la edición 2021 constituyen un proceso
%      propio (ES1/ES2), y así se representa.
%   4. El original sólo contemplaba el análisis de riesgos del proyecto;
%      la planificación abarca riesgos de producto y de proyecto.
% =====================================================================
\begin{figure}[htbp]
\centering
\begin{tikzpicture}[
  font=\scriptsize,
  proc/.style   ={draw=uvazul, thick, fill=uvazul!8, rounded corners=2pt,
                  text width=2.8cm, align=center, minimum height=1.0cm, inner sep=3pt},
  procm/.style  ={draw=polinaranja, very thick, fill=polfondo, rounded corners=2pt,
                  text width=2.9cm, align=center, minimum height=1.5cm, inner sep=3pt},
  salida/.style ={draw=gray!70, fill=gray!8, rounded corners=2pt,
                  text width=2.6cm, align=center, minimum height=0.7cm, inner sep=3pt},
  conf/.style   ={draw=instrverde, thick, fill=instrverde!8, rounded corners=2pt,
                  text width=2.7cm, align=center, minimum height=0.9cm, inner sep=3pt},
  fl/.style     ={-{Latex[length=4pt]}, draw=uvazul!85, thick},
  fb/.style     ={-{Latex[length=4pt]}, draw=polinaranja, thick, dashed},
  mej/.style    ={-{Latex[length=4pt]}, draw=instrverde, thick, dashdotted},
  et/.style     ={font=\tiny, inner sep=1pt, fill=white}
]

% ============================ CAPA ORGANIZACIONAL =====================
\node[proc] (ot) at (-4.6,0) {\textbf{Proceso organizacional}\\ OT1 desarrollar · OT2 monitorear · OT3 actualizar};
\node[salida] (pol) at (-0.6,0) {Política de prueba\\ y estrategia organizacional};
\draw[fl] (ot) -- (pol);

% ============================ CAPA DE GESTIÓN =========================
\node[proc] (tp) at (-5.6,-2.6) {\textbf{Planificación}\\ TP1--TP9\\ \emph{riesgos de producto y de proyecto}};
\node[salida] (plan) at (-5.6,-4.3) {Plan de prueba};
\draw[fl] (tp) -- (plan);

\node[procm] (tmc) at (-0.4,-3.3) {\textbf{Monitoreo y control}\\ TMC1--TMC4\\ \emph{concurrente con la ejecución}};

\node[proc] (tc) at (4.6,-3.3) {\textbf{Finalización}\\ TC1 archivar · TC2 limpiar entorno · TC3 lecciones · TC4 reportar};

\draw[fl] (pol) to[out=250,in=70] node[et,pos=0.5]{gobierna} (tp.north);

% ============================ CAPA DINÁMICA ===========================
\node[proc] (es) at (-6.0,-6.6) {\textbf{Entorno y datos}\\ ES1 establecer\\ ES2 mantener};
\node[proc] (td) at (-2.6,-6.6) {\textbf{Diseño e implementación}\\ TD1--TD6};
\node[proc] (te) at (0.8,-6.6)  {\textbf{Ejecución}\\ TE1 ejecutar · TE2 comparar · TE3 registrar};
\node[proc] (ir) at (4.2,-6.6)  {\textbf{Reporte de incidentes}\\ IR1 analizar · IR2 reportar};

\draw[fl] (es) -- (td);
\draw[fl] (td) -- (te);
\draw[fl] (te) -- node[et,pos=0.5]{no aprobado} (ir);

\draw[fl] (plan.south) to[out=270,in=120] node[et,pos=0.25]{dirige} (es.north);

% ------------- bucle interno: corrección, confirmación, regresión -----
\node[salida] (nv) at (4.2,-8.5) {Nueva versión del elemento de prueba};
\draw[fl] (ir) -- node[et,pos=0.5]{corrección} (nv);

\node[conf] (pc) at (0.3,-8.5)  {\textbf{Prueba de confirmación}\\ ¿se corrigió \emph{ese} defecto?};
\node[conf] (pr) at (0.3,-10.4) {\textbf{Prueba de regresión}\\ ¿el cambio dañó lo que ya funcionaba?};
\draw[fl] (nv.west) -- (pc.east);
\draw[fl] (nv.south) to[out=250,in=340] (pr.east);

\draw[fl] (pc.north) -- node[et,pos=0.5]{nueva ejecución} (te.south);
\draw[fl] (pr.west) to[out=180,in=250,looseness=1.2] (te.south west);

% ======================= BUCLES DE RETROALIMENTACIÓN ==================
% 1) bucle de control corto: dinámica -> monitoreo -> dinámica
\draw[fb] (te.north) to[out=100,in=280] node[et,pos=0.5]{medidas y estado} (tmc.south);
\draw[fb] (tmc.west) to[out=200,in=80] node[et,pos=0.55]{directivas de control} (td.north);

% 2) bucle de replanificación: monitoreo -> planificación
\draw[fb] (tmc.north west) to[out=150,in=20] node[et,pos=0.5]{replanificación} (tp.east);

% 3) bucle de mejora organizacional: finalización -> capa organizacional
% se rodea por encima de la capa organizacional para no atravesar ningún nodo
\draw[mej] (tc.north) |- ($(ot.north)+(0,1.25)$)
             node[et,pos=0.72]{lecciones aprendidas} -- (ot.north);

% cierre hacia finalización
\draw[fl] (ir.east) to[out=20,in=250] node[et,pos=0.5]{criterios de salida} (tc.south);

% ============================ BANDAS DE CAPA ==========================
\begin{scope}[on background layer]
  \node[fit=(ot)(pol), draw=gray!45, dashed, inner sep=7pt, rounded corners] (b1) {};
  \node[fit=(tp)(plan)(tmc)(tc), draw=gray!45, dashed, inner sep=7pt, rounded corners] (b2) {};
  \node[fit=(es)(td)(te)(ir)(nv)(pc)(pr), draw=gray!45, dashed, inner sep=7pt, rounded corners] (b3) {};
\end{scope}
\node[et, text=gray!60, anchor=west, font=\tiny\bfseries] at (b1.north west) {CAPA ORGANIZACIONAL};
\node[et, text=gray!60, anchor=west, font=\tiny\bfseries] at (b2.north west) {CAPA DE GESTIÓN DE LA PRUEBA};
\node[et, text=gray!60, anchor=west, font=\tiny\bfseries] at (b3.north west) {CAPA DINÁMICA (iterativa)};

% ============================== LEYENDA ===============================
\node[align=left, font=\scriptsize, text width=14.8cm] at (0,-12.4)
  {\textbf{Cómo leer la figura.} Las flechas azules continuas indican el flujo principal; las
   naranjas discontinuas, los bucles de control y de replanificación; la verde de trazo mixto,
   el bucle lento de mejora organizacional. El monitoreo y control \emph{no} ocupa un lugar fijo
   en la secuencia: acompaña a la ejecución de principio a fin. Los pasos de la capa dinámica se
   repiten por conjunto de pruebas, incremento o iteración hasta cumplir los criterios de salida.};

\end{tikzpicture}
\caption[Flujo del proceso de prueba con sus bucles de retroalimentación]%
{Flujo del proceso de prueba en las tres capas de \normap{2}, con el monitoreo y control
representado como proceso concurrente y con los bucles de control, replanificación y mejora.}
\label{fig:procesos}
\end{figure}

\begin{observacion}[Diagrama del flujo del proceso]
\obsproblema{El diagrama de referencia rotulaba «verificar correcciones» como prueba de
regresión, se leía como una secuencia lineal, situaba el entorno y los datos dentro de la
preparación de pruebas y contemplaba sólo el análisis de riesgos del proyecto.}
\obscambio{Confirmación y regresión aparecen como cajas distintas; el monitoreo y control se
dibuja como proceso concurrente con tres bucles de retroalimentación explícitos; entorno y datos
constituyen un proceso propio (ES1/ES2); la planificación abarca riesgos de producto y de
proyecto.}
\obsfuente{OBS §5.5 y §5.6 · PDF §2.2 · DIAG}
\end{observacion}


La Figura~\ref{fig:procesos} resume el flujo del proceso de prueba en las tres capas que define
\normap{2} —organizacional, de gestión y dinámica— y, sobre todo, los bucles de
retroalimentación que impiden leerlo como una cascada: el seguimiento y control opera de forma
\emph{concurrente} con la ejecución, los incidentes y los riesgos nuevos pueden reabrir la
planificación, y las lecciones aprendidas del cierre retroalimentan la capa organizacional. Este
plan corresponde a la capa de gestión; los anexos documentan lo que ocurre en la capa dinámica.
\fi

\begin{notanormativa}[\segunvariante{Declaración de alcance de esta plantilla}{Declaración de alcance de este plan de prueba}]
\segunvariante{Esta plantilla}{Este plan} \textbf{no declara conformidad} con \norma. Organiza
la planificación y el cierre de un esfuerzo de prueba tomando como referencia el modelo de
procesos de la parte~2 y la estructura documental de la parte~3. Demostrar conformidad exigiría,
además, evidencia de que las actividades se realizaron y de que sus resultados fueron los
especificados, y una revisión del texto normativo completo de la edición elegida.
\end{notanormativa}

% ---------------------------------------------------------------------
\subsection{Referentes utilizados}

\begin{instruccion}
Relacione los documentos que sirvieron de base al plan: referencias normativas y documentación
técnica del proyecto a probar. Incluya como mínimo nombre, descripción y, para documentos
digitales, la URL y el nombre del archivo con su extensión.

Cite siempre la \emph{edición} (año) de cada norma y la \emph{versión} de cada fuente técnica.
Una referencia como «\normap{4}, Anexo~A, pp.~32--38» es ambigua sin el año, porque la
paginación cambia entre ediciones. Cuando cite una cláusula o anexo, verifíquelo en la edición
que realmente consultó.
\end{instruccion}

\begin{observacion}[Ediciones normativas]
\obsproblema{La plantilla anterior remitía a «\normap{4}, Anexo~A, páginas 32--38» sin indicar
la edición, y no pedía versionar las fuentes técnicas del proyecto.}
\obscambio{El cuadro de referentes exige edición o versión y asigna una clave que se reutiliza
para identificar la base de prueba.}
\obsjustificacion{La paginación y la numeración de anexos cambian entre ediciones, de modo que
la referencia original no es verificable.}
\obsfuente{OBS §2 y §5.2}
\end{observacion}

\begin{table}[H]
\centering
\caption{Referentes normativos y técnicos utilizados en la elaboración del plan.}
\label{tab:referentes}
\begin{tabularx}{\textwidth}{|C{1.3cm}|L{4.6cm}|Y|C{2.1cm}|}
  \hline
  \filaenc \textbf{Clave} & \textbf{Documento (con edición/versión)} & \textbf{Descripción y uso en el plan} & \textbf{Ubicación} \\ \hline
  \campo{REF-01} & \campo{Norma, edición y cláusula verificada} & \campo{Para qué se usó} & \campo{URL/archivo} \\ \hline
  \campo{REF-02} & \campo{Especificación de requisitos v1.2} & \campo{Base de prueba} & \campo{URL/archivo} \\ \hline
  & & & \\ \hline
  & & & \\ \hline
\end{tabularx}
\notatabla{Distinga los referentes \textit{normativos} (normas y estándares) de los \textit{técnicos} (documentación del proyecto a probar). Ambos son citables desde los apartados posteriores mediante su clave.}
\end{table}

El Cuadro~\ref{tab:referentes} concentra las fuentes en que se apoya el plan y les asigna una
clave; esa clave se reutiliza en el apartado~\ref{sec:objeto-base} para identificar cuáles de
ellas constituyen la \emph{base de prueba}.

% ---------------------------------------------------------------------
\subsection{Glosario}

\begin{instruccion}
Defina los términos, acrónimos y abreviaturas relevantes cuyo significado deba explicitarse.
Incluya los términos del dominio de la aplicación que faciliten comprender el caso de estudio.
Se recomienda separar acrónimos y términos en dos subapartados, en formato de lista tipo
diccionario. Indique la fuente de las definiciones tomadas de una norma o de un cuerpo de
conocimiento.
\end{instruccion}

\subsubsection*{Acrónimos}

\begin{description}[style=nextline,leftmargin=2.6cm,labelwidth=2.4cm,font=\bfseries]
  \item[\campo{SIL}] \campo{Software Integrity Level. Nivel de integridad del software; véase el apartado~\ref{sec:criticidad}.}
  \item[\campo{ACR-02}] \campo{Significado}
\end{description}

\subsubsection*{Términos}

\begin{instruccion}
Los términos siguientes se ofrecen ya redactados porque son los que con más frecuencia se
confunden entre sí al documentar pruebas. Consérvelos, ajústelos si su proyecto usa otra
terminología y añada los términos propios del dominio de la aplicación.
\end{instruccion}

\begin{observacion}[Términos predefinidos del glosario]
\obsproblema{La plantilla anterior pedía un glosario sin indicar qué términos convenía definir,
y a lo largo del documento empleaba «elemento de prueba», «cobertura» y «regresión» con
significados que se solapaban.}
\obscambio{Se predefinen los ocho términos cuya confusión originó correcciones en el resto del
documento.}
\obsfuente{OBS §4.3, §4.4 y §5.5}
\end{observacion}

\begin{description}[style=nextline,leftmargin=2.6cm,labelwidth=2.4cm,font=\bfseries]
  \item[Objeto de prueba] Elemento que se somete a prueba: sistema, subsistema, módulo,
        componente, servicio, API o interfaz. Responde a \emph{qué se prueba}.
  \item[Base de prueba] Conjunto de documentos y fuentes de los que se deriva lo que debe
        verificarse: requisitos, historias de usuario, casos de uso, reglas de negocio,
        arquitectura o diseño. Responde a \emph{contra qué se juzga el objeto}.
  \item[Condición de prueba] Aspecto verificable del objeto de prueba, derivado de la base de
        prueba y de los riesgos, a partir del cual se diseñan casos.
  \item[Cobertura] Proporción de elementos de un conjunto identificado —requisitos,
        particiones, decisiones, transiciones— que han quedado \emph{ejercitados} por las
        pruebas. Se mide con independencia del veredicto obtenido: un elemento probado por un
        caso que falló está cubierto. No debe confundirse con el \emph{avance de ejecución}
        (cuánto del trabajo previsto se ha realizado) ni con el \emph{veredicto} (si lo
        ejercitado se comportó como se esperaba). Véase el apartado~\ref{sec:cobertura}.
  \item[Veredicto] Juicio emitido al comparar el resultado real de una ejecución con el
        esperado: aprobado, no aprobado o no concluyente. Pertenece a la ejecución, no al caso
        de prueba.
  \item[Defecto] Imperfección en un producto de trabajo que puede hacer que no cumpla sus
        requisitos. Es estático y reside en el artefacto.
  \item[Fallo] Evento en el que el servicio entregado se desvía del correcto; es la
        manifestación externa observable de un defecto activado.
  \item[Incidente de prueba] Evento ocurrido durante la prueba que requiere investigación o
        acción adicional. Es un registro documental: un incidente puede deberse a un defecto
        del objeto, del entorno, de los datos o del propio caso de prueba.
  \item[\campo{Término del dominio}] \campo{Definición}
\end{description}

\begin{notanormativa}[Fuente de las definiciones anteriores]
Las definiciones de defecto, fallo e incidente siguen la delimitación empleada en
ISO/IEC/IEEE 24765 y en \normap{1}, recogida en la sección~3.1 del documento
\emph{Plan de trabajo 29119} incluido en \texttt{sources/}. Un mismo defecto puede producir
varios fallos distintos y un mismo fallo puede tener varios defectos como causa.
\end{notanormativa}

% =====================================================================
%  2. Contexto del esfuerzo de prueba   (actividad TP1 de ISO/IEC/IEEE 29119-2)
% =====================================================================
\section{Contexto}

\begin{notanormativa}[Correspondencia con el modelo de procesos]
Este apartado desarrolla la actividad \textbf{TP1 «Comprender el contexto»} de \normap{2}:
identificar partes interesadas, analizar el elemento de prueba y la base de prueba, consultar
la política y estrategia aplicables y registrar las restricciones de tiempo, recursos y
regulación.
\end{notanormativa}

% ---------------------------------------------------------------------
\subsection{Proyecto caso de estudio}

\begin{instruccion}
Dé antecedentes del proyecto a probar: autores originales, objetivos, alcances y límites,
plataformas en que está desarrollado, funcionalidades generales que aborda y cualquier otra
información pertinente. Debe ser lo bastante explícito para que un lector externo entienda
qué hace el sistema sin consultar otros documentos.
\end{instruccion}

\campo{Describa aquí el proyecto caso de estudio}

% ---------------------------------------------------------------------
\subsection{Objetivo de la prueba}

\begin{instruccion}
Indique el objetivo general \textbf{de la prueba}, no del sistema. Un objetivo de prueba
enuncia qué se pretende averiguar o demostrar y para qué decisión sirve; no describe las
funciones del software.

Formule también objetivos específicos, cada uno asociado a los riesgos que pretende reducir
(apartado~\ref{sec:riesgos}). Esa asociación es la que permite justificar después la estrategia.
\end{instruccion}

\begin{ejemplo}
\textbf{Objetivo de prueba (correcto):} «Determinar si el módulo de facturación calcula
correctamente totales, impuestos y descuentos bajo las combinaciones de reglas de negocio
vigentes, con el fin de sustentar la decisión de liberación.»

\textbf{Objetivo del sistema (incorrecto en este apartado):} «Permitir a los usuarios emitir
facturas electrónicas.»
\end{ejemplo}

\campo{Objetivo general de la prueba}

\begin{table}[H]
\centering
\caption{Objetivos específicos de la prueba y riesgos que pretenden reducir.}
\label{tab:objetivos}
\begin{tabularx}{\textwidth}{|C{1.6cm}|Y|C{2.2cm}|C{3.2cm}|}
  \hline
  \filaenc \textbf{Clave} & \textbf{Objetivo específico de la prueba} & \textbf{Riesgo(s) asociado(s)} & \textbf{Cómo se evidenciará} \\ \hline
  \campo{OBJ-01} & \campo{Qué se pretende averiguar} & \campo{RP-01} & \campo{Registro o métrica} \\ \hline
  & & & \\ \hline
  & & & \\ \hline
\end{tabularx}
\end{table}

El Cuadro~\ref{tab:objetivos} vincula cada objetivo específico con los riesgos del
apartado~\ref{sec:riesgos} y con la evidencia que permitirá comprobar si el objetivo se
alcanzó; sin esa columna, el objetivo no es verificable al cierre.

% ---------------------------------------------------------------------
\subsection{Objeto de prueba y base de prueba}
\label{sec:objeto-base}

\begin{notanormativa}[Corrección respecto de la plantilla original]
La plantilla original equiparaba los «elementos de prueba» con los \emph{casos de uso}. Son
cosas distintas y conviene registrarlas por separado:
\begin{itemize}
  \item El \textbf{objeto de prueba} es lo que se somete a prueba: sistema, subsistema,
        módulo, componente, servicio, API o interfaz.
  \item La \textbf{base de prueba} es la fuente contra la que se juzga ese objeto: requisitos,
        historias de usuario, casos de uso, reglas de negocio, arquitectura o diseño.
\end{itemize}
Un caso de uso es, casi siempre, base de prueba y no objeto de prueba. Separarlos permite
aplicar esta plantilla a proyectos que no estén documentados mediante casos de uso.
\end{notanormativa}

\begin{instruccion}
Identifique primero los objetos bajo prueba con su versión exacta, y después la base de prueba
que se usará para derivar las condiciones de prueba. Si alguna parte del sistema carece de base
documentada, dígalo: es una restricción real que condiciona las técnicas aplicables
(apartado~\ref{sec:tecnicas}) y suele constituir un riesgo de proyecto.
\end{instruccion}

\begin{table}[H]
\centering
\caption{Objetos bajo prueba, con su versión y criticidad asignada.}
\label{tab:objetos}
\begin{tabularx}{\textwidth}{|C{1.6cm}|L{3.4cm}|C{2.1cm}|C{1.9cm}|Y|}
  \hline
  \filaenc \textbf{Clave} & \textbf{Objeto de prueba} & \textbf{Tipo} & \textbf{Versión} & \textbf{Criticidad asignada (\ref{sec:criticidad})} \\ \hline
  \campo{OP-01} & \campo{Módulo de facturación} & \campo{Módulo} & \campo{v1.2.0} & \campo{Crítico} \\ \hline
  \campo{OP-02} & \campo{API de autenticación} & \campo{Servicio} & \campo{v1.2.0} & \campo{Catastrófico} \\ \hline
  & & & & \\ \hline
\end{tabularx}
\notatabla{Tipo sugerido: sistema · subsistema · módulo · componente · servicio · API · interfaz de usuario · interfaz externa · base de datos.}
\end{table}

\begin{table}[H]
\centering
\caption{Base de prueba utilizada para derivar las condiciones de prueba.}
\label{tab:base-prueba}
\begin{tabularx}{\textwidth}{|C{1.6cm}|L{3.9cm}|C{2.3cm}|C{2.0cm}|Y|}
  \hline
  \filaenc \textbf{Clave} & \textbf{Elemento de la base} & \textbf{Tipo} & \textbf{Referente (\ref{tab:referentes})} & \textbf{Objeto(s) que sustenta} \\ \hline
  \campo{BP-01} & \campo{RF-12 Cálculo de total} & \campo{Requisito} & \campo{REF-02} & \campo{OP-01} \\ \hline
  \campo{BP-02} & \campo{CU-03 Emitir factura} & \campo{Caso de uso} & \campo{REF-02} & \campo{OP-01} \\ \hline
  & & & & \\ \hline
\end{tabularx}
\notatabla{Tipo sugerido: requisito funcional · requisito no funcional · historia de usuario · caso de uso · regla de negocio · documento de arquitectura · especificación de interfaz · norma aplicable.}
\end{table}

Los Cuadros~\ref{tab:objetos} y~\ref{tab:base-prueba} responden a dos preguntas distintas:
\emph{qué se prueba} y \emph{contra qué se juzga}. Sus claves son el punto de partida de la
matriz de trazabilidad del \ref{anx:trazabilidad}, que encadena base de prueba, riesgo,
condición, caso, ejecución, evidencia e incidente.

% ---------------------------------------------------------------------
\subsection{Criticidad asignada a los objetos de prueba}
\label{sec:criticidad}

\begin{politicaacademica}[Política académica: escala de criticidad (SIL)]
El curso emplea una escala de cuatro niveles para clasificar la criticidad de los objetos de
prueba: \textbf{Catastrófico}, \textbf{Crítico}, \textbf{Marginal} y \textbf{Despreciable}.
Esta escala es una \emph{regla local de la Experiencia Educativa}. Describa en el
Cuadro~\ref{tab:sil} el significado que su equipo da a cada nivel, quién lo asigna y con qué
criterio.
\end{politicaacademica}

\begin{notanormativa}[Por qué esta escala no se atribuye a \norma]
La revisión señaló que las etiquetas anteriores se presentaban como si procedieran de \norma.
No es así: \norma\ no define esa escala de cuatro niveles. El marco internacional para niveles
de integridad es \textbf{ISO/IEC/IEEE 15026-3} (\emph{Systems and software assurance —
Integrity levels}), que exige definir los niveles, los requisitos asociados a cada uno y el
método de asignación. Mientras esos elementos no se definan y verifiquen contra dicha norma,
la escala debe tratarse como una clasificación académica interna y no como un nivel de
integridad en sentido normativo.
\end{notanormativa}

\begin{instruccion}
Complete el cuadro siguiente antes de usar la escala. Un nivel de criticidad sin definición ni
método de asignación no puede justificar decisiones de cobertura ni de técnicas.
\end{instruccion}

\begin{table}[H]
\centering
\caption{Definición local de la escala de criticidad y su efecto sobre la estrategia.}
\label{tab:sil}
\begin{tabularx}{\textwidth}{|C{2.3cm}|Y|C{2.6cm}|L{3.5cm}|}
  \hline
  \filaenc \textbf{Nivel} & \textbf{Significado acordado por el equipo (consecuencia de un fallo)} & \textbf{Quién lo asigna y cuándo} & \textbf{Efecto sobre la estrategia} \\ \hline
  Catastrófico & \campo{p.\,ej. pérdida económica irreversible, daño legal o de seguridad} & \campo{Rol y momento} & \campo{Cobertura objetivo, técnicas y niveles} \\ \hline
  Crítico      & \campo{Significado} & \campo{Rol y momento} & \campo{Efecto} \\ \hline
  Marginal     & \campo{Significado} & \campo{Rol y momento} & \campo{Efecto} \\ \hline
  Despreciable & \campo{Significado} & \campo{Rol y momento} & \campo{Efecto} \\ \hline
\end{tabularx}
\end{table}

El Cuadro~\ref{tab:sil} hace explícito qué significa cada nivel y cómo modifica la estrategia;
la asignación concreta a cada objeto se registra en el Cuadro~\ref{tab:objetos}. El criterio de
asignación debe ser la \emph{consecuencia de un fallo}, no la dificultad de probar el elemento.

% ---------------------------------------------------------------------
\subsection{Alcance de la prueba}
\label{sec:alcance}

\begin{notanormativa}[Corrección respecto de la plantilla original]
La plantilla original orientaba las características a probar únicamente hacia los
\emph{requisitos no funcionales}. El alcance debe declarar explícitamente ambos tipos:
características \textbf{funcionales} (calcular un total, autorizar una operación) y
\textbf{no funcionales} (rendimiento, seguridad, usabilidad, compatibilidad). Además, toda
exclusión debe justificarse y relacionarse con el riesgo que se asume al excluirla.
\end{notanormativa}

\begin{instruccion}
Llene los tres cuadros siguientes. Para las características no funcionales puede apoyarse en el
modelo de calidad de producto de \textbf{ISO/IEC 25010:2023} (adecuación funcional, eficiencia
de desempeño, compatibilidad, capacidad de interacción, fiabilidad, seguridad, mantenibilidad,
flexibilidad y seguridad física) y en los requisitos no funcionales del propio sistema. Si cita
un anexo de \normap{4}, indique la edición consultada.
\end{instruccion}

\begin{table}[H]
\centering
\caption{Características funcionales incluidas en el alcance.}
\label{tab:alcance-f}
\begin{tabularx}{\textwidth}{|C{1.6cm}|C{1.7cm}|Y|C{2.2cm}|C{2.1cm}|}
  \hline
  \filaenc \textbf{Clave} & \textbf{Objeto (\ref{tab:objetos})} & \textbf{Característica funcional a probar} & \textbf{Base (\ref{tab:base-prueba})} & \textbf{Prioridad} \\ \hline
  \campo{CF-01} & \campo{OP-01} & \campo{Cálculo de total con descuentos} & \campo{BP-01} & \campo{Alta} \\ \hline
  & & & & \\ \hline
  & & & & \\ \hline
\end{tabularx}
\end{table}

\begin{table}[H]
\centering
\caption{Características no funcionales incluidas en el alcance.}
\label{tab:alcance-nf}
\begin{tabularx}{\textwidth}{|C{1.6cm}|C{1.7cm}|L{3.5cm}|Y|C{2.1cm}|}
  \hline
  \filaenc \textbf{Clave} & \textbf{Objeto} & \textbf{Característica de calidad} & \textbf{Criterio o umbral acordado} & \textbf{Prioridad} \\ \hline
  \campo{CNF-01} & \campo{OP-02} & \campo{Seguridad — autenticación} & \campo{Umbral verificable} & \campo{Alta} \\ \hline
  & & & & \\ \hline
\end{tabularx}
\notatabla{Una característica no funcional sin criterio o umbral acordado no es verificable: acuérdelo con los interesados antes de diseñar los casos.}
\end{table}

\begin{table}[H]
\centering
\caption{Exclusiones del alcance, su motivo y el riesgo que se asume.}
\label{tab:exclusiones}
\begin{tabularx}{\textwidth}{|C{1.6cm}|L{3.7cm}|Y|C{2.6cm}|}
  \hline
  \filaenc \textbf{Clave} & \textbf{Qué se excluye} & \textbf{Justificación de la exclusión} & \textbf{Riesgo asumido (\ref{sec:riesgos})} \\ \hline
  \campo{EX-01} & \campo{Característica u objeto} & \campo{Por qué no se prueba} & \campo{RP-04} \\ \hline
  & & & \\ \hline
\end{tabularx}
\end{table}

Los Cuadros~\ref{tab:alcance-f} y~\ref{tab:alcance-nf} delimitan lo que sí se probará,
distinguiendo características funcionales de no funcionales; el Cuadro~\ref{tab:exclusiones}
declara lo que queda fuera y, sobre todo, el riesgo que el proyecto acepta al excluirlo. Ese
riesgo debe reaparecer entre los riesgos residuales del apartado~\ref{sec:residuales}.

% ---------------------------------------------------------------------
\subsection{Supuestos y límites}
\label{sec:supuestos}

\begin{instruccion}
Describa los supuestos y limitaciones considerados al elaborar el plan: estándares a cumplir,
política o estrategia de prueba aplicable, tiempo, costo, disponibilidad de personal,
capacitación, herramientas y entornos. Un supuesto es algo que se da por cierto sin haberlo
verificado; si resulta falso, suele convertirse en un riesgo de proyecto.
\end{instruccion}

\begin{table}[H]
\centering
\caption{Supuestos y restricciones del esfuerzo de prueba.}
\label{tab:supuestos}
\begin{tabularx}{\textwidth}{|C{1.6cm}|C{2.2cm}|Y|C{2.6cm}|}
  \hline
  \filaenc \textbf{Clave} & \textbf{Tipo} & \textbf{Enunciado} & \textbf{Si no se cumple} \\ \hline
  \campo{SUP-01} & \campo{Supuesto} & \campo{El entorno estará disponible desde la semana 3} & \campo{RPy-02} \\ \hline
  \campo{LIM-01} & \campo{Restricción} & \campo{No se dispone de datos reales de producción} & \campo{Efecto en el plan} \\ \hline
  & & & \\ \hline
\end{tabularx}
\notatabla{Tipo: \textit{Supuesto} (se da por cierto y puede fallar) · \textit{Restricción} (límite conocido y firme).}
\end{table}

El Cuadro~\ref{tab:supuestos} explicita las condiciones bajo las cuales el plan es válido y
enlaza cada supuesto con el riesgo de proyecto en que se convierte si no se cumple.

\begin{politicaacademica}[Política académica: política general de prueba del curso]
La Experiencia Educativa fija la siguiente política de cobertura mínima \emph{a nivel de
planeación}, en función de la criticidad asignada a cada objeto de prueba
(apartado~\ref{sec:criticidad}):

\medskip
\begin{center}
\begin{tabular}{|L{3.4cm}|C{3.2cm}|}
  \hline
  \filaenc \textbf{Criticidad} & \textbf{Cobertura mínima} \\ \hline
  Catastrófico  & 100\,\% \\ \hline
  Crítico       & 80\,\% \\ \hline
  Marginal      & 60\,\% \\ \hline
  Despreciable  & 40\,\% \\ \hline
\end{tabular}
\end{center}
\medskip

A nivel de ejecución, deberá someterse a prueba la totalidad de los elementos con criticidad
\emph{Catastrófico}. Las demás técnicas abordadas en el curso se aplicarán según corresponda a
ese mismo nivel de integridad, justificando su elección según las características de las
unidades a probar.

\textbf{Estos porcentajes son una regla del curso, no un umbral de \norma.} Para que sean
interpretables debe declararse sobre \emph{qué elemento} se calcula la cobertura: no es lo
mismo cubrir el 80\,\% de los requisitos que el 80\,\% de las particiones de equivalencia o de
las decisiones del código. Defina ese elemento en el apartado~\ref{sec:cobertura}.
\end{politicaacademica}

\begin{instruccion}
Si reutiliza esta plantilla fuera del curso, sustituya la caja anterior por la política de
prueba de su organización, o por umbrales acordados con los interesados y justificados por el
riesgo. Deje constancia de esa adaptación en el apartado~\ref{sec:desviaciones}.
\end{instruccion}

% ---------------------------------------------------------------------
\subsection{Interesados en la prueba}

\begin{instruccion}
Relacione a los interesados (\emph{stakeholders}) en la prueba y su relevancia para el proceso,
indicando los objetivos y expectativas que tienen sobre la prueba en función de su rol.
Identifique también qué información necesita recibir cada uno y con qué periodicidad: ese dato
alimenta directamente el plan de comunicación del apartado~\ref{sec:informes}.
\end{instruccion}

\begin{table}[H]
\centering
\caption{Interesados en la prueba, expectativas y necesidades de información.}
\label{tab:interesados}
\begin{tabularx}{\textwidth}{|L{3.2cm}|L{2.9cm}|Y|C{2.9cm}|}
  \hline
  \filaenc \textbf{Interesado} & \textbf{Rol frente a la prueba} & \textbf{Objetivo o expectativa} & \textbf{Información y frecuencia} \\ \hline
  \campo{Nombre o grupo} & \campo{Decide / informa / provee} & \campo{Qué espera obtener de la prueba} & \campo{Informe de estado semanal} \\ \hline
  & & & \\ \hline
  & & & \\ \hline
\end{tabularx}
\end{table}

El Cuadro~\ref{tab:interesados} identifica a quién sirve el esfuerzo de prueba y qué espera
cada parte; la última columna determina los destinatarios y la periodicidad de los informes de
estado definidos en el apartado~\ref{sec:informes} y formalizados en el \ref{anx:informes}.

% =====================================================================
%  3. Registro de riesgos   (actividades TP3 y TP4 de ISO/IEC/IEEE 29119-2)
% =====================================================================
\section{Registro de riesgos}
\label{sec:riesgos}

\begin{notanormativa}[Correspondencia con el modelo de procesos]
Este apartado desarrolla las actividades \textbf{TP3 «Identificar y analizar riesgos»} y
\textbf{TP4 «Identificar enfoques de mitigación»} de \normap{2}. La planificación es el único
proceso de gestión que integra explícitamente la gestión de riesgos como paso previo al diseño
de la estrategia: los riesgos priorizados aquí son la justificación de las decisiones del
apartado~\ref{sec:estrategia}.
\end{notanormativa}

\begin{notanormativa}[Corrección respecto de la plantilla original]
La plantilla original explicaba el \emph{riesgo del producto} por su efecto sobre el tiempo y
el desarrollo de la prueba. Eso corresponde en realidad a un riesgo de proyecto. La distinción
correcta es:
\begin{itemize}
  \item \textbf{Riesgo del producto}: consecuencia potencial de un fallo o deficiencia del
        \emph{software} sobre los usuarios, el negocio, los datos, la seguridad o la operación.
        \emph{Ejemplo: un cálculo erróneo cobra de más al cliente.}
  \item \textbf{Riesgo del proyecto}: amenaza a la capacidad de \emph{ejecutar el esfuerzo de
        prueba}: tiempo, recursos, personal, entorno, herramientas o acceso a la información.
        \emph{Ejemplo: la indisponibilidad del entorno retrasa la ejecución.}
\end{itemize}
Ambos se evalúan por probabilidad e impacto, pero el «impacto» no significa lo mismo en cada
caso: en el producto se mide sobre el usuario o el negocio; en el proyecto, sobre el plan.
\end{notanormativa}

\begin{instruccion}
Use una escala explícita y aplíquela igual en los dos registros. Si emplea valores numéricos,
diga cómo se combinan (por ejemplo, exposición $=$ probabilidad $\times$ impacto) y qué umbral
convierte un riesgo en prioritario. Los riesgos priorizados deben poder rastrearse hasta las
decisiones de estrategia que los mitigan.
\end{instruccion}

\begin{table}[H]
\centering
\caption{Escala acordada para valorar probabilidad, impacto y exposición.}
\label{tab:escala-riesgo}
\begin{tabularx}{\textwidth}{|C{2.0cm}|Y|Y|}
  \hline
  \filaenc \textbf{Valor} & \textbf{Probabilidad — significado} & \textbf{Impacto — significado} \\ \hline
  \campo{Alta (3)}  & \campo{Criterio observable} & \campo{Criterio observable} \\ \hline
  \campo{Media (2)} & \campo{Criterio observable} & \campo{Criterio observable} \\ \hline
  \campo{Baja (1)}  & \campo{Criterio observable} & \campo{Criterio observable} \\ \hline
\end{tabularx}
\notatabla{Fórmula de exposición: \campo{p.\,ej. exposición = probabilidad × impacto}. Umbral de priorización: \campo{p.\,ej. exposición ≥ 6}.}
\end{table}

El Cuadro~\ref{tab:escala-riesgo} fija el significado de cada valor para que dos personas
distintas puedan valorar el mismo riesgo de forma comparable.

% ---------------------------------------------------------------------
\subsection{Riesgos del producto}

\begin{instruccion}
Identifique qué puede salir mal \textbf{en el software} y qué consecuencia tendría para
usuarios, negocio, datos, seguridad u operación. Estos riesgos son la principal justificación
para decidir qué se prueba primero, con qué profundidad y con qué técnicas.
\end{instruccion}

\begin{table}[H]
\centering
\caption{Riesgos del producto: consecuencias potenciales de un fallo del software.}
\label{tab:riesgos-producto}
\begin{tabularx}{\textwidth}{|C{1.5cm}|C{1.6cm}|Y|C{1.4cm}|C{1.4cm}|C{1.7cm}|}
  \hline
  \filaenc \textbf{Clave} & \textbf{Objeto (\ref{tab:objetos})} & \textbf{Descripción: qué falla y qué consecuencia tiene} & \textbf{Prob.} & \textbf{Impacto} & \textbf{Exposición} \\ \hline
  \campo{RP-01} & \campo{OP-01} & \campo{Si el cálculo de impuestos es incorrecto, se factura de más al cliente y se incurre en incumplimiento fiscal} & \campo{2} & \campo{3} & \campo{6} \\ \hline
  & & & & & \\ \hline
  & & & & & \\ \hline
\end{tabularx}
\end{table}

El Cuadro~\ref{tab:riesgos-producto} enumera las consecuencias que la prueba pretende hacer
visibles antes de la liberación; su columna de exposición ordena la prioridad de diseño y
ejecución.

% ---------------------------------------------------------------------
\subsection{Riesgos del proyecto}

\begin{instruccion}
Identifique lo que puede impedir o degradar la \emph{realización} del esfuerzo de prueba:
planeación, recursos, personal, competencias, entorno, herramientas, dependencias externas o
acceso a la base de prueba. Los supuestos del Cuadro~\ref{tab:supuestos} que puedan fallar
deben aparecer aquí.
\end{instruccion}

\begin{table}[H]
\centering
\caption{Riesgos del proyecto: amenazas a la ejecución del esfuerzo de prueba.}
\label{tab:riesgos-proyecto}
\begin{tabularx}{\textwidth}{|C{1.5cm}|C{2.0cm}|Y|C{1.4cm}|C{1.4cm}|C{1.7cm}|}
  \hline
  \filaenc \textbf{Clave} & \textbf{Categoría} & \textbf{Descripción del riesgo} & \textbf{Prob.} & \textbf{Impacto} & \textbf{Exposición} \\ \hline
  \campo{RPy-01} & \campo{Entorno} & \campo{El entorno de prueba no está disponible en la fecha prevista y la ejecución se retrasa} & \campo{3} & \campo{2} & \campo{6} \\ \hline
  & & & & & \\ \hline
  & & & & & \\ \hline
\end{tabularx}
\notatabla{Categorías sugeridas: planeación · recursos · personal y competencias · entorno · herramientas · datos · dependencias externas · calidad de la base de prueba.}
\end{table}

El Cuadro~\ref{tab:riesgos-proyecto} recoge las amenazas al plan mismo; se vigilan durante el
seguimiento (apartado~\ref{sec:seguimiento}) y pueden disparar acciones de control o
replanificación.

% ---------------------------------------------------------------------
\subsection{Tratamiento de los riesgos}

\begin{instruccion}
Para cada riesgo priorizado, indique cómo se tratará. No todos los riesgos se mitigan mediante
prueba: algunos se aceptan, se transfieren o se evitan cambiando el alcance. Declare
explícitamente cuáles se mitigarán \emph{probando} y cuáles no, porque esa decisión determina
el contenido de la estrategia.

Los riesgos que se acepten sin mitigar deben reaparecer en el apartado~\ref{sec:residuales}
al cierre.
\end{instruccion}

\begin{table}[H]
\centering
\caption{Tratamiento acordado para cada riesgo priorizado.}
\label{tab:tratamiento-riesgos}
\begin{tabularx}{\textwidth}{|C{1.5cm}|C{2.1cm}|Y|C{2.6cm}|C{1.9cm}|}
  \hline
  \filaenc \textbf{Riesgo} & \textbf{Tratamiento} & \textbf{Acción concreta} & \textbf{Se mitiga mediante prueba} & \textbf{Responsable} \\ \hline
  \campo{RP-01} & \campo{Mitigar} & \campo{Aplicar tablas de decisión y valores límite sobre el cálculo} & \campo{Sí — CF-01} & \campo{Rol} \\ \hline
  \campo{RPy-01} & \campo{Mitigar} & \campo{Preparar entorno alterno y verificarlo en la semana 2} & \campo{No} & \campo{Rol} \\ \hline
  & & & & \\ \hline
\end{tabularx}
\notatabla{Tratamiento: \textit{Mitigar} · \textit{Aceptar} · \textit{Transferir} · \textit{Evitar}. Sólo la columna «se mitiga mediante prueba» justifica contenido de la estrategia.}
\end{table}

El Cuadro~\ref{tab:tratamiento-riesgos} cierra el vínculo entre riesgo y estrategia: cada
riesgo que se decide mitigar probando debe corresponderse con al menos una condición de prueba
y una técnica en el apartado~\ref{sec:tecnicas}, y con una fila de la matriz de trazabilidad
del \ref{anx:trazabilidad}.

% =====================================================================
%  4. Estrategia de prueba   (actividad TP5 de ISO/IEC/IEEE 29119-2)
% =====================================================================
\section{Estrategia de prueba}
\label{sec:estrategia}

\begin{notanormativa}[Correspondencia con el modelo de procesos]
Este apartado desarrolla la actividad \textbf{TP5 «Diseñar la estrategia de prueba»} de
\normap{2}: seleccionar niveles, tipos y técnicas coherentes con los riesgos; definir criterios
de entrada, salida, suspensión y reanudación; y establecer el enfoque de repetición y
regresión, los requisitos de entorno y las métricas. Cada decisión de este apartado debería
poder justificarse con un riesgo del apartado~\ref{sec:riesgos}.
\end{notanormativa}

% ---------------------------------------------------------------------
\subsection{Niveles de prueba}
\label{sec:niveles}

\begin{instruccion}
Identifique los niveles (subprocesos) de prueba que se ejecutarán, en qué consisten y en qué
orden se realizarán.

\textbf{Corrección respecto de la plantilla original:} la redacción anterior sugería que los
cuatro niveles habituales debían ejecutarse siempre. No es así: \emph{elija y justifique} los
niveles aplicables a su proyecto. Un nivel que no se ejecuta debe declararse como no aplicable
con su motivo, no omitirse en silencio. Separe con claridad componente, integración, sistema y
aceptación: no son sinónimos ni etapas intercambiables.
\end{instruccion}

\begin{table}[H]
\centering
\caption{Niveles de prueba seleccionados, con su justificación y orden de realización.}
\label{tab:niveles}
\begin{tabularx}{\textwidth}{|L{2.6cm}|C{1.9cm}|Y|C{1.5cm}|C{2.2cm}|}
  \hline
  \filaenc \textbf{Nivel} & \textbf{¿Se aplica?} & \textbf{En qué consiste en este proyecto / por qué se aplica o no} & \textbf{Orden} & \textbf{Objetos (\ref{tab:objetos})} \\ \hline
  Componente   & \campo{Sí/No} & \campo{Justificación} & \campo{1} & \campo{OP-01} \\ \hline
  Integración  & \campo{Sí/No} & \campo{Justificación} & \campo{2} & \campo{OP-01, OP-02} \\ \hline
  Sistema      & \campo{Sí/No} & \campo{Justificación} & \campo{3} & \campo{OP-01} \\ \hline
  Aceptación   & \campo{Sí/No} & \campo{Justificación} & \campo{4} & \campo{OP-01} \\ \hline
  \campo{Otro} & \campo{Sí/No} & \campo{Justificación} & \campo{} & \campo{} \\ \hline
\end{tabularx}
\end{table}

El Cuadro~\ref{tab:niveles} declara qué niveles se ejecutan y cuáles se descartan con su
motivo; el orden indicado condiciona el calendario del apartado~\ref{sec:calendario}.

% ---------------------------------------------------------------------
\subsection{Tipos de prueba}

\begin{instruccion}
Indique los tipos de prueba que se aplicarán (funcional, rendimiento, seguridad, usabilidad,
compatibilidad, entre otros) y con qué característica del alcance no funcional
(Cuadro~\ref{tab:alcance-nf}) o funcional (Cuadro~\ref{tab:alcance-f}) se corresponde cada uno.
Un tipo de prueba sin característica asociada suele indicar que el alcance está incompleto.
\end{instruccion}

\begin{table}[H]
\centering
\caption{Tipos de prueba a aplicar y su correspondencia con el alcance.}
\label{tab:tipos}
\begin{tabularx}{\textwidth}{|L{3.0cm}|C{2.4cm}|Y|C{2.3cm}|}
  \hline
  \filaenc \textbf{Tipo de prueba} & \textbf{Nivel(es)} & \textbf{Qué pretende evidenciar} & \textbf{Característica} \\ \hline
  \campo{Funcional} & \campo{Sistema} & \campo{Objetivo} & \campo{CF-01} \\ \hline
  \campo{Seguridad} & \campo{Sistema} & \campo{Objetivo} & \campo{CNF-01} \\ \hline
  & & & \\ \hline
\end{tabularx}
\end{table}

El Cuadro~\ref{tab:tipos} evita una omisión frecuente: declarar características no funcionales
en el alcance (apartado~\ref{sec:alcance}) y no prever ningún tipo de prueba capaz de
evaluarlas en alguno de los niveles seleccionados en el apartado~\ref{sec:niveles}.

% ---------------------------------------------------------------------
\subsection{Diseño e implementación: condiciones y elementos de cobertura}
\label{sec:diseno}

\begin{notanormativa}[Cadena de derivación de \normap{2}]
Las actividades \textbf{TD1–TD6} definen una cadena trazable que va de la base de prueba al
procedimiento ejecutable: identificar conjuntos de características (TD1) $\rightarrow$ derivar
condiciones de prueba (TD2) $\rightarrow$ derivar elementos de cobertura (TD3) $\rightarrow$
derivar casos de prueba (TD4) $\rightarrow$ ensamblar conjuntos (TD5) $\rightarrow$ derivar
procedimientos (TD6). La revisión señaló que la plantilla original saltaba de los elementos de
prueba directamente a los casos, perdiendo los eslabones intermedios que hacen verificable la
cobertura.
\end{notanormativa}

\begin{instruccion}
Registre aquí las condiciones de prueba derivadas de la base de prueba y de los riesgos, y los
elementos de cobertura que cada técnica genera sobre ellas. Los casos concretos se documentan
en el \ref{anx:casos}; este cuadro es el eslabón que los justifica.
\end{instruccion}

\begin{table}[H]
\centering
\caption{Condiciones de prueba y elementos de cobertura derivados.}
\label{tab:condiciones}
\begin{tabularx}{\textwidth}{|C{1.6cm}|C{1.6cm}|C{1.4cm}|Y|C{2.6cm}|}
  \hline
  \filaenc \textbf{Clave} & \textbf{Base (\ref{tab:base-prueba})} & \textbf{Riesgo} & \textbf{Condición de prueba (enunciado verificable)} & \textbf{Elementos de cobertura} \\ \hline
  \campo{COND-01} & \campo{BP-01} & \campo{RP-01} & \campo{El total se calcula correctamente para cada combinación de descuento e impuesto} & \campo{4 particiones + 6 valores límite} \\ \hline
  & & & & \\ \hline
  & & & & \\ \hline
\end{tabularx}
\end{table}

El Cuadro~\ref{tab:condiciones} materializa las actividades TD2 y TD3 y es el origen de la
columna «condición» de la matriz de trazabilidad del \ref{anx:trazabilidad}.

% ---------------------------------------------------------------------
\subsection{Técnicas de diseño de prueba}
\label{sec:tecnicas}

\begin{instruccion}
Indique y justifique las técnicas específicas que se aplicarán en cada nivel y sobre cada
objeto, según su naturaleza. \textbf{No es obligatorio usar todas las técnicas}: hay que elegir
las pertinentes y explicar su relación con los objetivos y los riesgos. Para cada técnica
seleccionada complete: objetivo, elemento al que aplica, justificación, cobertura que produce y
riesgo que ayuda a reducir.
\end{instruccion}

% =====================================================================
%  diagrams/tecnicas-prueba.tex
%  Figura: mapa de técnicas de diseño de prueba.
%
%  CORRECCIONES respecto del diagrama original («tenicas de prueba.png»):
%   1. El original fijaba «cobertura mínima: 100 %» para sentencias y
%      decisiones, atribuyéndolo implícitamente a la norma. Los
%      porcentajes se han retirado: cada técnica indica ahora el
%      ELEMENTO DE COBERTURA que produce, y el objetivo numérico se
%      declara —con su justificación— en el apartado 4.5 del plan.
%   2. Se añade la advertencia de que el mapa es un menú de selección:
%      no hay obligación de aplicar todas las técnicas.
%   3. Se añade el enfoque por palabras clave, ausente del original, con
%      su carácter opcional explícito.
% =====================================================================
\begin{figure}[htbp]
\centering
\begin{tikzpicture}[
  font=\scriptsize,
  fam/.style  ={draw=uvazul, very thick, fill=uvazul!12, rounded corners=2pt,
                text width=4.4cm, align=center, minimum height=1.0cm, inner sep=3pt},
  tec/.style  ={draw=gray!70, fill=gray!6, rounded corners=2pt,
                text width=4.4cm, align=left, minimum height=0.9cm, inner sep=4pt},
  nota/.style ={draw=polinaranja, fill=polfondo, rounded corners=2pt,
                text width=14.6cm, align=left, inner sep=5pt}
]

% ------------------------- familias (columnas) ------------------------
\node[fam] (f1) at (-5.2,0) {\textbf{Basadas en especificación}\\ \emph{caja negra}};
\node[fam] (f2) at (0,0)    {\textbf{Basadas en estructura}\\ \emph{caja blanca}};
\node[fam] (f3) at (5.2,0)  {\textbf{Basadas en experiencia}};

% ------------------------- columna 1: especificación ------------------
\node[tec] (a1) at (-5.2,-1.35) {\textbf{Particiones de equivalencia}\\ \emph{cubre:} cada clase válida e inválida};
\node[tec] (a2) at (-5.2,-2.65) {\textbf{Análisis de valores límite}\\ \emph{cubre:} bordes de cada partición};
\node[tec] (a3) at (-5.2,-3.95) {\textbf{Tablas de decisión}\\ \emph{cubre:} cada regla de la tabla};
\node[tec] (a4) at (-5.2,-5.25) {\textbf{Transiciones de estado}\\ \emph{cubre:} estados y transiciones};
\node[tec] (a5) at (-5.2,-6.55) {\textbf{Basadas en casos de uso}\\ \emph{cubre:} flujo básico y alternativos};

% ------------------------- columna 2: estructura ----------------------
\node[tec] (b1) at (0,-1.35) {\textbf{Cobertura de sentencias}\\ \emph{cubre:} cada sentencia ejecutable};
\node[tec] (b2) at (0,-2.65) {\textbf{Cobertura de decisiones}\\ \emph{cubre:} cada salida de cada decisión};
\node[tec] (b3) at (0,-3.95) {\textbf{Cobertura de condiciones}\\ \emph{cubre:} cada condición simple};
\node[tec] (b4) at (0,-5.25) {\textbf{Condiciones múltiples}\\ \emph{cubre:} combinaciones booleanas};
\node[tec] (b5) at (0,-6.55) {\textbf{Cobertura de rutas y bucles}\\ \emph{cubre:} rutas seleccionadas; 0, 1, N iteraciones};

% ------------------------- columna 3: experiencia ---------------------
\node[tec] (c1) at (5.2,-1.35) {\textbf{Pruebas exploratorias}\\ \emph{útil cuando} la base de prueba es escasa};
\node[tec] (c2) at (5.2,-2.65) {\textbf{Basadas en listas de verificación}\\ \emph{cubre:} los puntos de la lista};
\node[tec] (c3) at (5.2,-3.95) {\textbf{Pruebas por ataque}\\ \emph{útil para} riesgos de seguridad};
\node[tec] (c4) at (5.2,-5.25) {\textbf{Conjetura de errores}\\ \emph{apoyada en} defectos históricos};
\node[tec] (c5) at (5.2,-6.55) {\textbf{Palabras clave} (\normap{5})\\ \emph{enfoque opcional de implementación}};

% ------------------------------ enlaces -------------------------------
\foreach \a/\b in {f1/a1, f2/b1, f3/c1} { \draw[-{Latex[length=4pt]}, uvazul!80] (\a) -- (\b); }
\foreach \a/\b in {a1/a2, a2/a3, a3/a4, a4/a5,
                   b1/b2, b2/b3, b3/b4, b4/b5,
                   c1/c2, c2/c3, c3/c4, c4/c5} { \draw[gray!55] (\a) -- (\b); }

% ------------------------------- nota ---------------------------------
\node[nota] at (0,-8.15) {%
  \textbf{Este mapa es un menú de selección, no una lista de obligaciones.} Se eligen las
  técnicas que el riesgo justifique (Cuadro~\ref{tab:tecnicas}). Los \emph{elementos de
  cobertura} indicados son lo que cada técnica genera; el \emph{objetivo numérico} de cobertura
  no procede de la norma y debe declararse, con su justificación, en el
  apartado~\ref{sec:cobertura}. Alcanzar una cobertura determinada no demuestra la ausencia de
  defectos.};

\end{tikzpicture}
\caption[Mapa de técnicas de diseño de prueba y su elemento de cobertura]%
{Mapa de técnicas de diseño de prueba, agrupadas en tres familias, con el elemento de cobertura
que cada una produce. Versión corregida del diagrama de referencia: se han retirado los
porcentajes de cobertura mínima que se atribuían a la norma.}
\label{fig:tecnicas}
\end{figure}


La Figura~\ref{fig:tecnicas} clasifica las técnicas de diseño más habituales en tres familias
—basadas en especificación, basadas en estructura y basadas en experiencia— e indica qué tipo
de elemento de cobertura genera cada una. Sirve como menú de selección: se elige de él lo que
el riesgo justifique, no todo su contenido.

\begin{notanormativa}[Corrección respecto del diagrama original de técnicas]
El diagrama de técnicas del que parte esta plantilla presentaba «100\,\% de cobertura mínima»
para sentencias y decisiones. Ese valor \textbf{no procede de \norma}: es, a lo sumo, un
objetivo que un proyecto puede fijarse. En la Figura~\ref{fig:tecnicas} los porcentajes se han
retirado y sustituido por el \emph{elemento de cobertura} que cada técnica produce; el objetivo
numérico se declara, con su justificación, en el apartado~\ref{sec:cobertura}. Alcanzar una
cobertura determinada tampoco demuestra la ausencia de defectos.
\end{notanormativa}

\begin{table}[H]
\centering
\caption{Técnicas de diseño seleccionadas y su justificación.}
\label{tab:tecnicas}
\begin{tabularx}{\textwidth}{|L{2.9cm}|C{1.6cm}|C{1.6cm}|Y|C{1.5cm}|}
  \hline
  \filaenc \textbf{Técnica} & \textbf{Familia} & \textbf{Nivel / objeto} & \textbf{Justificación y elemento de cobertura que produce} & \textbf{Riesgo} \\ \hline
  \campo{Valores límite} & \campo{Especificación} & \campo{Componente / OP-01} & \campo{Por qué esta técnica y qué cubre} & \campo{RP-01} \\ \hline
  & & & & \\ \hline
  & & & & \\ \hline
\end{tabularx}
\notatabla{Familias: \textit{basada en especificación} (caja negra) · \textit{basada en estructura} (caja blanca) · \textit{basada en experiencia}.}
\end{table}

El Cuadro~\ref{tab:tecnicas} justifica cada técnica frente a un riesgo concreto; una técnica
sin riesgo asociado suele indicar esfuerzo no priorizado.

\begin{politicaacademica}
El curso exige someter a prueba la totalidad de los elementos con criticidad
\emph{Catastrófico} empleando pruebas basadas en casos de uso. Las demás técnicas se
seleccionan justificando su elección. Registre esa obligación en la columna de justificación
del Cuadro~\ref{tab:tecnicas} indicando que el motivo es la política del curso, para no
confundirla con una decisión basada en riesgo.
\end{politicaacademica}

\begin{instruccion}
Si emplea pruebas por palabras clave (\emph{keyword-driven testing}, \normap{5}) u otro enfoque
no listado en la Figura~\ref{fig:tecnicas}, agréguelo al Cuadro~\ref{tab:tecnicas} con la misma
justificación. Su ausencia del diagrama no lo excluye.
\end{instruccion}

% ---------------------------------------------------------------------
\subsection{Definición de la cobertura}
\label{sec:cobertura}

\begin{notanormativa}[Corrección respecto de la plantilla original]
Un porcentaje de cobertura es ininteligible si no se declara sobre qué se mide. La revisión
señaló además un error frecuente: \textbf{la proporción de casos ejecutados es avance de
ejecución, no cobertura}. Ejecutar el 90\,\% de los casos no significa haber cubierto el 90\,\%
de los requisitos, de las particiones ni de las decisiones del código.
\end{notanormativa}

\begin{instruccion}
Para cada medida de cobertura que vaya a reportar, declare las siete columnas del cuadro
siguiente. Si no puede llenar el numerador y el denominador, la medida no es utilizable como
criterio de salida.
\end{instruccion}

\begin{table}[H]
\centering
\caption{Definición operativa de cada medida de cobertura utilizada.}
\label{tab:cobertura}
\begin{tabularx}{\textwidth}{|C{1.5cm}|L{2.4cm}|Y|C{1.5cm}|C{2.0cm}|}
  \hline
  \filaenc \textbf{Clave} & \textbf{Elemento medido} & \textbf{Numerador / denominador y fórmula} & \textbf{Objetivo} & \textbf{Fuente del dato} \\ \hline
  \campo{COB-01} & \campo{Requisitos} & \campo{Requisitos con al menos un caso aprobado ÷ requisitos en el alcance} & \campo{100 \%} & \campo{Anexo G} \\ \hline
  \campo{COB-02} & \campo{Particiones de equivalencia} & \campo{Particiones ejercitadas ÷ particiones identificadas} & \campo{80 \%} & \campo{Anexo A} \\ \hline
  & & & & \\ \hline
\end{tabularx}
\notatabla{Elementos posibles: requisitos · escenarios · particiones de equivalencia · valores límite · sentencias · decisiones · ramas · estados y transiciones · reglas de una tabla de decisión.}
\end{table}

\begin{instruccion}
Añada bajo el cuadro la \textbf{interpretación} de cada medida: qué conclusión permite extraer
y cuál no. Por ejemplo: «COB-01 al 100\,\% indica que ningún requisito del alcance quedó sin
ejercitar; no indica que el software esté libre de defectos ni que los casos hayan sido
adecuados.»
\end{instruccion}

\campo{Interpretación de cada medida de cobertura}

El Cuadro~\ref{tab:cobertura} convierte los porcentajes de la política académica
(apartado~\ref{sec:criticidad}) en medidas calculables y auditables. Sus objetivos se retoman
como criterios de salida en el apartado~\ref{sec:criterios-salida}.

% ---------------------------------------------------------------------
\subsection{Criterios de entrada}
\label{sec:criterios-entrada}

\begin{notanormativa}[Apartado nuevo respecto de la plantilla original]
La plantilla original carecía de criterios de entrada. Sin ellos no puede decidirse cuándo es
legítimo empezar a ejecutar, y las ejecuciones fallidas por causas evitables (versión
equivocada, entorno no listo, datos ausentes) se confunden con defectos del producto. La
actividad \textbf{TE1} de \normap{2} incluye explícitamente «verificar el cumplimiento de los
criterios de entrada» antes de ejecutar.
\end{notanormativa}

\begin{notanormativa}[Criterio de entrada $\neq$ precondición de un caso de prueba]
Un \textbf{criterio de entrada} es una condición de \emph{proceso}: habilita el inicio de un
ciclo, nivel o campaña de ejecución y se verifica una vez, por un responsable identificado. Una
\textbf{precondición} es una condición de \emph{estado del sistema} necesaria para que un caso
concreto sea ejecutable, y se establece en cada ejecución de ese caso (véase el
\ref{anx:casos}). Confundirlas hace imposible auditar por qué se inició una campaña.
\end{notanormativa}

\begin{instruccion}
Defina los criterios que deben cumplirse antes de iniciar cada ciclo o nivel de prueba, y quién
los verifica. El registro de la verificación efectiva se lleva en el \ref{anx:entorno}.
\end{instruccion}

\begin{table}[H]
\centering
\caption{Criterios de entrada para iniciar un ciclo o nivel de prueba.}
\label{tab:criterios-entrada}
\begin{tabularx}{\textwidth}{|C{1.5cm}|Y|C{2.6cm}|C{2.3cm}|}
  \hline
  \filaenc \textbf{Clave} & \textbf{Criterio: qué debe estar listo} & \textbf{Cómo se verifica (evidencia)} & \textbf{Quién verifica} \\ \hline
  CE-01 & La versión del objeto de prueba está identificada y desplegada en el entorno. & \campo{Evidencia} & \campo{Rol} \\ \hline
  CE-02 & Los casos de prueba del ciclo están especificados y revisados. & \campo{Evidencia} & \campo{Rol} \\ \hline
  CE-03 & Los procedimientos de prueba están disponibles y ordenados. & \campo{Evidencia} & \campo{Rol} \\ \hline
  CE-04 & El entorno de prueba está preparado y verificado. & \campo{Evidencia} & \campo{Rol} \\ \hline
  CE-05 & Los datos de prueba están disponibles, cargados y en su versión correcta. & \campo{Evidencia} & \campo{Rol} \\ \hline
  CE-06 & Las personas responsables de ejecutar están asignadas y disponibles. & \campo{Evidencia} & \campo{Rol} \\ \hline
  \campo{CE-07} & \campo{Criterio adicional propio del proyecto} & \campo{Evidencia} & \campo{Rol} \\ \hline
\end{tabularx}
\notatabla{Estos seis criterios son el mínimo recomendado por la revisión; amplíelos según el proyecto. Si un ciclo se inicia sin cumplir alguno, regístrelo como desviación (\ref{sec:desviaciones}).}
\end{table}

El Cuadro~\ref{tab:criterios-entrada} fija las condiciones que habilitan el inicio de la
ejecución y nombra a quién las comprueba; su verificación efectiva, ciclo por ciclo, se
documenta en el \ref{anx:entorno}.

% ---------------------------------------------------------------------
\subsection{Criterios de finalización (criterios de salida)}
\label{sec:criterios-salida}

\begin{notanormativa}[Corrección respecto de la plantilla original]
La plantilla original resolvía la finalización por cobertura alcanzada, argumentando que en el
curso no se reparan los defectos encontrados. Eso mezcla tres cosas distintas que la revisión
pide separar: (i)~haber \emph{completado las actividades} planificadas, (ii)~los
\emph{resultados obtenidos}, y (iii)~la \emph{decisión de aceptación} del producto. Este
apartado define únicamente cuándo se dan por terminadas las actividades de prueba. La decisión
de aceptación se trata en el apartado~\ref{sec:aceptacion}.
\end{notanormativa}

\begin{instruccion}
Describa las condiciones bajo las cuales se considerarán completadas las actividades de prueba.
Combine al menos un criterio de cobertura (del Cuadro~\ref{tab:cobertura}) con criterios sobre
los resultados y sobre los incidentes abiertos. Indique qué ocurre si un criterio no se alcanza:
suspender, ampliar el esfuerzo o cerrar documentando la desviación.
\end{instruccion}

\begin{table}[H]
\centering
\caption{Criterios de finalización de las actividades de prueba.}
\label{tab:criterios-salida}
\begin{tabularx}{\textwidth}{|C{1.5cm}|C{2.4cm}|Y|C{3.0cm}|}
  \hline
  \filaenc \textbf{Clave} & \textbf{Tipo de criterio} & \textbf{Enunciado y valor objetivo} & \textbf{Si no se alcanza} \\ \hline
  \campo{CS-01} & \campo{Cobertura} & \campo{COB-01 ≥ 100 \% en objetos Catastróficos} & \campo{Ampliar / documentar desviación} \\ \hline
  \campo{CS-02} & \campo{Ejecución} & \campo{Todos los procedimientos planificados ejecutados o formalmente cancelados} & \campo{Acción} \\ \hline
  \campo{CS-03} & \campo{Incidentes} & \campo{Ningún incidente de severidad crítica sin analizar} & \campo{Acción} \\ \hline
  & & & \\ \hline
\end{tabularx}
\end{table}

El Cuadro~\ref{tab:criterios-salida} distingue criterios de cobertura, de ejecución y de
incidentes, de modo que el cierre no dependa de una sola cifra. Su evaluación efectiva se
reporta en el apartado~\ref{sec:evaluacion-salida}.

% ---------------------------------------------------------------------
\subsection{Criterios de suspensión y reanudación}

\begin{instruccion}
Especifique los criterios para suspender o reanudar las actividades de prueba, quién autoriza
cada decisión y qué actividades deberán repetirse tras una suspensión. Un caso típico de
suspensión es un entorno inestable que invalida los veredictos obtenidos.
\end{instruccion}

\begin{table}[H]
\centering
\caption{Criterios de suspensión y reanudación de las actividades de prueba.}
\label{tab:suspension}
\begin{tabularx}{\textwidth}{|C{2.3cm}|Y|C{2.4cm}|C{3.0cm}|}
  \hline
  \filaenc \textbf{Situación} & \textbf{Criterio} & \textbf{Autoriza} & \textbf{Qué se repite al reanudar} \\ \hline
  Suspensión & \campo{p.\,ej. el entorno deja de ser representativo} & \campo{Rol} & \campo{Alcance de la repetición} \\ \hline
  Reanudación & \campo{Condición que debe restablecerse} & \campo{Rol} & \campo{Alcance de la repetición} \\ \hline
\end{tabularx}
\end{table}

El Cuadro~\ref{tab:suspension} evita una discusión frecuente al reanudar: qué resultados
obtenidos antes de la suspensión siguen siendo válidos y cuáles deben volver a obtenerse.

% ---------------------------------------------------------------------
\subsection{Prueba de confirmación y prueba de regresión}
\label{sec:confirmacion-regresion}

\begin{notanormativa}[Corrección: no son lo mismo]
\begin{itemize}
  \item \textbf{Prueba de confirmación} (\emph{retesting}): se vuelve a ejecutar el caso que
        reveló un defecto, sobre la versión corregida, para comprobar que \emph{ese defecto
        concreto} quedó resuelto.
  \item \textbf{Prueba de regresión}: se ejecuta un conjunto de pruebas sobre funcionalidad que
        \emph{ya funcionaba}, para detectar efectos adversos introducidos por el cambio.
\end{itemize}
La revisión detectó que el diagrama de procesos original rotulaba la actividad «verificar
correcciones» como prueba de regresión. Es prueba de confirmación. Llamar «regresión» a
cualquier verificación de una corrección oculta que la segunda actividad —comprobar que no se
rompió nada más— puede no estarse realizando. La Figura~\ref{fig:procesos} representa ambas
por separado.
\end{notanormativa}

\begin{instruccion}
Especifique, por separado, cuándo se ejecuta cada una, qué conjunto de pruebas se selecciona y
con qué criterio. Para la regresión, indique el criterio de selección: todo el conjunto,
selección por riesgo, por área afectada por el cambio, o priorización.
\end{instruccion}

\begin{table}[H]
\centering
\caption{Enfoque de prueba de confirmación y de prueba de regresión.}
\label{tab:confirmacion-regresion}
\begin{tabularx}{\textwidth}{|L{3.0cm}|Y|Y|}
  \hline
  \filaenc  & \textbf{Prueba de confirmación} & \textbf{Prueba de regresión} \\ \hline
  Qué verifica & Que el defecto reportado quedó corregido. & Que el cambio no introdujo efectos adversos en lo que ya funcionaba. \\ \hline
  Cuándo se ejecuta & \campo{Disparador} & \campo{Disparador} \\ \hline
  Qué se selecciona & \campo{Criterio de selección} & \campo{Criterio de selección} \\ \hline
  Responsable & \campo{Rol} & \campo{Rol} \\ \hline
\end{tabularx}
\end{table}

El Cuadro~\ref{tab:confirmacion-regresion} obliga a decidir las dos actividades por separado;
cada nueva ejecución resultante se registra como una ejecución independiente en el
\ref{anx:ejecuciones}, sin sobrescribir la anterior.

\begin{politicaacademica}
El curso establece que el equipo de prueba \emph{no repara} los defectos encontrados. Esto no
elimina este apartado, pero condiciona su aplicación: declare si el equipo desarrollador
entregará correcciones durante el periodo y, en caso afirmativo, cómo se confirmarán. Si no
habrá correcciones, regístrelo como limitación y deje constancia de los defectos pendientes al
cierre (apartado~\ref{sec:incidentes-abiertos}).
\end{politicaacademica}

\campo{Declare aquí si habrá correcciones durante el periodo y cómo se tratarán}

% ---------------------------------------------------------------------
\subsection{Requisitos de entorno de prueba}
\label{sec:entorno}

\begin{instruccion}
Especifique los requisitos del entorno: hardware, software, redes, herramientas de prueba,
interfaces simuladas y fichas técnicas de los equipos necesarios. Incluya el procedimiento de
instalación y configuración.

\textbf{Ampliación respecto de la plantilla original:} además de \emph{describir} el entorno,
prevea cómo se \emph{comprueba} que está listo, quién lo mantiene, cómo se controlan sus
versiones y cómo se restaura a un estado conocido entre ejecuciones. Estas tareas corresponden
a las actividades \textbf{ES1 «Establecer el entorno»} y \textbf{ES2 «Mantener el entorno»} de
\normap{2}. El registro de cada preparación concreta se lleva en el \ref{anx:entorno}.
\end{instruccion}

\begin{table}[H]
\centering
\caption{Requisitos del entorno de prueba.}
\label{tab:req-entorno}
\begin{tabularx}{\textwidth}{|C{1.5cm}|C{2.2cm}|Y|C{2.2cm}|C{2.1cm}|}
  \hline
  \filaenc \textbf{Clave} & \textbf{Componente} & \textbf{Especificación técnica requerida} & \textbf{Versión objetivo} & \textbf{Responsable} \\ \hline
  \campo{ENT-01} & \campo{Servidor de aplicación} & \campo{Especificación} & \campo{v0.0} & \campo{Rol} \\ \hline
  \campo{ENT-02} & \campo{Base de datos} & \campo{Especificación} & \campo{v0.0} & \campo{Rol} \\ \hline
  & & & & \\ \hline
\end{tabularx}
\end{table}

\begin{table}[H]
\centering
\caption{Gestión del entorno: preparación, verificación, mantenimiento y restauración.}
\label{tab:gestion-entorno}
\begin{tabularx}{\textwidth}{|L{3.4cm}|Y|C{2.6cm}|}
  \hline
  \filaenc \textbf{Aspecto} & \textbf{Enfoque acordado} & \textbf{Responsable} \\ \hline
  Preparación (ES1) & \campo{Procedimiento de instalación y configuración} & \campo{Rol} \\ \hline
  Verificación de disponibilidad & \campo{Cómo se comprueba que el entorno está listo (criterio CE-04)} & \campo{Rol} \\ \hline
  Control de versiones del entorno & \campo{Cómo se identifica la configuración usada en cada ejecución} & \campo{Rol} \\ \hline
  Mantenimiento (ES2) & \campo{Gestión de cambios durante el periodo} & \campo{Rol} \\ \hline
  Restauración entre ejecuciones & \campo{Cuándo y cómo se devuelve el entorno a un estado conocido} & \campo{Rol} \\ \hline
  Liberación al cierre & \campo{Qué se desmonta y qué se conserva (\ref{sec:archivo})} & \campo{Rol} \\ \hline
\end{tabularx}
\end{table}

El Cuadro~\ref{tab:req-entorno} describe \emph{qué} entorno se necesita y el
Cuadro~\ref{tab:gestion-entorno}, \emph{cómo se gobierna} a lo largo del esfuerzo. Sin el
segundo, los veredictos registrados no son reproducibles.

% ---------------------------------------------------------------------
\subsection{Requisitos de datos de prueba}
\label{sec:datos}

\begin{instruccion}
Especifique los datos fuente relevantes para la prueba: qué conjuntos se necesitan, de dónde
provienen, dónde se almacenan y cómo se cargan. Indique si deben anonimizarse o sustituirse por
razones de confidencialidad o protección de datos personales, y quién es responsable de ello.

\textbf{Ampliación respecto de la plantilla original:} los datos también tienen versión. Debe
poder saberse qué conjunto de datos, en qué versión, se usó en cada ejecución; de lo contrario
un resultado no puede reproducirse ni discutirse.
\end{instruccion}

\begin{table}[H]
\centering
\caption{Conjuntos de datos de prueba requeridos.}
\label{tab:datos}
\begin{tabularx}{\textwidth}{|C{1.5cm}|L{2.9cm}|Y|C{1.8cm}|C{2.4cm}|}
  \hline
  \filaenc \textbf{Clave} & \textbf{Conjunto de datos} & \textbf{Contenido, origen y forma de obtención} & \textbf{Versión} & \textbf{Tratamiento de confidencialidad} \\ \hline
  \campo{DAT-01} & \campo{Catálogo de clientes} & \campo{Origen y volumen} & \campo{v1} & \campo{Anonimizado / sintético} \\ \hline
  & & & & \\ \hline
\end{tabularx}
\notatabla{Si algún conjunto contiene datos personales reales, indique la base legal de su uso y el procedimiento de eliminación segura al cierre (\ref{sec:archivo}).}
\end{table}

El Cuadro~\ref{tab:datos} identifica y versiona los datos; su preparación efectiva, ciclo por
ciclo, se registra junto con la del entorno en el \ref{anx:entorno}.

% ---------------------------------------------------------------------
\subsection{Gestión de configuración de los activos de prueba}
\label{sec:configuracion}

\begin{notanormativa}[Ampliación respecto de la plantilla original]
La plantilla original limitaba el control de versiones al propio documento. Para que un
resultado sea reproducible debe poder identificarse la versión de \emph{todos} los activos que
intervinieron en él. La actividad \textbf{TP9} de \normap{2} exige colocar el plan bajo gestión
de configuración, y \textbf{TE3} exige registrar versión y configuración en cada ejecución.
\end{notanormativa}

\begin{instruccion}
Indique dónde se almacena cada tipo de activo, cómo se identifica su versión y quién lo
mantiene. No es necesario un repositorio distinto por tipo: basta con que el plan diga dónde
está cada cosa y cómo se nombra.
\end{instruccion}

\begin{table}[H]
\centering
\caption{Control de configuración de los activos de prueba.}
\label{tab:configuracion}
\begin{tabularx}{\textwidth}{|L{3.2cm}|Y|C{3.0cm}|C{2.1cm}|}
  \hline
  \filaenc \textbf{Activo} & \textbf{Ubicación / repositorio} & \textbf{Esquema de identificación de versión} & \textbf{Responsable} \\ \hline
  Software bajo prueba & \campo{Ubicación} & \campo{p.\,ej. etiqueta de versión o commit} & \campo{Rol} \\ \hline
  Plan de prueba & \campo{Ubicación} & \campo{Ver. del documento} & \campo{Rol} \\ \hline
  Casos de prueba & \campo{Ubicación} & \campo{Esquema} & \campo{Rol} \\ \hline
  Procedimientos & \campo{Ubicación} & \campo{Esquema} & \campo{Rol} \\ \hline
  Guiones automatizados & \campo{Ubicación} & \campo{Esquema} & \campo{Rol} \\ \hline
  Datos de prueba & \campo{Ubicación} & \campo{Esquema} & \campo{Rol} \\ \hline
  Configuración del entorno & \campo{Ubicación} & \campo{Esquema} & \campo{Rol} \\ \hline
  Evidencias y registros & \campo{Ubicación} & \campo{Esquema} & \campo{Rol} \\ \hline
\end{tabularx}
\end{table}

El Cuadro~\ref{tab:configuracion} es la referencia que hace utilizables las columnas «versión»
y «configuración» del registro de ejecución (\ref{anx:ejecuciones}) y sostiene el archivo de
activos al cierre (apartado~\ref{sec:archivo}).

% ---------------------------------------------------------------------
\subsection{Métricas a recopilar}
\label{sec:metricas}

\begin{instruccion}
Describa las métricas que se calcularán a partir de la información recopilada durante la
prueba. Para cada una indique fórmula, fuente del dato, periodicidad y —sobre todo— cómo se
interpretará el resultado. Una métrica cuya interpretación no se ha acordado de antemano
suele terminar usándose para justificar una decisión ya tomada.

Distinga las métricas de las medidas de cobertura del apartado~\ref{sec:cobertura}: la
cobertura es un tipo particular de métrica y ya está definida allí.
\end{instruccion}

\begin{table}[H]
\centering
\caption{Métricas de prueba a recopilar, con su fórmula e interpretación.}
\label{tab:metricas}
\begin{tabularx}{\textwidth}{|C{1.5cm}|C{1.7cm}|L{3.1cm}|Y|C{2.0cm}|}
  \hline
  \filaenc \textbf{Clave} & \textbf{Tipo} & \textbf{Fórmula de cálculo} & \textbf{Interpretación acordada: qué indica y qué no} & \textbf{Fuente} \\ \hline
  \campo{MET-01} & \campo{Producto} & \campo{Fórmula} & \campo{Interpretación} & \campo{Anexo D} \\ \hline
  \campo{MET-02} & \campo{Proceso} & \campo{Fórmula} & \campo{Interpretación} & \campo{Anexo C} \\ \hline
  \campo{MET-03} & \campo{Proyecto} & \campo{Fórmula} & \campo{Interpretación} & \campo{Anexo E} \\ \hline
  & & & & \\ \hline
\end{tabularx}
\end{table}

El Cuadro~\ref{tab:metricas} define las métricas antes de recolectarlas; sus valores finales se
presentan y discuten en el apartado~\ref{sec:metricas-finales}.

\begin{politicaacademica}
Como requisito de acreditación, deben incluirse al menos \textbf{dos métricas de cada tipo}:
\begin{itemize}
  \item \textbf{De proceso}: permiten identificar mejoras en la eficiencia del proceso de
        desarrollo o de prueba.
  \item \textbf{De producto}: permiten medir la calidad del software.
  \item \textbf{De proyecto}: permiten medir la eficiencia del equipo o de las herramientas.
\end{itemize}
El número mínimo es una regla del curso. Fuera de él, la cantidad de métricas debe responder a
las preguntas que los interesados necesitan responder (Cuadro~\ref{tab:interesados}), no a una
cuota.
\end{politicaacademica}

% ---------------------------------------------------------------------
\subsection{Entregables de la prueba}
\label{sec:entregables}

\begin{instruccion}
Identifique los documentos y registros que se generarán. Indique para cada uno su formato, dónde
se mantendrá y quién es responsable. Los anexos de esta plantilla ya proporcionan formatos
listos para usar; si emplea una herramienta externa (gestor de incidencias, hoja de cálculo,
repositorio), basta con indicar dónde está y quién la mantiene, tal como admite la revisión.
\end{instruccion}

\begin{table}[H]
\centering
\caption{Entregables del esfuerzo de prueba.}
\label{tab:entregables}
\begin{tabularx}{\textwidth}{|L{4.4cm}|C{2.4cm}|Y|C{1.9cm}|}
  \hline
  \filaenc \textbf{Entregable} & \textbf{Formato} & \textbf{Ubicación / herramienta} & \textbf{Responsable} \\ \hline
  Plan de prueba (Parte I) & Este documento & \campo{Ubicación} & \campo{Rol} \\ \hline
  Reporte de finalización (Parte II) & Este documento & \campo{Ubicación} & \campo{Rol} \\ \hline
  Especificación de casos de prueba & \ref{anx:casos} & \campo{Ubicación} & \campo{Rol} \\ \hline
  Especificación de procedimientos & \ref{anx:procedimientos} & \campo{Ubicación} & \campo{Rol} \\ \hline
  Registro de ejecución & \ref{anx:ejecuciones} & \campo{Ubicación} & \campo{Rol} \\ \hline
  Registro de incidentes & \ref{anx:incidentes} & \campo{Ubicación} & \campo{Rol} \\ \hline
  Informes de estado & \ref{anx:informes} & \campo{Ubicación} & \campo{Rol} \\ \hline
  Registro de preparación de entorno y datos & \ref{anx:entorno} & \campo{Ubicación} & \campo{Rol} \\ \hline
  Matriz de trazabilidad & \ref{anx:trazabilidad} & \campo{Ubicación} & \campo{Rol} \\ \hline
\end{tabularx}
\end{table}

El Cuadro~\ref{tab:entregables} reúne los productos documentales del esfuerzo de prueba; los
seis últimos no existían en la plantilla original y son los que permiten reconstruir qué se
probó, en qué condiciones y con qué resultado.

% ---------------------------------------------------------------------
\subsection{Desviaciones respecto de la estrategia acordada}
\label{sec:desviaciones}

\begin{instruccion}
Indique cualquier aspecto de este plan que se aparte de la política o estrategia de prueba
aplicable (apartado~\ref{sec:supuestos}), con su justificación y la persona que la autoriza.
Registre aquí también las adaptaciones realizadas al reutilizar esta plantilla fuera del curso:
sustitución de umbrales académicos, cambio de la autoridad de aprobación o exclusión de algún
proceso.
\end{instruccion}

\begin{table}[H]
\centering
\caption{Desviaciones respecto de la estrategia acordada.}
\label{tab:desviaciones}
\begin{tabularx}{\textwidth}{|C{1.5cm}|L{3.0cm}|Y|C{2.4cm}|C{2.5cm}|}
  \hline
  \filaenc \textbf{Clave} & \textbf{Aspecto que se desvía} & \textbf{Justificación} & \textbf{Autoriza} & \textbf{Fecha} \\ \hline
  \campo{DES-01} & \campo{Aspecto} & \campo{Motivo} & \autoridadAprobacion & \campo{dd/mm/aaaa} \\ \hline
  & & & & \\ \hline
\end{tabularx}
\end{table}

El Cuadro~\ref{tab:desviaciones} deja constancia de lo que se hizo distinto y con qué
autorización; sin él, una adaptación legítima resulta indistinguible de un incumplimiento.

% =====================================================================
%  5. Seguimiento y control   (proceso TMC de ISO/IEC/IEEE 29119-2)
% =====================================================================
\section{Seguimiento y control de la prueba}
\label{sec:seguimiento}

\begin{notanormativa}[Ampliación respecto de la plantilla original]
La plantilla original resolvía este apartado con un listado de «puntos de revisión»: las
entregas del curso y sus fechas. Eso cubre los hitos, pero no el \emph{seguimiento}: no dice
qué se mide, con qué frecuencia, quién compara lo planificado con lo observado, ni qué se hace
cuando aparece una desviación.

Este apartado desarrolla el proceso de \textbf{monitoreo y control (TMC1–TMC4)} de \normap{2}:
preparación (establecer indicadores y línea base), monitorear (recolectar y comparar), controlar
(emitir directivas y, si procede, replanificar) y reportar. Como muestra la
Figura~\ref{fig:procesos}, este proceso opera de forma \emph{continua y concurrente} con la
ejecución: no es una etapa que ocurra entre dos entregas.
\end{notanormativa}

% ---------------------------------------------------------------------
\subsection{Enfoque de monitoreo (TMC1)}

\begin{instruccion}
Defina qué se va a observar, con qué frecuencia, quién lo hace y contra qué línea base se
compara. La línea base es el plan aprobado: sin ella no puede hablarse de desviación.
\end{instruccion}

\begin{table}[H]
\centering
\caption{Enfoque de monitoreo: indicadores, frecuencia y responsables.}
\label{tab:enfoque-monitoreo}
\begin{tabularx}{\textwidth}{|L{3.6cm}|Y|C{2.4cm}|C{2.1cm}|}
  \hline
  \filaenc \textbf{Qué se monitorea} & \textbf{Indicador y fuente del dato} & \textbf{Frecuencia} & \textbf{Responsable} \\ \hline
  Avance de la ejecución & \campo{Procedimientos ejecutados vs. planificados — \ref{anx:ejecuciones}} & \campo{Semanal} & \campo{Rol} \\ \hline
  Cobertura alcanzada & \campo{Medidas COB-xx del Cuadro~\ref{tab:cobertura}} & \campo{Semanal} & \campo{Rol} \\ \hline
  Incidentes & \campo{Altas, estado y severidad — \ref{anx:incidentes}} & \campo{Continua} & \campo{Rol} \\ \hline
  Riesgos & \campo{Riesgos materializados y riesgos nuevos} & \campo{Semanal} & \campo{Rol} \\ \hline
  Consumo de esfuerzo & \campo{Horas empleadas vs. estimadas — \ref{sec:actividades}} & \campo{Semanal} & \campo{Rol} \\ \hline
  Bloqueos & \campo{Pruebas bloqueadas y su causa} & \campo{Continua} & \campo{Rol} \\ \hline
\end{tabularx}
\notatabla{Línea base para la comparación: \campo{versión del plan aprobada y fecha}.}
\end{table}

El Cuadro~\ref{tab:enfoque-monitoreo} materializa la actividad TMC1 y determina qué columnas
debe traer cada informe de estado del \ref{anx:informes}.

% ---------------------------------------------------------------------
\subsection{Puntos de revisión e hitos}

\begin{instruccion}
Describa los informes planificados y los hitos en que se presentarán resultados parciales o
finales, indicando qué se entrega y en qué fecha.
\end{instruccion}

\begin{table}[H]
\centering
\caption{Puntos de revisión e hitos de entrega.}
\label{tab:hitos}
\begin{tabularx}{\textwidth}{|C{1.5cm}|C{2.5cm}|Y|C{2.2cm}|C{2.0cm}|}
  \hline
  \filaenc \textbf{Hito} & \textbf{Fecha} & \textbf{Qué se entrega} & \textbf{Destinatario} & \textbf{Estatus} \\ \hline
  \campo{H-01} & \campo{dd/mm/aaaa} & \campo{Contenido de la entrega} & \campo{Interesado} & \campo{Pendiente} \\ \hline
  \campo{H-02} & \campo{dd/mm/aaaa} & \campo{Contenido de la entrega} & \campo{Interesado} & \campo{Pendiente} \\ \hline
  \campo{H-03} & \campo{dd/mm/aaaa} & \campo{Contenido de la entrega} & \campo{Interesado} & \campo{Pendiente} \\ \hline
\end{tabularx}
\end{table}

El Cuadro~\ref{tab:hitos} recoge los momentos de entrega formal; se refleja en el cronograma del
apartado~\ref{sec:calendario}.

\begin{politicaacademica}
El plan de clase de la Experiencia Educativa establece \textbf{tres entregas} a lo largo del
periodo. Consígnelas como hitos en el Cuadro~\ref{tab:hitos} junto con lo que se presentará en
cada una y sus fechas.

El número de entregas es una regla del curso. Al reutilizar la plantilla, sustitúyalo por los
hitos que el proyecto requiera: los puntos de revisión deberían fijarse por necesidad de
decisión de los interesados, no por un número predefinido.
\end{politicaacademica}

% ---------------------------------------------------------------------
\subsection{Registro de seguimiento (TMC2–TMC3)}

\begin{instruccion}
Este es el registro operativo del seguimiento: una fila por periodo. Compare lo previsto con lo
real, anote lo que se observó y —lo más importante— las acciones correctivas acordadas con su
responsable y fecha objetivo. La última columna obliga a pronunciarse sobre si la desviación
exige replanificar.

Si emplea una herramienta de gestión, puede sustituir este cuadro por una referencia a ella,
siempre que indique dónde está, quién la mantiene y que contenga la misma información.
\end{instruccion}

{\footnotesize
\begin{xltabular}{\textwidth}{|L{2.9cm}|Y|Y|}
\caption{Registro de seguimiento por periodo.}\label{tab:seguimiento}\\ \hline
\endfirsthead
\multicolumn{3}{@{}l@{}}{\footnotesize\itshape\tablename~\thetable{} (continuación)}\\ \hline
\endhead
\multicolumn{3}{@{}r@{}}{\footnotesize\itshape continúa en la página siguiente}\\
\endfoot
\endlastfoot
  \filaenc \textbf{Campo} & \textbf{Periodo \campo{1}} & \textbf{Periodo \campo{2}} \\ \hline
  Periodo (fechas)            & \campo{dd/mm – dd/mm} & \campo{} \\ \hline
  Avance previsto             & \campo{\% o hitos}    & \campo{} \\ \hline
  Avance real                 & \campo{\% o hitos}    & \campo{} \\ \hline
  Pruebas ejecutadas          & \campo{n}             & \campo{} \\ \hline
  Aprobadas                   & \campo{n}             & \campo{} \\ \hline
  Fallidas                    & \campo{n}             & \campo{} \\ \hline
  Bloqueadas                  & \campo{n y causa}     & \campo{} \\ \hline
  Cobertura alcanzada         & \campo{COB-xx = \%}   & \campo{} \\ \hline
  Incidentes relevantes       & \campo{INC-xx}        & \campo{} \\ \hline
  Riesgos nuevos              & \campo{Clave y descripción} & \campo{} \\ \hline
  Esfuerzo utilizado          & \campo{h reales / h estimadas} & \campo{} \\ \hline
  Bloqueos y su causa         & \campo{Descripción}   & \campo{} \\ \hline
  Acción correctiva acordada  & \campo{Qué se hará}   & \campo{} \\ \hline
  Responsable de la acción    & \campo{Rol}           & \campo{} \\ \hline
  Fecha objetivo              & \campo{dd/mm/aaaa}    & \campo{} \\ \hline
  ¿Requiere replanificación?  & \campo{Sí/No + motivo}& \campo{} \\ \hline
\end{xltabular}
}

El Cuadro~\ref{tab:seguimiento} es el instrumento que conecta lo planificado con lo observado
periodo a periodo. Duplique la columna tantas veces como periodos tenga el esfuerzo de prueba,
o use una fila por periodo si prefiere una disposición horizontal.

% ---------------------------------------------------------------------
\subsection{Acciones de control y replanificación (TMC3)}

\begin{instruccion}
Defina de antemano qué desviación dispara qué respuesta y quién la autoriza. Distinga las
acciones que \emph{corrigen la ejecución} (repriorizar, reasignar, ampliar un ciclo) de las que
\emph{reabren la planificación} (cambiar el alcance, los niveles, los criterios de salida o el
calendario). Sin esta regla acordada, cada desviación se discute desde cero.
\end{instruccion}

\begin{table}[H]
\centering
\caption{Reglas de control y disparadores de replanificación.}
\label{tab:control}
\begin{tabularx}{\textwidth}{|L{3.6cm}|Y|C{2.6cm}|C{2.0cm}|}
  \hline
  \filaenc \textbf{Situación observada} & \textbf{Respuesta acordada} & \textbf{¿Reabre la planificación?} & \textbf{Autoriza} \\ \hline
  \campo{Desviación de avance > 20 \%} & \campo{Acción} & \campo{No} & \campo{Rol} \\ \hline
  \campo{Entorno indisponible > 3 días} & \campo{Acción} & \campo{Sí} & \campo{Rol} \\ \hline
  \campo{Incidente crítico bloquea un objeto Catastrófico} & \campo{Acción} & \campo{Sí} & \campo{Rol} \\ \hline
  & & & \\ \hline
\end{tabularx}
\end{table}

El Cuadro~\ref{tab:control} cierra el bucle de control representado en la
Figura~\ref{fig:procesos}: toda replanificación acordada aquí debe producir una nueva versión
del plan en el Cuadro~\ref{tab:historial} y, si procede, una entrada en el
Cuadro~\ref{tab:desviaciones}.

% ---------------------------------------------------------------------
\subsection{Informes de estado (TMC4)}
\label{sec:informes}

\begin{instruccion}
Indique quién recibe informes de estado, con qué periodicidad y por qué medio; los
destinatarios y su frecuencia salen del Cuadro~\ref{tab:interesados}. El formato reutilizable
del informe está en el \ref{anx:informes}: una copia por periodo.
\end{instruccion}

\begin{table}[H]
\centering
\caption{Plan de comunicación de los informes de estado.}
\label{tab:comunicacion}
\begin{tabularx}{\textwidth}{|L{3.4cm}|C{2.4cm}|Y|C{2.1cm}|}
  \hline
  \filaenc \textbf{Destinatario} & \textbf{Periodicidad} & \textbf{Contenido mínimo} & \textbf{Medio} \\ \hline
  \campo{Interesado} & \campo{Semanal} & \campo{Avance, cobertura, incidentes, riesgos, acciones} & \campo{Medio} \\ \hline
  & & & \\ \hline
\end{tabularx}
\end{table}

El Cuadro~\ref{tab:comunicacion} evita que el seguimiento quede confinado al equipo de prueba:
determina qué información sale del equipo, hacia quién y cuándo.

% =====================================================================
%  6. Gestión de incidentes   (proceso IR de ISO/IEC/IEEE 29119-2)
% =====================================================================
\section{Gestión de incidentes de prueba}
\label{sec:gestion-incidentes}

\begin{notanormativa}[Apartado nuevo respecto de la plantilla original]
La plantilla original no preveía gestión de incidentes. Sin ella, un resultado fallido no
produce ninguna consecuencia documentada: no puede comunicarse, priorizarse, seguirse ni
cerrarse. Este apartado desarrolla el proceso de \textbf{reporte de incidentes (IR1–IR2)} de
\normap{2}. El formato de registro está en el \ref{anx:incidentes}.
\end{notanormativa}

% ---------------------------------------------------------------------
\subsection{Análisis previo al reporte (IR1)}

\begin{notanormativa}[Un resultado fallido no es automáticamente un defecto del producto]
La actividad IR1 exige, antes de crear un reporte, determinar el origen de lo observado. Un
resultado inesperado puede deberse a:
\begin{itemize}
  \item un \textbf{defecto del objeto de prueba} —el caso que interesa al producto—;
  \item un \textbf{defecto del entorno} (configuración, versión desplegada, servicio caído);
  \item un \textbf{defecto de los datos} de prueba;
  \item un \textbf{error del propio caso de prueba} (resultado esperado mal derivado);
  \item un \textbf{incidente ya conocido}, que debe actualizarse en lugar de duplicarse.
\end{itemize}
Omitir esta discriminación infla el conteo de defectos y desvía el esfuerzo de corrección.
\end{notanormativa}

\begin{instruccion}
Indique quién realiza este análisis, en qué momento y qué evidencia debe acompañarlo. Registre
el origen determinado en la columna correspondiente del \ref{anx:incidentes}.
\end{instruccion}

\campo{Describa el procedimiento de análisis previo al reporte}

% ---------------------------------------------------------------------
\subsection{Severidad y prioridad}

\begin{notanormativa}[Son dimensiones independientes]
La \textbf{severidad} mide el impacto técnico o de negocio del incidente. La \textbf{prioridad}
mide la urgencia de atenderlo. Son independientes: un incidente puede ser de severidad baja y
prioridad alta (un error ortográfico en la pantalla de inicio antes de una demostración), o de
severidad alta y prioridad baja (un fallo grave en una función que no se usará en el periodo).
Colapsarlas en un solo campo impide priorizar el trabajo de corrección de forma defendible.
\end{notanormativa}

\begin{instruccion}
Defina las escalas antes de registrar el primer incidente y aplíquelas de forma consistente.
Indique también quién asigna cada valor: normalmente la severidad la propone quien prueba y la
prioridad la decide quien decide sobre el producto.
\end{instruccion}

\begin{table}[H]
\centering
\caption{Escala de severidad acordada.}
\label{tab:severidad}
\begin{tabularx}{\textwidth}{|C{2.2cm}|Y|C{2.4cm}|}
  \hline
  \filaenc \textbf{Nivel} & \textbf{Significado: impacto observable} & \textbf{Quién la asigna} \\ \hline
  \campo{Crítica} & \campo{Criterio} & \campo{Rol} \\ \hline
  \campo{Alta}    & \campo{Criterio} & \campo{Rol} \\ \hline
  \campo{Media}   & \campo{Criterio} & \campo{Rol} \\ \hline
  \campo{Baja}    & \campo{Criterio} & \campo{Rol} \\ \hline
\end{tabularx}
\end{table}

\begin{table}[H]
\centering
\caption{Escala de prioridad acordada.}
\label{tab:prioridad}
\begin{tabularx}{\textwidth}{|C{2.2cm}|Y|C{2.4cm}|}
  \hline
  \filaenc \textbf{Nivel} & \textbf{Significado: urgencia de atención} & \textbf{Quién la asigna} \\ \hline
  \campo{Inmediata} & \campo{Criterio y plazo} & \campo{Rol} \\ \hline
  \campo{Alta}      & \campo{Criterio y plazo} & \campo{Rol} \\ \hline
  \campo{Normal}    & \campo{Criterio y plazo} & \campo{Rol} \\ \hline
  \campo{Baja}      & \campo{Criterio y plazo} & \campo{Rol} \\ \hline
\end{tabularx}
\end{table}

Los Cuadros~\ref{tab:severidad} y~\ref{tab:prioridad} definen dos ejes distintos de
clasificación; ambos se registran, por separado, en cada incidente del \ref{anx:incidentes}.

% ---------------------------------------------------------------------
\subsection{Estados y flujo de tratamiento}

\begin{instruccion}
Defina los estados por los que puede pasar un incidente y quién autoriza cada transición. Un
estado sin responsable es un incidente que nadie mueve. Incluya explícitamente el estado
terminal que corresponda cuando el defecto \emph{no} se corrige, para que los incidentes
abiertos al cierre queden documentados en lugar de desaparecer.
\end{instruccion}

\begin{table}[H]
\centering
\caption{Estados de un incidente y responsables de cada transición.}
\label{tab:estados-incidente}
\begin{tabularx}{\textwidth}{|C{2.4cm}|Y|C{2.4cm}|C{2.4cm}|}
  \hline
  \filaenc \textbf{Estado} & \textbf{Significado} & \textbf{Transiciona a} & \textbf{Responsable} \\ \hline
  Nuevo        & Registrado, aún no analizado por el equipo receptor. & \campo{Abierto / Rechazado / Duplicado} & \campo{Rol} \\ \hline
  Abierto      & Aceptado como incidente válido y pendiente de tratamiento. & \campo{En corrección / Diferido} & \campo{Rol} \\ \hline
  En corrección& El equipo desarrollador trabaja en la resolución. & \campo{Resuelto} & \campo{Rol} \\ \hline
  Resuelto     & Se entregó una corrección; pendiente de confirmación. & \campo{Cerrado / Reabierto} & \campo{Rol} \\ \hline
  Cerrado      & Prueba de confirmación aprobada (\ref{sec:confirmacion-regresion}). & --- & \campo{Rol} \\ \hline
  Diferido     & Válido, pero no se corregirá en este periodo. & \campo{Riesgo residual} & \campo{Rol} \\ \hline
  Rechazado    & No constituye un defecto del objeto de prueba (véase IR1). & --- & \campo{Rol} \\ \hline
  Duplicado    & Ya registrado; se actualiza el incidente original. & --- & \campo{Rol} \\ \hline
\end{tabularx}
\end{table}

El Cuadro~\ref{tab:estados-incidente} define el ciclo de vida del incidente. Los estados
\emph{Diferido} y \emph{Abierto} al cierre del esfuerzo alimentan directamente los incidentes
abiertos (apartado~\ref{sec:incidentes-abiertos}) y los riesgos residuales
(apartado~\ref{sec:residuales}).

\begin{politicaacademica}
Dado que el curso establece que el equipo de prueba no repara los defectos encontrados, es
previsible que muchos incidentes terminen el periodo en estado \emph{Abierto} o
\emph{Diferido}. Eso no es un fallo del proceso: lo que sí debe evitarse es cerrarlos sin
corrección ni confirmación, porque haría irreconocible la diferencia entre un defecto resuelto
y uno simplemente no atendido.
\end{politicaacademica}

% ---------------------------------------------------------------------
\subsection{Comunicación de incidentes}

\begin{instruccion}
Indique cómo se comunican los incidentes al equipo desarrollador y a los demás interesados: por
qué medio, con qué plazo según la severidad, y quién es el punto de contacto de cada parte.
\end{instruccion}

\begin{table}[H]
\centering
\caption{Comunicación de incidentes según severidad.}
\label{tab:comunicacion-incidentes}
\begin{tabularx}{\textwidth}{|C{2.2cm}|C{3.0cm}|Y|C{2.4cm}|}
  \hline
  \filaenc \textbf{Severidad} & \textbf{Destinatario} & \textbf{Medio} & \textbf{Plazo} \\ \hline
  \campo{Crítica} & \campo{Rol} & \campo{Medio} & \campo{Plazo} \\ \hline
  \campo{Otras}   & \campo{Rol} & \campo{Medio} & \campo{Plazo} \\ \hline
\end{tabularx}
\end{table}

El Cuadro~\ref{tab:comunicacion-incidentes} fija los plazos de aviso; su cumplimiento es uno de
los aspectos que el seguimiento (apartado~\ref{sec:seguimiento}) puede monitorear.

% =====================================================================
%  7. Actividades de prueba y estimaciones   (actividad TP6)
% =====================================================================
\section{Actividades de prueba y estimaciones}
\label{sec:actividades}

\begin{instruccion}
Identifique las actividades necesarias para el proceso de prueba que se aplicará al caso de
estudio y estime el esfuerzo neto de cada una, en días u horas.

\textbf{Ampliación respecto de la plantilla original:} añada las \emph{dependencias} entre
actividades y el \emph{producto} que cada una genera. Las dependencias son las que hacen
realista el cronograma del apartado~\ref{sec:calendario}; el producto es lo que permite
comprobar que la actividad terminó.

Las actividades deberían cubrir el ciclo completo: planificación, preparación de entorno y
datos, diseño e implementación, ejecución, reporte de incidentes, seguimiento y finalización.
Puede apoyarse en la Figura~\ref{fig:procesos} para no omitir ninguna.
\end{instruccion}

\begin{table}[H]
\centering
\caption{Actividades de prueba, esfuerzo estimado, dependencias y productos.}
\label{tab:actividades}
\begin{tabularx}{\textwidth}{|C{1.4cm}|Y|C{1.8cm}|C{1.8cm}|C{2.5cm}|}
  \hline
  \filaenc \textbf{Clave} & \textbf{Actividad} & \textbf{Esfuerzo estimado} & \textbf{Depende de} & \textbf{Producto que genera} \\ \hline
  \campo{ACT-01} & \campo{Elaborar el plan de prueba} & \campo{h/días} & \campo{---} & \campo{Parte I} \\ \hline
  \campo{ACT-02} & \campo{Preparar entorno y datos} & \campo{h/días} & \campo{ACT-01} & \campo{Anexo F} \\ \hline
  \campo{ACT-03} & \campo{Diseñar condiciones y casos} & \campo{h/días} & \campo{ACT-01} & \campo{Anexo A} \\ \hline
  \campo{ACT-04} & \campo{Elaborar procedimientos} & \campo{h/días} & \campo{ACT-03} & \campo{Anexo B} \\ \hline
  \campo{ACT-05} & \campo{Ejecutar y registrar} & \campo{h/días} & \campo{ACT-02, ACT-04} & \campo{Anexo C} \\ \hline
  \campo{ACT-06} & \campo{Reportar y dar seguimiento a incidentes} & \campo{h/días} & \campo{ACT-05} & \campo{Anexo D} \\ \hline
  \campo{ACT-07} & \campo{Monitorear, controlar e informar} & \campo{h/días} & \campo{Concurrente} & \campo{Anexo E} \\ \hline
  \campo{ACT-08} & \campo{Finalizar: archivar, limpiar, reportar} & \campo{h/días} & \campo{ACT-05, ACT-06} & \campo{Parte II} \\ \hline
  & & & & \\ \hline
\end{tabularx}
\notatabla{Las actividades anteriores son un punto de partida derivado del modelo de procesos; adáptelas al proyecto. La actividad de seguimiento se marca como concurrente porque no ocupa un lugar fijo en la secuencia.}
\end{table}

El Cuadro~\ref{tab:actividades} sostiene tanto la estimación de esfuerzo como el cronograma del
apartado~\ref{sec:calendario}; su columna de esfuerzo estimado es además la línea base contra
la que el seguimiento compara el consumo real.

\begin{instruccion}
Indique el método de estimación empleado (juicio experto, analogía con proyectos previos,
descomposición) y los supuestos en que se apoya. Una estimación sin método declarado no puede
revisarse cuando se desvía.
\end{instruccion}

\campo{Método de estimación y supuestos}

% =====================================================================
%  8. Equipo de trabajo   (actividad TP6)
% =====================================================================
\section{Equipo de trabajo}
\label{sec:equipo}

\begin{instruccion}
Indique las personas integrantes del equipo de prueba, los roles que cubrirán, sus
responsabilidades y su autoridad dentro del proceso.

\textbf{Ampliación respecto de la plantilla original:} añada las \emph{competencias} requeridas
por cada rol, la \emph{disponibilidad} real de cada persona y las necesidades de
\emph{formación} detectadas. La tarea TP6 de \normap{2} incluye explícitamente definir
competencias y necesidades de formación; una brecha entre competencia requerida y disponible es
un riesgo de proyecto y debería figurar en el Cuadro~\ref{tab:riesgos-proyecto}.
\end{instruccion}

\begin{table}[H]
\centering
\caption{Integrantes del equipo de prueba, roles, responsabilidades y autoridad.}
\label{tab:equipo}
\begin{tabularx}{\textwidth}{|L{3.2cm}|C{2.3cm}|Y|C{2.6cm}|}
  \hline
  \filaenc \textbf{Integrante} & \textbf{Rol} & \textbf{Responsabilidades} & \textbf{Autoridad (qué decide)} \\ \hline
  \campo{Matrícula y nombre} & \campo{Líder de prueba} & \campo{Responsabilidades} & \campo{Aprueba / propone} \\ \hline
  \campo{Matrícula y nombre} & \campo{Diseñador de pruebas} & \campo{Responsabilidades} & \campo{} \\ \hline
  \campo{Matrícula y nombre} & \campo{Ejecutor} & \campo{Responsabilidades} & \campo{} \\ \hline
  & & & \\ \hline
\end{tabularx}
\end{table}

\begin{table}[H]
\centering
\caption{Competencias requeridas, disponibilidad y necesidades de formación.}
\label{tab:competencias}
\begin{tabularx}{\textwidth}{|C{2.3cm}|Y|C{2.4cm}|C{3.0cm}|}
  \hline
  \filaenc \textbf{Rol} & \textbf{Competencia requerida} & \textbf{Disponibilidad} & \textbf{Formación necesaria} \\ \hline
  \campo{Rol} & \campo{Competencia} & \campo{h/semana} & \campo{Acción y fecha} \\ \hline
  & & & \\ \hline
\end{tabularx}
\notatabla{Toda brecha de competencia o de disponibilidad que no pueda cubrirse debe registrarse como riesgo de proyecto.}
\end{table}

Los Cuadros~\ref{tab:equipo} y~\ref{tab:competencias} responden a preguntas complementarias:
quién hace qué y con qué autoridad, y si el equipo cuenta efectivamente con la capacidad y el
tiempo para hacerlo. La disponibilidad declarada debe ser coherente con el esfuerzo estimado en
el Cuadro~\ref{tab:actividades}.

% =====================================================================
%  9. Calendarización   (actividad TP6)
% =====================================================================
\section{Calendarización}
\label{sec:calendario}

\begin{instruccion}
Elabore el cronograma a partir de las actividades y duraciones del
Cuadro~\ref{tab:actividades}, de sus dependencias y de la disponibilidad declarada en el
apartado~\ref{sec:equipo}. Incluya revisiones, correcciones, entregas
y evaluaciones, así como los hitos del Cuadro~\ref{tab:hitos}.

Puede sustituir el cuadro siguiente por un diagrama de Gantt, siempre que éste lleve caption,
etiqueta y una mención desde el texto. Si lo incluye como imagen, colóquela en
\texttt{diagrams/} o en una carpeta \texttt{img/} y refiérala con ruta relativa.
\end{instruccion}

\begin{table}[H]
\centering
\caption{Cronograma planificado del esfuerzo de prueba.}
\label{tab:cronograma}
\begin{tabularx}{\textwidth}{|C{1.4cm}|Y|C{2.5cm}|C{2.5cm}|C{2.4cm}|}
  \hline
  \filaenc \textbf{Clave} & \textbf{Actividad o hito} & \textbf{Inicio} & \textbf{Fin} & \textbf{Responsable} \\ \hline
  \campo{ACT-01} & \campo{Actividad} & \campo{dd/mm/aaaa} & \campo{dd/mm/aaaa} & \campo{Rol} \\ \hline
  \campo{H-01}   & \campo{Primera entrega} & \campo{dd/mm/aaaa} & \campo{dd/mm/aaaa} & \campo{Rol} \\ \hline
  & & & & \\ \hline
  & & & & \\ \hline
\end{tabularx}
\end{table}

El Cuadro~\ref{tab:cronograma} sitúa en el tiempo las actividades y los hitos previstos, y
constituye la línea base temporal contra la que el apartado~\ref{sec:seguimiento} compara el
avance real. Al cierre, se contrasta con las fechas efectivamente ejecutadas en el
apartado~\ref{sec:resumen-final}.


% ---------------------------------------------------------------------
% PARTE II. REPORTE DE FINALIZACIÓN DE LA PRUEBA
% ---------------------------------------------------------------------
\parteplan{PARTE II}{REPORTE DE FINALIZACIÓN DE LA PRUEBA}

% =====================================================================
%  PARTE II. Reporte de finalización   (proceso TC de ISO/IEC/IEEE 29119-2)
% =====================================================================

\begin{notanormativa}[Correspondencia con el modelo de procesos]
Esta parte desarrolla el proceso de \textbf{finalización de prueba (TC1–TC4)} de \normap{2}:
archivar los activos reutilizables (TC1), limpiar y liberar el entorno (TC2), identificar
lecciones aprendidas (TC3) y reportar la finalización (TC4). Es, según el documento de procesos
consultado, el proceso más frecuentemente omitido en la práctica, pese a ser el que sostiene la
mejora continua.
\end{notanormativa}

\begin{notanormativa}[Seis preguntas distintas en el cierre]
El cierre responde a seis preguntas que conviene no mezclar, y que aquí ocupan apartados
distintos:
\begin{enumerate}
  \item qué actividades se completaron (\ref{sec:resumen-final});
  \item qué resultados se obtuvieron (\ref{sec:resultados});
  \item si se cumplieron los criterios de salida (\ref{sec:evaluacion-salida});
  \item qué incidentes quedan abiertos (\ref{sec:incidentes-abiertos});
  \item qué riesgos residuales persisten (\ref{sec:residuales});
  \item la decisión de aceptación y quién la toma (\ref{sec:aceptacion}).
\end{enumerate}
Haber ejecutado el 100\,\% de los casos —aun si muchos fallaron— no demuestra que el producto
sea aceptable.
\end{notanormativa}

\begin{observacion}[«Evaluación de la finalización» basada en cobertura]
\obsproblema{La plantilla anterior resolvía el cierre con un único apartado que determinaba «el
nivel de garantía de calidad logrado con base a la cobertura de las pruebas realizadas», sin
distinguir actividades completadas, resultados, incidentes abiertos, riesgo residual ni
autoridad de aceptación.}
\obscambio{El cierre se desglosa en los seis apartados enumerados arriba, con un cuadro de
decisión de aceptación que separa la recomendación técnica del equipo de la decisión de quien
asume el riesgo residual.}
\obsfuente{OBS §4.6}
\end{observacion}

% ---------------------------------------------------------------------
\section{Resumen de las pruebas realizadas}
\label{sec:resumen-final}

\begin{instruccion}
Compare las fechas planificadas con las realmente ejecutadas y presente el cronograma con
tiempos reales. Indique las incidencias que afectaron el progreso, sus causas y cómo afectaron
la ejecución del plan. Señale además qué actividades del Cuadro~\ref{tab:actividades} se
completaron, cuáles se completaron parcialmente y cuáles no se realizaron, con su motivo.
\end{instruccion}

\begin{table}[H]
\centering
\caption{Ejecución real frente a lo planificado.}
\label{tab:real-vs-plan}
\begin{tabularx}{\textwidth}{|C{1.4cm}|L{3.2cm}|C{2.0cm}|C{2.0cm}|C{1.9cm}|Y|}
  \hline
  \filaenc \textbf{Clave} & \textbf{Actividad} & \textbf{Fin planificado} & \textbf{Fin real} & \textbf{Estado} & \textbf{Causa de la desviación} \\ \hline
  \campo{ACT-01} & \campo{Actividad} & \campo{dd/mm} & \campo{dd/mm} & \campo{Completada} & \campo{---} \\ \hline
  & & & & & \\ \hline
\end{tabularx}
\notatabla{Estado: \textit{Completada} · \textit{Parcial} · \textit{No realizada} · \textit{Cancelada}.}
\end{table}

El Cuadro~\ref{tab:real-vs-plan} documenta qué se hizo realmente y por qué se desvió del plan.
Responde únicamente a la pregunta de si las actividades se completaron; no dice todavía nada
sobre la calidad del producto.

% ---------------------------------------------------------------------
\section{Resultados obtenidos}
\label{sec:resultados}

\begin{instruccion}
Presente los resultados agregados de la ejecución, tomados del \ref{anx:ejecuciones} y del
Cuadro~\ref{tab:cobertura-base}. El cuadro está dividido en los tres bloques que el
apartado~\ref{sec:cobertura} mantiene separados —avance, cobertura y veredicto—; respete la
división al rellenarlo y no traslade cifras de un bloque a otro.

En particular, la cobertura se calcula sobre lo \emph{ejercitado}: un elemento probado por un
caso que falló cuenta en el bloque de cobertura \textbf{y} en el de resultados. No lo descuente
del primero.
\end{instruccion}

\begin{table}[H]
\centering
\caption{Resultados agregados de la ejecución.}
\label{tab:resultados}
\begin{tabularx}{\textwidth}{|L{5.2cm}|C{2.6cm}|Y|}
  \hline
  \filaenc \textbf{Indicador} & \textbf{Valor} & \textbf{Comentario} \\ \hline
  \multicolumn{3}{|l|}{\cellcolor{uvgrisc}\textbf{Avance} — ¿cuánto del trabajo previsto se realizó?} \\ \hline
  Casos de prueba especificados        & \campo{n} & \campo{} \\ \hline
  Ejecuciones realizadas               & \campo{n} & \campo{Incluye repeticiones y regresiones} \\ \hline
  Casos bloqueados o no ejecutados     & \campo{n} & \campo{Motivo} \\ \hline
  Avance de ejecución (\%)             & \campo{\%} & Casos ejecutados ÷ casos planificados \\ \hline
  \multicolumn{3}{|l|}{\cellcolor{uvgrisc}\textbf{Cobertura} — ¿qué elementos quedaron ejercitados?} \\ \hline
  \campo{COB-01} — cobertura alcanzada & \campo{\%} & Elementos ejercitados, \emph{con independencia del veredicto}; fuente: Cuadro~\ref{tab:cobertura-base} \\ \hline
  \campo{COB-02} — cobertura alcanzada & \campo{\%} & \campo{Según su definición en el Cuadro~\ref{tab:cobertura}} \\ \hline
  Elementos del alcance sin ejercitar  & \campo{n} & \campo{Brechas de cobertura; deben reaparecer en \ref{sec:residuales}} \\ \hline
  \multicolumn{3}{|l|}{\cellcolor{uvgrisc}\textbf{Resultados} — de lo ejercitado, ¿qué se comportó como se esperaba?} \\ \hline
  Veredicto \emph{aprobado}            & \campo{n} & \campo{} \\ \hline
  Veredicto \emph{no aprobado}         & \campo{n} & \campo{Origina los incidentes del \ref{anx:incidentes}} \\ \hline
  Veredicto \emph{no concluyente}      & \campo{n} & \campo{} \\ \hline
\end{tabularx}
\end{table}

El Cuadro~\ref{tab:resultados} separa en tres bloques el avance, la cobertura y los veredictos.
Leídos juntos
responden a la pregunta que interesa al cierre: cuánto se probó, sobre qué y con qué resultado.
Ninguno de los tres bloques basta por sí solo para sostener la decisión de aceptación del
apartado~\ref{sec:aceptacion}.

% ---------------------------------------------------------------------
\section{Evaluación frente a los criterios de salida}
\label{sec:evaluacion-salida}

\begin{instruccion}
Evalúe uno a uno los criterios de finalización definidos en el
Cuadro~\ref{tab:criterios-salida}, indicando el valor alcanzado y si el criterio se cumple. No
reformule aquí los criterios: si alguno resultó inaplicable, declárelo como desviación en el
Cuadro~\ref{tab:desviaciones}.

Tome cada valor del bloque que le corresponde en el Cuadro~\ref{tab:resultados}: los criterios
de avance, de cobertura y de resultados se evalúan con cifras distintas. Un criterio de
cobertura cumplido al 100\,\% no compensa un criterio de resultados incumplido.
\end{instruccion}

\begin{instruccion}[Advertencia]
Si algún criterio de salida quedó redactado sobre «casos aprobados» pero clasificado como
criterio de \emph{cobertura}, corríjalo antes de evaluarlo: pertenece al tipo
\emph{Resultados}.
\end{instruccion}

\begin{table}[H]
\centering
\caption{Cumplimiento de los criterios de finalización.}
\label{tab:eval-salida}
\begin{tabularx}{\textwidth}{|C{1.5cm}|C{2.2cm}|Y|C{1.9cm}|C{1.9cm}|C{1.9cm}|}
  \hline
  \filaenc \textbf{Criterio} & \textbf{Tipo} & \textbf{Enunciado} & \textbf{Valor objetivo} & \textbf{Valor alcanzado} & \textbf{¿Se cumple?} \\ \hline
  \campo{CS-01} & \campo{Cobertura} & \campo{Enunciado} & \campo{Objetivo} & \campo{Real} & \campo{Sí/No} \\ \hline
  \campo{CS-02} & \campo{Avance} & \campo{Enunciado} & \campo{Objetivo} & \campo{Real} & \campo{Sí/No} \\ \hline
  \campo{CS-03} & \campo{Resultados} & \campo{Enunciado} & \campo{Objetivo} & \campo{Real} & \campo{Sí/No} \\ \hline
  \campo{CS-04} & \campo{Incidentes} & \campo{Enunciado} & \campo{Objetivo} & \campo{Real} & \campo{Sí/No} \\ \hline
  & & & & & \\ \hline
\end{tabularx}
\notatabla{La columna «Tipo» debe coincidir con la del Cuadro~\ref{tab:criterios-salida}. Sirve para comprobar de un vistazo que el cierre no se sostiene sobre un solo tipo de criterio.}
\end{table}

El Cuadro~\ref{tab:eval-salida} establece si las \emph{actividades de prueba} pueden darse por
concluidas, y la columna «Tipo» deja ver sobre qué se apoya esa conclusión. Es un insumo de la
decisión de aceptación del apartado~\ref{sec:aceptacion}, pero no la sustituye: cumplir los
criterios de salida significa que se hizo lo previsto, no que el resultado sea aceptable.

% ---------------------------------------------------------------------
\section{Incidentes abiertos al cierre}
\label{sec:incidentes-abiertos}

\begin{instruccion}
Relacione los incidentes que permanecen sin cerrar, con su severidad, prioridad y estado. Estos
incidentes son el principal insumo de los riesgos residuales y de la decisión de aceptación:
omitirlos convierte el cierre en una formalidad.
\end{instruccion}

\begin{table}[H]
\centering
\caption{Incidentes abiertos o diferidos al cierre del esfuerzo de prueba.}
\label{tab:incidentes-abiertos}
\begin{tabularx}{\textwidth}{|C{1.5cm}|Y|C{1.8cm}|C{1.8cm}|C{1.9cm}|C{2.2cm}|}
  \hline
  \filaenc \textbf{Clave} & \textbf{Resumen} & \textbf{Severidad} & \textbf{Prioridad} & \textbf{Estado} & \textbf{Disposición acordada} \\ \hline
  \campo{INC-01} & \campo{Resumen} & \campo{Alta} & \campo{Normal} & \campo{Diferido} & \campo{Qué se hará después} \\ \hline
  & & & & & \\ \hline
\end{tabularx}
\end{table}

El Cuadro~\ref{tab:incidentes-abiertos} es el extracto de cierre del \ref{anx:incidentes}; cada
fila debería tener su correspondencia en el Cuadro~\ref{tab:riesgos-residuales}.

% ---------------------------------------------------------------------
\section{Métricas de prueba}
\label{sec:metricas-finales}

\begin{instruccion}
Presente el cálculo final de las métricas definidas en el Cuadro~\ref{tab:metricas}, junto con
su análisis y discusión. Interprete cada valor según lo acordado en su definición y señale
explícitamente lo que la métrica \emph{no} permite concluir.
\end{instruccion}

\begin{table}[H]
\centering
\caption{Valores finales de las métricas de prueba.}
\label{tab:metricas-finales}
\begin{tabularx}{\textwidth}{|C{1.5cm}|C{1.7cm}|C{2.0cm}|Y|}
  \hline
  \filaenc \textbf{Clave} & \textbf{Tipo} & \textbf{Valor final} & \textbf{Análisis e interpretación} \\ \hline
  \campo{MET-01} & \campo{Producto} & \campo{Valor} & \campo{Qué indica y qué no} \\ \hline
  & & & \\ \hline
\end{tabularx}
\end{table}

El Cuadro~\ref{tab:metricas-finales} cierra el ciclo abierto en el apartado~\ref{sec:metricas}:
las métricas se interpretan según el criterio acordado \emph{antes} de conocer los resultados.

% ---------------------------------------------------------------------
\section{Riesgos residuales}
\label{sec:residuales}

\begin{instruccion}
Analice qué riesgos pudieron resolverse o evitarse y cuáles persisten y podrían afectar
procesos posteriores. Incluya al menos tres fuentes: los riesgos de producto no mitigados
(Cuadro~\ref{tab:riesgos-producto}), las exclusiones del alcance
(Cuadro~\ref{tab:exclusiones}) y los incidentes abiertos
(Cuadro~\ref{tab:incidentes-abiertos}).
\end{instruccion}

\begin{table}[H]
\centering
\caption{Riesgos residuales al cierre del esfuerzo de prueba.}
\label{tab:riesgos-residuales}
\begin{tabularx}{\textwidth}{|C{1.5cm}|C{2.0cm}|Y|C{1.9cm}|C{2.6cm}|}
  \hline
  \filaenc \textbf{Clave} & \textbf{Origen} & \textbf{Riesgo que persiste y su posible consecuencia} & \textbf{Exposición} & \textbf{Recomendación} \\ \hline
  \campo{RR-01} & \campo{RP-04 / EX-01 / INC-02} & \campo{Descripción} & \campo{Alta} & \campo{Acción sugerida} \\ \hline
  & & & & \\ \hline
\end{tabularx}
\end{table}

El Cuadro~\ref{tab:riesgos-residuales} declara lo que el esfuerzo de prueba \emph{no} llegó a
descartar. Es la información que la autoridad de aceptación necesita para decidir con
conocimiento de causa.

% ---------------------------------------------------------------------
\section{Decisión de aceptación}
\label{sec:aceptacion}

\begin{notanormativa}[Finalizar no es aceptar]
Dar por terminadas las actividades de prueba (\ref{sec:evaluacion-salida}) es una decisión
sobre el \emph{proceso}. Aceptar el producto es una decisión sobre el \emph{riesgo residual} que
la organización está dispuesta a asumir. Son decisiones distintas, se toman con información
distinta y, con frecuencia, corresponden a personas distintas, por lo que el plan debe nombrar
una autoridad de aceptación.
\end{notanormativa}

\begin{instruccion}
Registre la decisión, quién la toma y sobre qué base. La decisión debe considerar los resultados
obtenidos, la cobertura relevante alcanzada, los incidentes abiertos con su severidad y los
riesgos residuales. Si el equipo de prueba sólo emite una recomendación y la decisión
corresponde a otra parte, dígalo así.
\end{instruccion}

\begin{table}[H]
\centering
\caption{Decisión de aceptación del resultado de la prueba.}
\label{tab:aceptacion}
\begin{tabularx}{\textwidth}{|E{4.6cm}|Y|}
  \hline
  Recomendación del equipo de prueba & \campo{Aceptar / aceptar con condiciones / no aceptar} \\ \hline
  Base de la recomendación & \campo{Resultados, cobertura, incidentes abiertos y riesgos residuales considerados} \\ \hline
  Condiciones propuestas & \campo{Qué debería ocurrir antes de usar el software} \\ \hline
  Autoridad de aceptación & \autoridadAceptacion \\ \hline
  Decisión tomada & \campo{Decisión} \\ \hline
  Fecha & \campo{dd/mm/aaaa} \\ \hline
\end{tabularx}
\end{table}

El Cuadro~\ref{tab:aceptacion} separa explícitamente la recomendación técnica del equipo de
prueba de la decisión de quien tiene autoridad para asumir el riesgo residual.

\begin{politicaacademica}
En el contexto del curso, la autoridad de aceptación es \autoridadAceptacion, definida en
\texttt{config/metadata.tex}. Si el ejercicio académico no contempla una decisión real de
aceptación, indíquelo así en el cuadro y limite el contenido a la recomendación del equipo: es
preferible declarar que la decisión no aplica antes que simularla.
\end{politicaacademica}

% ---------------------------------------------------------------------
\section{Entregables de la prueba}

\begin{instruccion}
Relacione los entregables efectivamente producidos, comprometidos en el
apartado~\ref{sec:entregables} (Cuadro~\ref{tab:entregables}), indicando su
versión final y dónde se encuentran. Los procedimientos y casos de prueba aplicados se integran
en los anexos de este documento.
\end{instruccion}

\begin{table}[H]
\centering
\caption{Entregables producidos y su ubicación final.}
\label{tab:entregables-finales}
\begin{tabularx}{\textwidth}{|L{4.6cm}|C{2.0cm}|Y|C{1.9cm}|}
  \hline
  \filaenc \textbf{Entregable} & \textbf{Versión final} & \textbf{Ubicación} & \textbf{Estado} \\ \hline
  \campo{Entregable} & \campo{v1.0} & \campo{Ubicación} & \campo{Entregado} \\ \hline
  & & & \\ \hline
\end{tabularx}
\end{table}

El Cuadro~\ref{tab:entregables-finales} permite comprobar que lo prometido en el
Cuadro~\ref{tab:entregables} se produjo efectivamente.

% ---------------------------------------------------------------------
\section{Archivo, limpieza y activos reutilizables}
\label{sec:archivo}

\begin{notanormativa}[Correspondencia con el modelo de procesos]
Este apartado desarrolla las actividades \textbf{TC1 «Archivar los activos»} y \textbf{TC2
«Limpiar el entorno»} de \normap{2}: no basta con identificar qué activos son reutilizables, hay
que decir dónde se conservan, quién responde por ellos, qué ocurre con los datos sensibles y en
qué estado se deja el entorno.
\end{notanormativa}

\begin{observacion}[Archivo y limpieza sin custodia]
\obsproblema{La plantilla anterior pedía identificar los activos reutilizables, pero no su
ubicación, su responsable ni su periodo de conservación, y no contemplaba la liberación del
entorno ni el tratamiento de los datos sensibles al cierre.}
\obscambio{Dos cuadros: custodia de los activos archivados (TC1) y limpieza, liberación de
recursos, revocación de accesos y eliminación segura de datos (TC2).}
\obsfuente{OBS §3}
\end{observacion}

\begin{instruccion}
Complete los dos cuadros siguientes. El primero registra qué se conserva y bajo la
responsabilidad de quién; el segundo, qué se desmonta, libera o elimina. Ambos deben ser
coherentes con el control de configuración del Cuadro~\ref{tab:configuracion}.
\end{instruccion}

\begin{table}[H]
\centering
\caption{Activos de prueba conservados y su custodia.}
\label{tab:archivo}
\begin{tabularx}{\textwidth}{|L{3.4cm}|C{1.9cm}|Y|C{2.1cm}|C{2.0cm}|}
  \hline
  \filaenc \textbf{Activo} & \textbf{¿Reutilizable?} & \textbf{Dónde se archiva} & \textbf{Responsable} & \textbf{Periodo de conservación} \\ \hline
  \campo{Casos de prueba} & \campo{Sí} & \campo{Repositorio} & \campo{Rol} & \campo{Plazo} \\ \hline
  \campo{Procedimientos} & \campo{Sí} & \campo{Repositorio} & \campo{Rol} & \campo{Plazo} \\ \hline
  \campo{Guiones automatizados} & \campo{Sí/No} & \campo{Repositorio} & \campo{Rol} & \campo{Plazo} \\ \hline
  \campo{Evidencias de ejecución} & \campo{---} & \campo{Repositorio} & \campo{Rol} & \campo{Plazo} \\ \hline
  \campo{Registros de incidentes} & \campo{---} & \campo{Repositorio} & \campo{Rol} & \campo{Plazo} \\ \hline
  \campo{Datos de prueba} & \campo{Sí/No} & \campo{Repositorio} & \campo{Rol} & \campo{Plazo} \\ \hline
\end{tabularx}
\end{table}

\begin{table}[H]
\centering
\caption{Limpieza del entorno y liberación de recursos.}
\label{tab:limpieza}
\begin{tabularx}{\textwidth}{|L{3.8cm}|Y|C{2.1cm}|C{2.0cm}|}
  \hline
  \filaenc \textbf{Acción} & \textbf{Qué se hace concretamente} & \textbf{Responsable} & \textbf{Fecha} \\ \hline
  Restaurar el entorno & \campo{Estado al que se devuelve} & \campo{Rol} & \campo{dd/mm} \\ \hline
  Liberar recursos y licencias & \campo{Qué se libera} & \campo{Rol} & \campo{dd/mm} \\ \hline
  Tratar los datos sensibles & \campo{Eliminación segura, anonimización o devolución} & \campo{Rol} & \campo{dd/mm} \\ \hline
  Retirar accesos y credenciales & \campo{Qué se revoca} & \campo{Rol} & \campo{dd/mm} \\ \hline
\end{tabularx}
\notatabla{Si algún conjunto de datos contenía información personal, la eliminación segura no es opcional: documente el procedimiento aplicado.}
\end{table}

Los Cuadros~\ref{tab:archivo} y~\ref{tab:limpieza} distinguen lo que se \emph{conserva} de lo
que se \emph{desmonta}; las mejoras que se deriven de este cierre se registran en el
apartado~\ref{sec:lecciones}. Sin el segundo, un entorno de prueba con datos reales puede sobrevivir
al proyecto sin dueño ni control.

% ---------------------------------------------------------------------
\section{Lecciones aprendidas}
\label{sec:lecciones}

\begin{instruccion}
Discuta los resultados y plantee las conclusiones principales sobre el producto y sobre el
proceso. Evalúe el propio proceso de prueba e identifique áreas de mejora. Liste
recomendaciones justificadas y trazables hasta los resultados que las sustentan.

Indique, para cada lección, a quién se comunica y qué debería cambiar como consecuencia. Según
el modelo de procesos (actividad TC3), las lecciones aprendidas retroalimentan la capa
organizacional, tal como muestra el bucle externo de la Figura~\ref{fig:procesos}; una lección
que no llega a nadie no cierra ese bucle.
\end{instruccion}

\begin{table}[H]
\centering
\caption{Lecciones aprendidas y mejoras propuestas.}
\label{tab:lecciones}
\begin{tabularx}{\textwidth}{|C{1.4cm}|C{1.8cm}|Y|C{2.6cm}|C{2.0cm}|}
  \hline
  \filaenc \textbf{Clave} & \textbf{Ámbito} & \textbf{Lección y evidencia que la sustenta} & \textbf{Mejora propuesta} & \textbf{Se comunica a} \\ \hline
  \campo{LA-01} & \campo{Proceso} & \campo{Qué se aprendió y con qué evidencia} & \campo{Qué debería cambiar} & \campo{Destinatario} \\ \hline
  \campo{LA-02} & \campo{Producto} & \campo{Lección} & \campo{Mejora} & \campo{Destinatario} \\ \hline
  \campo{LA-03} & \campo{Proyecto} & \campo{Lección} & \campo{Mejora} & \campo{Destinatario} \\ \hline
  & & & & \\ \hline
\end{tabularx}
\notatabla{Ámbito: \textit{producto} (qué aprendimos sobre el software) · \textit{proceso} (sobre cómo probamos) · \textit{proyecto} (sobre cómo nos organizamos).}
\end{table}

El Cuadro~\ref{tab:lecciones} exige que cada lección esté sustentada en evidencia registrada y
tenga un destinatario concreto; así el cierre produce información utilizable en lugar de una
lista de buenas intenciones.


% ---------------------------------------------------------------------
% ANEXOS · formatos reutilizables de registro
% ---------------------------------------------------------------------
\iniciaranexos

% =====================================================================
%  ANEXO A. Especificación de casos de prueba
% =====================================================================
\section{Especificación de casos de prueba}
\label{anx:casos}

\begin{notanormativa}[Qué es y qué no es un caso de prueba]
Un caso de prueba describe \textbf{qué debería ejecutarse y qué debería ocurrir}. Es un
artefacto de diseño y se mantiene estable entre ejecuciones. Lo que \emph{sucedió} al ejecutarlo
—resultado real, veredicto, evidencia— no se escribe aquí: se registra en el
\ref{anx:ejecuciones}. Esa separación es la que permite ejecutar el mismo caso muchas veces sin
sobrescribir su historial.
\end{notanormativa}

\begin{instruccion}
Reproduzca la ficha siguiente una vez por cada caso de prueba. Un caso queda derivado de una
condición de prueba (apartado~\ref{sec:diseno}, Cuadro~\ref{tab:condiciones}), que a su vez
procede de la base de prueba y
de un riesgo: por eso la ficha pide esas tres claves. El Cuadro~\ref{tab:indice-casos} sirve
como índice de todos los casos especificados.

Si mantiene los casos en una herramienta externa, conserve al menos los mismos campos e indique
su ubicación en el Cuadro~\ref{tab:configuracion}.
\end{instruccion}

\begin{observacion}[Anexo sin formato]
\obsproblema{El anexo equivalente de la plantilla anterior decía únicamente «incluir los casos
de prueba aplicados», sin campos ni estructura, y el apartado que lo invocaba mezclaba el diseño
del caso con el resultado de su ejecución.}
\obscambio{Ficha completa con base, riesgo, condición, técnica, elemento de cobertura,
precondiciones, datos, pasos, resultado esperado, postcondiciones y limpieza, más un índice de
casos. Los campos de resultado se trasladan al \ref{anx:ejecuciones}.}
\obsfuente{OBS §2 y §6}
\end{observacion}

\subsection*{Ficha de caso de prueba}

\begin{xltabular}{\textwidth}{|E{4.4cm}|Y|}
\caption{Formato de especificación de un caso de prueba.}\label{tab:ficha-caso}\\ \hline
\endfirsthead
\multicolumn{2}{@{}l@{}}{\footnotesize\itshape\tablename~\thetable{} (continuación)}\\ \hline
\endhead
\multicolumn{2}{@{}r@{}}{\footnotesize\itshape continúa en la página siguiente}\\
\endfoot
\endlastfoot
  Identificador del caso        & \campo{CP-001} \\ \hline
  Versión del caso              & \campo{v1} \\ \hline
  Objeto de prueba              & \campo{OP-01 (Cuadro~\ref{tab:objetos})} \\ \hline
  Base de prueba                & \campo{BP-01 (Cuadro~\ref{tab:base-prueba})} \\ \hline
  Riesgo asociado               & \campo{RP-01 (Cuadro~\ref{tab:riesgos-producto})} \\ \hline
  Condición de prueba           & \campo{COND-01 (Cuadro~\ref{tab:condiciones})} \\ \hline
  Objetivo del caso             & \campo{Qué pretende evidenciar este caso en particular} \\ \hline
  Nivel de prueba               & \campo{Componente / integración / sistema / aceptación} \\ \hline
  Técnica aplicada              & \campo{Cuadro~\ref{tab:tecnicas}} \\ \hline
  Elemento(s) de cobertura      & \campo{Qué partición, valor límite, decisión o transición ejercita} \\ \hline
  Prioridad                     & \campo{Alta / media / baja, según el riesgo} \\ \hline
  Precondiciones                & \campo{Estado que debe tener el sistema antes de ejecutar} \\ \hline
  Datos de entrada              & \campo{Valores concretos y conjunto de datos (DAT-xx)} \\ \hline
  Acciones / pasos              & \campo{1. …\quad 2. …\quad 3. …} \\ \hline
  Resultado esperado            & \campo{Qué debe ocurrir; indique el oráculo del que se deriva} \\ \hline
  Postcondiciones               & \campo{Estado en que queda el sistema tras la ejecución} \\ \hline
  Limpieza requerida            & \campo{Qué debe deshacerse o restaurarse después} \\ \hline
  Dependencias                  & \campo{Otros casos que deban ejecutarse antes} \\ \hline
  Autor y fecha                 & \campo{Nombre — dd/mm/aaaa} \\ \hline
\end{xltabular}

El Cuadro~\ref{tab:ficha-caso} es el formato que debe replicarse por cada caso. Sus campos de
base, riesgo y condición son los que alimentan la matriz de trazabilidad del
\ref{anx:trazabilidad}; el campo «resultado esperado» es el que se referencia —no se copia—
desde cada ejecución.

\subsection*{Índice de casos de prueba especificados}

\begin{longtable}{|C{1.6cm}|L{4.0cm}|C{1.6cm}|C{1.6cm}|C{1.9cm}|C{1.6cm}|}
\caption{Índice de casos de prueba especificados.}
\label{tab:indice-casos} \\
\hline
\filaenc \textbf{ID} & \textbf{Título del caso} & \textbf{Objeto} & \textbf{Condición} & \textbf{Técnica} & \textbf{Prioridad} \\ \hline
\endfirsthead
\hline
\filaenc \textbf{ID} & \textbf{Título del caso} & \textbf{Objeto} & \textbf{Condición} & \textbf{Técnica} & \textbf{Prioridad} \\ \hline
\endhead
\hline \multicolumn{6}{|r|}{\footnotesize\textit{continúa en la página siguiente}} \\ \hline
\endfoot
\hline
\endlastfoot
\campo{CP-001} & \campo{Título} & \campo{OP-01} & \campo{COND-01} & \campo{Valores límite} & \campo{Alta} \\ \hline
 & & & & & \\ \hline
 & & & & & \\ \hline
 & & & & & \\ \hline
\end{longtable}

El Cuadro~\ref{tab:indice-casos} permite ver de un vistazo el conjunto diseñado y detectar
condiciones sin casos o técnicas declaradas que no se usaron.

% =====================================================================
%  ANEXO B. Especificación de procedimientos de prueba
% =====================================================================
\section{Especificación de procedimientos de prueba}
\label{anx:procedimientos}

\begin{notanormativa}[Diferencia con el caso de prueba]
Un caso de prueba especifica una verificación concreta. Un \textbf{procedimiento} agrupa y
\emph{ordena} varios casos en una secuencia ejecutable, añadiendo la preparación previa y la
limpieza posterior. Corresponde a la actividad \textbf{TD6} de \normap{2}: es el artefacto que
realmente se ejecuta.
\end{notanormativa}

\begin{observacion}[Anexo sin formato]
\obsproblema{El anexo equivalente de la plantilla anterior pedía «incluir los procedimientos de
prueba aplicados» sin proporcionar campos ni estructura.}
\obscambio{Ficha con identificador, propósito, casos incluidos y justificación de su orden,
entorno y datos requeridos, preparación, pasos, guion asociado y limpieza, más un índice.}
\obsfuente{OBS §2 y §6 · PDF §1.4 (TD6)}
\end{observacion}

\begin{instruccion}
Reproduzca la ficha siguiente por cada procedimiento. El orden de los casos importa: indique si
existen dependencias que obliguen a una secuencia concreta. La preparación y la limpieza son las
que hacen repetible el procedimiento; si se omiten, la segunda ejecución parte de un estado
distinto de la primera y los veredictos no son comparables.
\end{instruccion}

\subsection*{Ficha de procedimiento de prueba}

\begin{xltabular}{\textwidth}{|E{4.4cm}|Y|}
\caption{Formato de especificación de un procedimiento de prueba.}\label{tab:ficha-procedimiento}\\ \hline
\endfirsthead
\multicolumn{2}{@{}l@{}}{\footnotesize\itshape\tablename~\thetable{} (continuación)}\\ \hline
\endhead
\multicolumn{2}{@{}r@{}}{\footnotesize\itshape continúa en la página siguiente}\\
\endfoot
\endlastfoot
  Identificador del procedimiento & \campo{PP-001} \\ \hline
  Versión                         & \campo{v1} \\ \hline
  Propósito                       & \campo{Qué conjunto de verificaciones agrupa y con qué fin} \\ \hline
  Nivel de prueba                 & \campo{Componente / integración / sistema / aceptación} \\ \hline
  Casos incluidos y orden         & \campo{1.~CP-001 \quad 2.~CP-004 \quad 3.~CP-002} \\ \hline
  Justificación del orden         & \campo{Dependencias entre casos o restricciones del entorno} \\ \hline
  Entorno requerido               & \campo{ENT-xx (Cuadro~\ref{tab:req-entorno})} \\ \hline
  Datos requeridos                & \campo{DAT-xx (Cuadro~\ref{tab:datos})} \\ \hline
  Preparación previa              & \campo{Pasos para dejar el sistema en el estado inicial} \\ \hline
  Pasos de ejecución              & \campo{Secuencia detallada, manual o automatizada} \\ \hline
  Guion asociado                  & \campo{Ruta e identificación de versión del script, si aplica} \\ \hline
  Limpieza y restauración         & \campo{Qué se deshace al terminar para poder repetir} \\ \hline
  Duración estimada               & \campo{minutos/horas} \\ \hline
  Autor y fecha                   & \campo{Nombre — dd/mm/aaaa} \\ \hline
\end{xltabular}

El Cuadro~\ref{tab:ficha-procedimiento} define la unidad que se planifica y se ejecuta; su
identificador es el que aparece en el registro de ejecución del \ref{anx:ejecuciones} y en el
cronograma del apartado~\ref{sec:calendario}.

\subsection*{Índice de procedimientos}

\begin{longtable}{|C{1.8cm}|L{4.6cm}|L{4.0cm}|C{2.0cm}|}
\caption{Índice de procedimientos de prueba especificados.}
\label{tab:indice-procedimientos} \\
\hline
\filaenc \textbf{ID} & \textbf{Propósito} & \textbf{Casos incluidos} & \textbf{Nivel} \\ \hline
\endfirsthead
\hline
\filaenc \textbf{ID} & \textbf{Propósito} & \textbf{Casos incluidos} & \textbf{Nivel} \\ \hline
\endhead
\hline \multicolumn{4}{|r|}{\footnotesize\textit{continúa en la página siguiente}} \\ \hline
\endfoot
\hline
\endlastfoot
\campo{PP-001} & \campo{Propósito} & \campo{CP-001, CP-004} & \campo{Sistema} \\ \hline
 & & & \\ \hline
 & & & \\ \hline
\end{longtable}

El Cuadro~\ref{tab:indice-procedimientos} permite comprobar que todos los casos especificados
en el Cuadro~\ref{tab:indice-casos} están asignados a algún procedimiento; un caso huérfano
nunca llegará a ejecutarse.

% =====================================================================
%  ANEXO C. Registro de ejecución de pruebas
% =====================================================================
\section{Registro de ejecución de pruebas}
\label{anx:ejecuciones}

\begin{notanormativa}[Registro nuevo respecto de la plantilla original]
Este registro corresponde a la actividad \textbf{TE3 «Registrar la ejecución de la prueba»} de
\normap{2}, y es independiente de la especificación del caso:

\begin{itemize}
  \item El \textbf{caso de prueba} (\ref{anx:casos}) describe \emph{qué debería ejecutarse}.
  \item La \textbf{ejecución} registra \emph{qué ocurrió una vez determinada}: en qué versión,
        en qué entorno, con qué datos, quién la realizó y con qué veredicto.
\end{itemize}

Un mismo caso puede ejecutarse muchas veces —en distintas versiones, tras una corrección, en
una campaña de regresión—. Cada una es una \textbf{ejecución distinta} con su propio
identificador. \textbf{Nunca sobrescriba el diseño del caso con el resultado de una ejecución}:
al hacerlo se pierde el historial y deja de poder demostrarse cuándo se corrigió un defecto.
\end{notanormativa}

\begin{instruccion}
Use la ficha detallada para las ejecuciones que requieran documentación amplia (típicamente las
no aprobadas) y la tabla resumen para el conjunto. Los campos de versión, configuración y
entorno son los que hacen reproducible el resultado: tómelos del Cuadro~\ref{tab:configuracion}.
\end{instruccion}

\subsection*{Ficha de ejecución}

\begin{xltabular}{\textwidth}{|E{4.4cm}|Y|}
\caption{Formato de registro de una ejecución de prueba.}\label{tab:ficha-ejecucion}\\ \hline
\endfirsthead
\multicolumn{2}{@{}l@{}}{\footnotesize\itshape\tablename~\thetable{} (continuación)}\\ \hline
\endhead
\multicolumn{2}{@{}r@{}}{\footnotesize\itshape continúa en la página siguiente}\\
\endfoot
\endlastfoot
  Identificador de ejecución      & \campo{EJ-014} \\ \hline
  Caso o procedimiento ejecutado  & \campo{CP-001 (v1) / PP-001 (v1)} \\ \hline
  Ciclo o campaña                 & \campo{Ciclo 2 — regresión} \\ \hline
  Versión del software probado    & \campo{v1.2.1 — identificación exacta de la compilación} \\ \hline
  Configuración utilizada         & \campo{Parámetros, banderas, navegador, sistema operativo} \\ \hline
  Entorno                         & \campo{ENT-xx y estado verificado (Anexo~F)} \\ \hline
  Datos utilizados                & \campo{DAT-xx, versión del conjunto} \\ \hline
  Fecha y hora                    & \campo{dd/mm/aaaa hh:mm} \\ \hline
  Ejecutor                        & \campo{Nombre} \\ \hline
  Entradas efectivamente usadas   & \campo{Valores concretos, si difieren de los especificados} \\ \hline
  Resultado esperado (referencia) & \campo{Referencia al campo del caso; no se transcribe} \\ \hline
  Resultado real observado        & \campo{Qué ocurrió realmente} \\ \hline
  Veredicto                       & \campo{Aprobado / no aprobado / no concluyente / bloqueado} \\ \hline
  Evidencia                       & \campo{EVID-014 — ruta, captura, registro o archivo} \\ \hline
  Incidente asociado              & \campo{INC-05, si el veredicto no fue aprobado} \\ \hline
  Observaciones                   & \campo{Anomalías del entorno, desviaciones del procedimiento} \\ \hline
\end{xltabular}
\notatabla{Veredicto \textit{no concluyente}: no pudo determinarse si el comportamiento fue correcto. \textit{Bloqueado}: la ejecución no pudo completarse por una causa externa al objeto de prueba.}

El Cuadro~\ref{tab:ficha-ejecucion} conserva, para cada intento, las condiciones exactas en que
se obtuvo el veredicto. Es el eslabón que une el caso diseñado con la evidencia y el incidente.

\subsection*{Registro consolidado de ejecuciones}

\begin{longtable}{|C{1.5cm}|C{1.6cm}|C{1.5cm}|C{1.7cm}|C{1.5cm}|C{1.9cm}|C{1.5cm}|C{1.5cm}|}
\caption{Registro consolidado de ejecuciones de prueba.}
\label{tab:registro-ejecuciones} \\
\hline
\filaenc \textbf{ID ejec.} & \textbf{Caso/proc.} & \textbf{Versión SW} & \textbf{Entorno} & \textbf{Fecha} & \textbf{Ejecutor} & \textbf{Veredicto} & \textbf{Incidente} \\ \hline
\endfirsthead
\hline
\filaenc \textbf{ID ejec.} & \textbf{Caso/proc.} & \textbf{Versión SW} & \textbf{Entorno} & \textbf{Fecha} & \textbf{Ejecutor} & \textbf{Veredicto} & \textbf{Incidente} \\ \hline
\endhead
\hline \multicolumn{8}{|r|}{\footnotesize\textit{continúa en la página siguiente}} \\ \hline
\endfoot
\hline
\endlastfoot
\campo{EJ-001} & \campo{CP-001} & \campo{v1.2.0} & \campo{ENT-01} & \campo{dd/mm} & \campo{Nombre} & \campo{No aprob.} & \campo{INC-05} \\ \hline
\campo{EJ-014} & \campo{CP-001} & \campo{v1.2.1} & \campo{ENT-01} & \campo{dd/mm} & \campo{Nombre} & \campo{Aprobado} & \campo{INC-05 (conf.)} \\ \hline
 & & & & & & & \\ \hline
 & & & & & & & \\ \hline
\end{longtable}

El Cuadro~\ref{tab:registro-ejecuciones} muestra el patrón esperado: el mismo caso CP-001
aparece dos veces, con identificadores de ejecución distintos y versiones distintas del
software. La segunda es la \emph{prueba de confirmación} del incidente INC-05
(apartado~\ref{sec:confirmacion-regresion}); la primera no se borra ni se modifica.

% =====================================================================
%  ANEXO D. Registro de incidentes
% =====================================================================
\section{Registro de incidentes}
\label{anx:incidentes}

\begin{notanormativa}[Registro nuevo respecto de la plantilla original]
Corresponde a la actividad \textbf{IR2 «Crear o actualizar el reporte de incidente»} de
\normap{2}. Las reglas de análisis previo, las escalas de severidad y prioridad y el flujo de
estados se definen en el apartado~\ref{sec:gestion-incidentes}; aquí se registran los
incidentes concretos.
\end{notanormativa}

\begin{instruccion}
Reproduzca la ficha por cada incidente reportado. Los pasos de reproducción son el campo que
determina si el incidente podrá ser atendido: escríbalos de forma que una persona ajena al
equipo pueda reproducir el fallo sin preguntar nada.

Recuerde que \textbf{severidad y prioridad son campos distintos} y que un incidente puede tener
como origen el entorno, los datos o el propio caso de prueba, y no el objeto probado.
\end{instruccion}

\subsection*{Ficha de incidente}

\begin{xltabular}{\textwidth}{|E{4.4cm}|Y|}
\caption{Formato de registro de un incidente de prueba.}\label{tab:ficha-incidente}\\ \hline
\endfirsthead
\multicolumn{2}{@{}l@{}}{\footnotesize\itshape\tablename~\thetable{} (continuación)}\\ \hline
\endhead
\multicolumn{2}{@{}r@{}}{\footnotesize\itshape continúa en la página siguiente}\\
\endfoot
\endlastfoot
  Identificador                & \campo{INC-05} \\ \hline
  Resumen                      & \campo{Una línea que identifique el problema} \\ \hline
  Descripción                  & \campo{Explicación detallada de lo observado} \\ \hline
  Origen determinado (IR1)     & \campo{Objeto de prueba / entorno / datos / caso de prueba / duplicado} \\ \hline
  Ejecución que lo reveló      & \campo{EJ-001 (\ref{anx:ejecuciones})} \\ \hline
  Objeto de prueba afectado    & \campo{OP-01} \\ \hline
  Resultado esperado           & \campo{Qué debía ocurrir} \\ \hline
  Resultado obtenido           & \campo{Qué ocurrió} \\ \hline
  Pasos de reproducción        & \campo{1. …\quad 2. …\quad 3. …} \\ \hline
  ¿Reproducible?               & \campo{Siempre / intermitente / no reproducido} \\ \hline
  Versión del software         & \campo{v1.2.0} \\ \hline
  Entorno y configuración      & \campo{ENT-xx y parámetros} \\ \hline
  Datos utilizados             & \campo{DAT-xx} \\ \hline
  Evidencia                    & \campo{Capturas, registros, archivos} \\ \hline
  Severidad                    & \campo{Cuadro~\ref{tab:severidad}} \\ \hline
  Prioridad                    & \campo{Cuadro~\ref{tab:prioridad}} \\ \hline
  Riesgo de producto asociado  & \campo{RP-01} \\ \hline
  Reportado por / fecha        & \campo{Nombre — dd/mm/aaaa} \\ \hline
  Responsable asignado         & \campo{Nombre o equipo} \\ \hline
  Estado                       & \campo{Cuadro~\ref{tab:estados-incidente}} \\ \hline
  Fecha de última actualización& \campo{dd/mm/aaaa} \\ \hline
  Resolución o disposición     & \campo{Qué se hizo, o por qué se difiere o rechaza} \\ \hline
  Ejecución de confirmación    & \campo{EJ-014, si ya se confirmó la corrección} \\ \hline
\end{xltabular}

El Cuadro~\ref{tab:ficha-incidente} contiene los campos mínimos para que un incidente pueda
comunicarse, priorizarse y cerrarse de forma verificable. El campo de origen evita que un
problema del entorno se contabilice como defecto del producto.

\subsection*{Registro consolidado de incidentes}

\begin{longtable}{|C{1.4cm}|L{3.6cm}|C{1.5cm}|C{1.5cm}|C{1.6cm}|C{1.8cm}|C{1.6cm}|}
\caption{Registro consolidado de incidentes.}
\label{tab:registro-incidentes} \\
\hline
\filaenc \textbf{ID} & \textbf{Resumen} & \textbf{Sever.} & \textbf{Prior.} & \textbf{Estado} & \textbf{Responsable} & \textbf{Ejecución} \\ \hline
\endfirsthead
\hline
\filaenc \textbf{ID} & \textbf{Resumen} & \textbf{Sever.} & \textbf{Prior.} & \textbf{Estado} & \textbf{Responsable} & \textbf{Ejecución} \\ \hline
\endhead
\hline \multicolumn{7}{|r|}{\footnotesize\textit{continúa en la página siguiente}} \\ \hline
\endfoot
\hline
\endlastfoot
\campo{INC-05} & \campo{Resumen} & \campo{Alta} & \campo{Normal} & \campo{Cerrado} & \campo{Nombre} & \campo{EJ-001} \\ \hline
 & & & & & & \\ \hline
 & & & & & & \\ \hline
\end{longtable}

El Cuadro~\ref{tab:registro-incidentes} es la vista de conjunto que alimenta el seguimiento
(apartado~\ref{sec:seguimiento}) y, al cierre, el
Cuadro~\ref{tab:incidentes-abiertos}.

% =====================================================================
%  ANEXO E. Informe de estado de la prueba
% =====================================================================
\section{Informe de estado de la prueba}
\label{anx:informes}

\begin{notanormativa}[Correspondencia con el modelo de procesos]
Corresponde a la actividad \textbf{TMC4 «Reportar»} de \normap{2}. Es el producto periódico del
seguimiento definido en el apartado~\ref{sec:seguimiento} y se dirige a los destinatarios del
Cuadro~\ref{tab:comunicacion}.
\end{notanormativa}

\begin{observacion}[No existía informe de estado]
\obsproblema{La plantilla anterior no preveía ningún reporte periódico: entre una entrega del
curso y la siguiente no quedaba constancia del avance, la cobertura, los incidentes ni el
esfuerzo consumido.}
\obscambio{Formato reutilizable con cinco bloques —identificación, avance, resultados de
ejecución, incidentes y riesgos, recursos, decisiones y acciones—, una copia por periodo.}
\obsfuente{OBS §3 y §6 · PDF §1.3 (TMC4)}
\end{observacion}

\begin{instruccion}
Reproduzca este formato una vez por periodo. Un informe de estado útil no se limita a decir
cuánto se avanzó: contrasta lo previsto con lo real, señala lo que está bloqueado y termina con
decisiones y acciones que tienen responsable y fecha. Si un periodo no requiere acciones,
escríbalo explícitamente en lugar de dejar la sección vacía.
\end{instruccion}

\begin{xltabular}{\textwidth}{|E{4.6cm}|Y|}
\caption{Formato de informe de estado de la prueba.}\label{tab:ficha-informe}\\ \hline
\endfirsthead
\multicolumn{2}{@{}l@{}}{\footnotesize\itshape\tablename~\thetable{} (continuación)}\\ \hline
\endhead
\multicolumn{2}{@{}r@{}}{\footnotesize\itshape continúa en la página siguiente}\\
\endfoot
\endlastfoot
  \multicolumn{2}{|l|}{\cellcolor{uvgris}\textbf{Identificación}} \\ \hline
  Informe número / periodo      & \campo{IE-03 — dd/mm/aaaa a dd/mm/aaaa} \\ \hline
  Elaborado por / fecha         & \campo{Nombre — dd/mm/aaaa} \\ \hline
  Destinatarios                 & \campo{Cuadro~\ref{tab:comunicacion}} \\ \hline
  \multicolumn{2}{|l|}{\cellcolor{uvgris}\textbf{Avance}} \\ \hline
  Avance previsto en el periodo & \campo{\% o hitos comprometidos} \\ \hline
  Avance real alcanzado         & \campo{\% o hitos logrados} \\ \hline
  Desviación y su causa         & \campo{Diferencia y explicación} \\ \hline
  \multicolumn{2}{|l|}{\cellcolor{uvgris}\textbf{Resultados de ejecución}} \\ \hline
  Pruebas ejecutadas            & \campo{n} \\ \hline
  Aprobadas                     & \campo{n} \\ \hline
  Fallidas                      & \campo{n} \\ \hline
  Bloqueadas                    & \campo{n y causa del bloqueo} \\ \hline
  Cobertura alcanzada           & \campo{COB-xx = \% (indique la medida, no sólo el número). Mide lo ejercitado; las pruebas fallidas cuentan igualmente} \\ \hline
  \multicolumn{2}{|l|}{\cellcolor{uvgris}\textbf{Incidentes y riesgos}} \\ \hline
  Incidentes nuevos             & \campo{n; claves de los relevantes} \\ \hline
  Incidentes cerrados           & \campo{n} \\ \hline
  Incidentes abiertos críticos  & \campo{Claves y estado} \\ \hline
  Riesgos materializados        & \campo{Claves} \\ \hline
  Riesgos nuevos identificados  & \campo{Descripción y exposición} \\ \hline
  \multicolumn{2}{|l|}{\cellcolor{uvgris}\textbf{Recursos}} \\ \hline
  Esfuerzo consumido / estimado & \campo{h reales de h estimadas} \\ \hline
  Bloqueos y dependencias       & \campo{Qué impide avanzar y de quién depende} \\ \hline
  \multicolumn{2}{|l|}{\cellcolor{uvgris}\textbf{Decisiones y acciones}} \\ \hline
  Decisiones tomadas            & \campo{Qué se decidió y quién} \\ \hline
  Acciones acordadas            & \campo{Acción — responsable — fecha objetivo} \\ \hline
  ¿Requiere replanificación?    & \campo{Sí/No; si sí, según qué regla del Cuadro~\ref{tab:control}} \\ \hline
\end{xltabular}

El Cuadro~\ref{tab:ficha-informe} agrupa la información en cinco bloques —identificación,
avance, resultados, incidentes y riesgos, recursos, decisiones— de modo que el lector pueda
localizar rápidamente lo que le concierne. Cada informe emitido debería quedar archivado según
el Cuadro~\ref{tab:configuracion}.

% =====================================================================
%  ANEXO F. Registro de preparación de entorno y datos
% =====================================================================
\section{Registro de preparación de entorno y datos}
\label{anx:entorno}

\begin{notanormativa}[Registro nuevo respecto de la plantilla original]
Corresponde a las actividades \textbf{ES1 «Establecer el entorno de prueba»} y \textbf{ES2
«Mantener el entorno de prueba»} de \normap{2}, e implementa la verificación de los criterios de
entrada CE-04 y CE-05 del Cuadro~\ref{tab:criterios-entrada}. La plantilla original describía
los requisitos del entorno, pero no dejaba constancia de que se hubieran cumplido antes de
ejecutar.
\end{notanormativa}

\begin{instruccion}
Complete un registro por cada ciclo de prueba, o cada vez que el entorno cambie. Los requisitos
que aquí se verifican se definieron en los apartados~\ref{sec:entorno} y~\ref{sec:datos}, y su
cumplimiento habilita el inicio del ciclo según el apartado~\ref{sec:criterios-entrada}. Es el documento
que permite responder, meses después, a la pregunta «¿en qué condiciones exactas se obtuvo este
resultado?». Sin él, el registro de ejecución del \ref{anx:ejecuciones} queda incompleto.
\end{instruccion}

\subsection*{Verificación de disponibilidad}

\begin{table}[H]
\centering
\caption{Registro de preparación y verificación del entorno y los datos.}
\label{tab:registro-entorno}
\begin{tabularx}{\textwidth}{|C{1.5cm}|L{2.9cm}|C{1.8cm}|Y|C{1.7cm}|C{1.8cm}|}
  \hline
  \filaenc \textbf{Clave} & \textbf{Componente o conjunto} & \textbf{Versión instalada} & \textbf{Configuración aplicada} & \textbf{Estado} & \textbf{Verificado por / fecha} \\ \hline
  \campo{ENT-01} & \campo{Servidor de aplicación} & \campo{v0.0} & \campo{Parámetros} & \campo{Listo} & \campo{Nombre — dd/mm} \\ \hline
  \campo{ENT-02} & \campo{Base de datos} & \campo{v0.0} & \campo{Parámetros} & \campo{Listo} & \campo{Nombre — dd/mm} \\ \hline
  \campo{DAT-01} & \campo{Catálogo de clientes} & \campo{v1} & \campo{Cargado y anonimizado} & \campo{Listo} & \campo{Nombre — dd/mm} \\ \hline
  & & & & & \\ \hline
\end{tabularx}
\notatabla{Estado: \textit{Listo} · \textit{Listo con restricciones} (indíquelas) · \textit{No disponible}. Un ciclo que inicia con algún componente «no disponible» debe registrarse como desviación.}
\end{table}

El Cuadro~\ref{tab:registro-entorno} deja constancia firmada de que el entorno y los datos
estaban en el estado requerido antes de ejecutar, y de quién lo comprobó.

\subsection*{Cambios y restauración durante el ciclo}

\begin{table}[H]
\centering
\caption{Cambios aplicados al entorno o a los datos durante el ciclo.}
\label{tab:cambios-entorno}
\begin{tabularx}{\textwidth}{|C{1.9cm}|C{1.8cm}|Y|C{2.4cm}|C{2.2cm}|}
  \hline
  \filaenc \textbf{Fecha} & \textbf{Componente} & \textbf{Cambio aplicado y motivo} & \textbf{Ejecuciones afectadas} & \textbf{Responsable} \\ \hline
  \campo{dd/mm} & \campo{ENT-01} & \campo{Qué cambió y por qué} & \campo{EJ-010 en adelante} & \campo{Nombre} \\ \hline
  & & & & \\ \hline
\end{tabularx}
\end{table}

\begin{table}[H]
\centering
\caption{Restauración del entorno a un estado conocido.}
\label{tab:restauracion}
\begin{tabularx}{\textwidth}{|C{1.9cm}|Y|C{2.6cm}|C{2.2cm}|}
  \hline
  \filaenc \textbf{Fecha} & \textbf{Acción de restauración realizada} & \textbf{Estado resultante} & \textbf{Responsable} \\ \hline
  \campo{dd/mm} & \campo{Qué se restauró y cómo} & \campo{Estado base} & \campo{Nombre} \\ \hline
  & & & \\ \hline
\end{tabularx}
\end{table}

Los Cuadros~\ref{tab:cambios-entorno} y~\ref{tab:restauracion} documentan la actividad ES2. La
columna «ejecuciones afectadas» es la que permite detectar cuándo dos resultados no son
comparables porque el entorno cambió entre ellos.

% =====================================================================
%  ANEXO G. Matriz de trazabilidad
% =====================================================================
\section{Matriz de trazabilidad}
\label{anx:trazabilidad}

\begin{notanormativa}[Para qué sirve la trazabilidad]
La cadena debe poder recorrerse en ambos sentidos: desde un requisito hasta el resultado que lo
evidencia, y desde un incidente hasta el requisito y el riesgo que justificaron probarlo.
\end{notanormativa}

\begin{observacion}[Trazabilidad sin cadena definida]
\obsproblema{La plantilla anterior pedía «elaborar una tabla de trazabilidad» para los elementos
de prueba, sin indicar qué debía encadenar ni con qué identificadores.}
\obscambio{Matriz de nueve columnas —base, riesgo, condición, elemento de cobertura, caso,
procedimiento, ejecución, evidencia e incidente—, con la Figura~\ref{fig:trazabilidad} que
señala dónde se registra cada eslabón y un resumen de cobertura de la base de prueba.}
\obsfuente{OBS §6}
\end{observacion}

% =====================================================================
%  diagrams/cadena-trazabilidad.tex
%  Figura: cadena de trazabilidad que la plantilla permite representar,
%          con el apartado o anexo donde se registra cada eslabón.
% =====================================================================
\begin{figure}[htbp]
\centering
\begin{tikzpicture}[
  font=\scriptsize,
  esl/.style ={draw=uvazul, thick, fill=uvazul!8, rounded corners=2pt,
               text width=2.45cm, align=center, minimum height=0.95cm, inner sep=3pt},
  ori/.style ={font=\tiny\itshape, text=gray!65, align=center, text width=2.6cm},
  fl/.style  ={-{Latex[length=4pt]}, draw=uvazul!85, thick}
]

\node[esl] (r1) at (-5.6,0)  {\textbf{Base de prueba}\\ BP-xx};
\node[esl] (r2) at (-2.8,0)  {\textbf{Riesgo}\\ RP-xx};
\node[esl] (r3) at (0,0)     {\textbf{Condición de prueba}\\ COND-xx};
\node[esl] (r4) at (2.8,0)   {\textbf{Elemento de cobertura}};
\node[esl] (r5) at (5.6,0)   {\textbf{Caso de prueba}\\ CP-xxx};

% La segunda fila se dispone de derecha a izquierda («en boustrophedon»)
% para que el enlace entre CP-xxx y PP-xxx sea un salto vertical corto y
% no una línea que atraviese todo el diagrama.
\node[esl] (r6) at (5.6,-2.6)  {\textbf{Procedimiento}\\ PP-xxx};
\node[esl] (r7) at (2.8,-2.6)  {\textbf{Ejecución}\\ EJ-xxx};
\node[esl] (r8) at (0,-2.6)    {\textbf{Evidencia}\\ EVID-xxx};
\node[esl] (r9) at (-2.8,-2.6) {\textbf{Incidente}\\ INC-xx};

\foreach \a/\b in {r1/r2, r2/r3, r3/r4, r4/r5, r6/r7, r7/r8, r8/r9} { \draw[fl] (\a) -- (\b); }
\draw[fl] (r5.south) -- (r6.north);

\node[ori] at (-5.6,-1.0) {Cuadro~\ref{tab:base-prueba}};
\node[ori] at (-2.8,-1.0) {Cuadro~\ref{tab:riesgos-producto}};
\node[ori] at (0,-1.0)    {Cuadro~\ref{tab:condiciones}};
\node[ori] at (2.8,-1.0)  {Cuadro~\ref{tab:condiciones}};
\node[ori] at (5.6,-1.0)  {Anexo A};

\node[ori] at (5.6,-3.6)  {Anexo B};
\node[ori] at (2.8,-3.6)  {Anexo C};
\node[ori] at (0,-3.6)    {Anexo C};
\node[ori] at (-2.8,-3.6) {Anexo D};

\node[align=left, font=\scriptsize, text width=15.2cm] at (0,-4.8)
  {\textbf{Cómo leer la figura.} Cada eslabón se registra en el cuadro o anexo indicado bajo él.
   La cadena se recorre en ambos sentidos: hacia adelante, para comprobar que un requisito quedó
   verificado; hacia atrás, para saber qué requisito y qué riesgo justificaron una prueba que
   reveló un incidente. No todos los proyectos necesitan todos los eslabones; la plantilla
   permite representarlos, no obliga a usarlos.};

\end{tikzpicture}
\caption[Cadena de trazabilidad representable en la plantilla]%
{Cadena de trazabilidad desde la base de prueba hasta el incidente, con el cuadro o anexo en
que se registra cada eslabón.}
\label{fig:trazabilidad}
\end{figure}


La Figura~\ref{fig:trazabilidad} muestra la cadena que esta plantilla permite representar y el
apartado o anexo donde se registra cada eslabón. No todos los proyectos necesitan todos los
eslabones: un proyecto sin análisis de riesgo formal, o sin elementos de cobertura explícitos,
puede dejar esas columnas vacías. Lo que la plantilla garantiza es que sean \emph{expresables}.

\begin{instruccion}
Complete una fila por cada relación que quiera poder demostrar. Los identificadores provienen de
los cuadros ya definidos: base de prueba (Cuadro~\ref{tab:base-prueba}), riesgos
(Cuadro~\ref{tab:riesgos-producto}), condiciones (Cuadro~\ref{tab:condiciones}), casos
(Cuadro~\ref{tab:indice-casos}), procedimientos (Cuadro~\ref{tab:indice-procedimientos}),
ejecuciones (Cuadro~\ref{tab:registro-ejecuciones}) e incidentes
(Cuadro~\ref{tab:registro-incidentes}).

Use la matriz también como herramienta de control: una fila de la base de prueba sin caso
asociado indica una brecha de cobertura; un caso sin base ni riesgo indica esfuerzo que nadie
pidió.
\end{instruccion}

\begin{footnotesize}
\begin{longtable}{|C{1.35cm}|C{1.35cm}|C{1.45cm}|C{1.8cm}|C{1.3cm}|C{1.3cm}|C{1.3cm}|C{1.45cm}|C{1.3cm}|}
\caption{Matriz de trazabilidad de la base de prueba al incidente.}
\label{tab:trazabilidad} \\
\hline
\filaenc \textbf{Base} & \textbf{Riesgo} & \textbf{Condición} & \textbf{Elemento de cobertura} & \textbf{Caso} & \textbf{Proc.} & \textbf{Ejec.} & \textbf{Evidencia} & \textbf{Incid.} \\ \hline
\endfirsthead
\hline
\filaenc \textbf{Base} & \textbf{Riesgo} & \textbf{Condición} & \textbf{Elemento de cobertura} & \textbf{Caso} & \textbf{Proc.} & \textbf{Ejec.} & \textbf{Evidencia} & \textbf{Incid.} \\ \hline
\endhead
\hline \multicolumn{9}{|r|}{\textit{continúa en la página siguiente}} \\ \hline
\endfoot
\hline
\endlastfoot
BP-01 & RP-01 & COND-01 & Partición P3 & CP-007 & PP-001 & EJ-014 & EVID-014 & INC-05 \\ \hline
\campo{} & \campo{} & \campo{} & \campo{} & \campo{} & \campo{} & \campo{} & \campo{} & \campo{} \\ \hline
 & & & & & & & & \\ \hline
 & & & & & & & & \\ \hline
 & & & & & & & & \\ \hline
 & & & & & & & & \\ \hline
\end{longtable}
\end{footnotesize}

La primera fila del Cuadro~\ref{tab:trazabilidad} está rellena a modo de ejemplo y reproduce la
cadena completa: el requisito BP-01, motivado por el riesgo RP-01, generó la
condición COND-01, cuya partición P3 se ejercita mediante el caso CP-007, incluido en el
procedimiento PP-001; su ejecución EJ-014 produjo la evidencia EVID-014 y dio origen al
incidente INC-05. Sustituya esa fila por los datos reales del proyecto.

\begin{instruccion}
Si el volumen de filas hace inmanejable esta matriz en el documento, manténgala en una hoja de
cálculo o herramienta y deje aquí la referencia a su ubicación y responsable, igual que con el
resto de los registros. Lo que no es admisible es no poder reconstruir la cadena.
\end{instruccion}

\subsection*{Cobertura de la base de prueba}

\begin{instruccion}
Este cuadro resume la matriz anterior desde el punto de vista de la base de prueba. Es la
fuente del dato de la medida COB-01 del Cuadro~\ref{tab:cobertura} y sirve para detectar
elementos de la base que quedaron sin ejercitar.

Las dos últimas columnas responden a preguntas distintas y \textbf{no deben fundirse}:
«¿Cubierto?» dice si el elemento llegó a ejercitarse; «Veredictos» dice cómo se comportó. Un
requisito ejercitado por casos que fallaron está \emph{cubierto} y además tiene un
\emph{resultado adverso}: cuenta en el numerador de COB-01 y, a la vez, en los incidentes del
\ref{anx:incidentes}.
\end{instruccion}

\begin{table}[H]
\centering
\caption{Resumen de cobertura de la base de prueba y de sus resultados.}
\label{tab:cobertura-base}
\begin{tabularx}{\textwidth}{|C{1.5cm}|Y|C{1.9cm}|C{1.9cm}|C{2.2cm}|C{2.3cm}|}
  \hline
  \filaenc \textbf{Base} & \textbf{Descripción} & \textbf{Casos asociados} & \textbf{Casos ejecutados} & \textbf{¿Cubierto? (COB-01)} & \textbf{Veredictos (apr. / no apr.)} \\ \hline
  \campo{BP-01} & \campo{Descripción} & \campo{3} & \campo{3} & \campo{Sí} & \campo{3 / 0} \\ \hline
  \campo{BP-02} & \campo{Descripción} & \campo{2} & \campo{2} & \campo{Sí} & \campo{0 / 2} \\ \hline
  \campo{BP-03} & \campo{Descripción} & \campo{0} & \campo{0} & \campo{No — brecha} & \campo{---} \\ \hline
  & & & & & \\ \hline
\end{tabularx}
\notatabla{«¿Cubierto?» es \textit{Sí} en cuanto el elemento tiene al menos un caso ejecutado, cualquiera que sea el veredicto. Sólo las columnas de casos ejecutados y de veredictos cambian al repetir una prueba; la de cobertura no retrocede.}
\end{table}

El Cuadro~\ref{tab:cobertura-base} hace visibles dos situaciones que conviene no confundir. La
fila BP-03 es una \textbf{brecha de cobertura}: un elemento de la base que quedó sin ejercitar
y que, si no se justifica como exclusión en el Cuadro~\ref{tab:exclusiones}, debe aparecer
entre los riesgos residuales del apartado~\ref{sec:residuales}. La fila BP-02, en cambio,
está plenamente cubierta y aun así arroja un \textbf{resultado adverso}: sus dos casos
fallaron. Fundir ambas columnas en una sola haría desaparecer una de las dos situaciones, y
contar BP-02 como «no cubierta» ocultaría que la prueba sí se realizó y encontró un defecto.


\end{document}
